\documentclass[11pt,oneside]{book}
\usepackage[a4paper,margin=27mm,headheight=15pt]{geometry}
\usepackage{fontspec}
\setmainfont{TeX Gyre Pagella}
\setsansfont{TeX Gyre Heros}
\setmonofont{Latin Modern Mono}[Scale=MatchLowercase]
\usepackage{amsmath,amssymb,amsthm,mathtools}
\usepackage{microtype}
\usepackage{booktabs,longtable,array,enumitem}
\usepackage{xcolor,xurl,hyperref,bookmark}
\usepackage{fancyhdr}
\definecolor{GuideBlue}{HTML}{214B64}
\hypersetup{colorlinks=true,linkcolor=GuideBlue,citecolor=GuideBlue,urlcolor=GuideBlue,
  pdfauthor={},pdfcreator={},pdfsubject={Mathematical documentation for ModularRep}}
\urlstyle{tt}
\newcommand{\lean}[1]{{\small\nolinkurl{#1}}}
\newcommand{\module}[1]{\lean{#1}}
\newenvironment{implementation}{\par\begingroup\small\raggedright
  \noindent\textsf{Implementation.}\par\nobreak}{\par\endgroup}
\newcommand{\source}[1]{\textup{#1}}
\newcommand{\GAP}{GAP}
\newcommand{\Irr}{\operatorname{Irr}}
\newcommand{\IBr}{\operatorname{IBr}}
\newcommand{\Alp}{\operatorname{Alp}}
\newcommand{\Aut}{\operatorname{Aut}}
\newcommand{\Out}{\operatorname{Out}}
\newcommand{\Irrdz}{\operatorname{dz}}
\newcommand{\bl}{\operatorname{bl}}
\newcommand{\Stab}{\operatorname{Stab}}
\newcommand{\GL}{\operatorname{GL}}
\newcommand{\SL}{\operatorname{SL}}
\newcommand{\SU}{\operatorname{SU}}
\newcommand{\Sp}{\operatorname{Sp}}
\newcommand{\PSp}{\operatorname{PSp}}
\newcommand{\Spin}{\operatorname{Spin}}
\newcommand{\CSp}{\operatorname{CSp}}
\newcommand{\PCSp}{\operatorname{PCSp}}
\newcommand{\Cl}{\operatorname{Cl}}
\newcommand{\tr}{\operatorname{tr}}
\newcommand{\rad}{\operatorname{rad}}
\newcommand{\im}{\operatorname{im}}
\newcommand{\id}{\operatorname{id}}
\newcommand{\rk}{\operatorname{rk}}
\newcommand{\defect}{\operatorname{def}}
\newcommand{\Z}{\mathbb Z}
\newcommand{\Q}{\mathbb Q}
\newcommand{\F}{\mathbb F}
\newcommand{\N}{\mathbb N}
\newcommand{\C}{\mathbb C}
\newcommand{\sourcegrade}[1]{\textsf{#1}}
\newtheorem{theorem}{Theorem}[chapter]
\newtheorem{proposition}[theorem]{Proposition}
\newtheorem{lemma}[theorem]{Lemma}
\newtheorem{corollary}[theorem]{Corollary}
\theoremstyle{definition}
\newtheorem{definition}[theorem]{Definition}
\newtheorem{example}[theorem]{Example}
\newtheorem{hypothesis}[theorem]{Hypothesis}
\theoremstyle{remark}
\newtheorem{remark}[theorem]{Remark}
\setlist{topsep=5pt,itemsep=3pt,parsep=0pt}
\setlength{\parskip}{4pt}
\setlength{\parindent}{0pt}
\setlength{\emergencystretch}{2em}
\pagestyle{fancy}
\fancyhf{}
\fancyhead[L]{\small\sffamily\nouppercase{\leftmark}}
\fancyfoot[C]{\small\thepage}
\renewcommand{\headrulewidth}{0.3pt}
\renewcommand{\chaptermark}[1]{\markboth{#1}{}}
\fancypagestyle{plain}{\fancyhf{}\fancyfoot[C]{\small\thepage}\renewcommand{\headrulewidth}{0pt}}
\setcounter{tocdepth}{1}
\setcounter{secnumdepth}{2}

% One maintained definition for each bibliographic reference.
\newcommand{\DeclareManualReference}[2]{%
  \expandafter\def\csname ManualReferenceBody@#1\endcsname{#2}}
\newcommand{\ManualReference}[1]{%
  \ifcsname ManualReferenceBody@#1\endcsname
    \bibitem{#1}\csname ManualReferenceBody@#1\endcsname
  \else
    \PackageError{ModularRep}{Unknown bibliography key #1}{Add its definition to bibliography-entries.tex.}%
  \fi}
% Canonical bibliographic definitions shared by the report and manual.
% Document-specific selection files preserve each document's reference order.

\DeclareManualReference{Alperin1987}{%
J.~L. Alperin,
{\it Weights for finite groups},
in {\it The Arcata Conference on Representations of Finite Groups},
Proc.\ Sympos.\ Pure Math., vol.~47, part~1, Amer.\ Math.\ Soc.,
Providence, RI, 1987, 369--379.
}

\DeclareManualReference{An1993}{%
J. An,
{\it $2$-weights for classical groups},
J.\ Reine Angew.\ Math.\ \textbf{439} (1993), 159--204.
}

\DeclareManualReference{An1994}{%
J. An,
{\it Weights for classical groups},
Trans.\ Amer.\ Math.\ Soc.\ \textbf{342} (1994), no.~1, 1--42.
}

\DeclareManualReference{AnCannonOBrienUngerFi24}{%
J. An, J.~J. Cannon, E.~A. O'Brien, and W.~R. Unger,
{\it The Alperin weight conjecture and Dade's conjecture for the simple
group $Fi'_{24}$},
LMS J.\ Comput.\ Math.\ \textbf{11} (2008), 100--145.
}

\DeclareManualReference{AnDietrich2012}{%
J. An and H. Dietrich,
{\it The AWC-goodness and essential rank of sporadic simple groups},
J.\ Algebra \textbf{356} (2012), 325--354.
}

\DeclareManualReference{AnDietrichErratum}{%
J. An and H. Dietrich,
{\it Errata for ``The AWC-goodness and essential rank of sporadic simple
groups''},
version of 1 January 2013, available at
\url{https://users.monash.edu/~heikod/hddats/errataAWCgood.pdf}.
}

\DeclareManualReference{AnOBrienWilsonJ4}{%
J. An, E.~A. O'Brien, and R.~A. Wilson,
{\it The Alperin weight conjecture and Dade's conjecture for the simple
group $J_4$},
LMS J.\ Comput.\ Math.\ \textbf{6} (2003), 119--140.
}

\DeclareManualReference{AnWilsonBaby}{%
J. An and R.~A. Wilson,
{\it The Alperin weight conjecture and Uno's conjecture for the Baby
Monster $B$, $p$ odd},
LMS J.\ Comput.\ Math.\ \textbf{7} (2004), 120--166.
}

\DeclareManualReference{AnWilsonMonster}{%
J. An and R.~A. Wilson,
{\it The Alperin weight conjecture and Uno's conjecture for the Monster
$M$, $p$ odd},
LMS J.\ Comput.\ Math.\ \textbf{13} (2010), 320--356.
}

\DeclareManualReference{Bonnafe2005}{%
C. Bonnaf\'e,
{\it Quasi-isolated elements in reductive groups},
Comm.\ Algebra \textbf{33} (2005), no.~7, 2315--2337.
}

\DeclareManualReference{BonnafeDatRouquier2017}{%
C. Bonnaf\'e, J.-F. Dat, and R. Rouquier,
{\it Derived categories and Deligne--Lusztig varieties II},
Ann.\ of Math.\ (2) \textbf{185} (2017), no.~2, 609--670.
}

\DeclareManualReference{Bouc2000}{%
S. Bouc,
{\it Burnside rings},
in {\it Handbook of Algebra}, vol.~2, Elsevier, 2000, 739--804.
Corrected author version, 16 June 2016: \url{https://www.lamfa.u-picardie.fr/bouc/burnamscorr.pdf}.
}

\DeclareManualReference{BreuerSporadicOverview}{%
T. Breuer,
{\it Computations for some simple groups},
documentation distributed with CTBlocks, version~0.9.5 (31 May 2022),
\href{https://www.math.rwth-aachen.de/~Thomas.Breuer/ctblocks/doc/overview.html}
{online overview} and
\href{https://www.math.rwth-aachen.de/~Thomas.Breuer/ctblocks/ctblocks-0.9.5.tar.gz}
{versioned source archive}, accessed 19 September 2026.
}

\DeclareManualReference{CTblLib2025}{%
T. Breuer,
{\it CTblLib, the \textup{\GAP} Character Table Library, Version 1.3.11},
2025, available at
\href{https://www.gap-system.org/Packages/ctbllib.html}
{www.gap-system.org/Packages/ctbllib.html}.
}

\DeclareManualReference{BreuerMagaardWilsonBaby2020}{%
T. Breuer, K. Magaard, and R.~A. Wilson,
{\it Verification of the ordinary character table of the Baby Monster},
J.\ Algebra \textbf{561} (2020), 111--130.
}

\DeclareManualReference{BreuerMagaardWilsonMonster2026}{%
T. Breuer, K. Magaard, and R.~A. Wilson,
{\it Verification of the conjugacy classes and ordinary character table of
the Monster},
J.\ Algebra \textbf{703} (2026), 452--463.
}

\DeclareManualReference{BroueMichel1989}{%
M. Brou\'e and J. Michel,
{\it Blocs et s\'eries de Lusztig dans un groupe r\'eductif fini},
J.\ Reine Angew.\ Math.\ \textbf{395} (1989), 56--67.
}

\DeclareManualReference{BroughSchaefferFry2020}{%
J. Brough and A.~A. Schaeffer Fry,
{\it Radical subgroups and the inductive blockwise Alperin weight
conditions for $\operatorname{PSp}_4(q)$},
Rocky Mountain J.\ Math.\ \textbf{50} (2020), 1181--1205.
}

\DeclareManualReference{BroughSpath2022}{%
J. Brough and B. Sp\"ath,
{\it A criterion for the inductive Alperin weight condition},
Bull.\ London Math.\ Soc.\ \textbf{54} (2022), 466--481.
}

\DeclareManualReference{CabanesEnguehard1994}{%
M. Cabanes and M. Enguehard,
{\it On unipotent blocks and their ordinary characters},
Invent.\ Math.\ \textbf{117} (1994), 149--164.
}

\DeclareManualReference{CabanesEnguehard2004}{%
M. Cabanes and M. Enguehard,
{\it Representation Theory of Finite Reductive Groups},
New Mathematical Monographs, vol.~1, Cambridge Univ.\ Press, Cambridge,
2004.
}

\DeclareManualReference{CabanesSpath2013}{%
M. Cabanes and B. Sp\"ath,
{\it Equivariance and extendibility in finite reductive groups with
connected center},
Math.\ Z.\ \textbf{275} (2013), 689--713.
}

\DeclareManualReference{CabanesSpath2017TypeC}{%
M. Cabanes and B. Sp\"ath,
{\it Inductive McKay condition for finite simple groups of type $\mathsf C$},
Represent.\ Theory \textbf{21} (2017), 61--81.
}

\DeclareManualReference{Chaneb2021}{%
R. Chaneb,
{\it Basic sets for unipotent blocks of finite reductive groups in bad
characteristic},
Int. Math. Res. Not. IMRN \textbf{2021} (2021), no.~16, 12613--12638.
}

\DeclareManualReference{Conlon1968}{%
S.~B. Conlon,
{\it Decompositions induced from the Burnside algebra},
J.\ Algebra \textbf{10} (1968), 102--122.
}

\DeclareManualReference{Atlas1985}{%
J.~H. Conway, R.~T. Curtis, S.~P. Norton, R.~A. Parker, and R.~A. Wilson,
{\it Atlas of Finite Groups},
Clarendon Press, Oxford, 1985.
}

\DeclareManualReference{FengLiZhang2021Equivariant}{%
Z. Feng, C. Li, and J. Zhang,
{\it Equivariant correspondences and the inductive Alperin weight condition
for type $\mathsf A$},
Trans.\ Amer.\ Math.\ Soc.\ \textbf{374} (2021), 8365--8433.
}

\DeclareManualReference{FengLiZhang2022TypeB}{%
Z. Feng, C. Li, and J. Zhang,
{\it Inductive blockwise Alperin weight condition for type $\mathsf B$ and odd
primes},
J.\ Algebra \textbf{604} (2022), 533--576.
}

\DeclareManualReference{FengLiZhang2023LowRankA}{%
Z. Feng, C. Li, and J. Zhang,
{\it On the inductive blockwise Alperin weight condition for type $\mathsf A$},
J.\ Algebra \textbf{631} (2023), 287--305.
}

\DeclareManualReference{FengLiZhang2019}{%
Z. Feng, Z. Li, and J. Zhang,
{\it On the inductive blockwise Alperin weight condition for classical
groups},
J.\ Algebra \textbf{537} (2019), 381--434.
}

\DeclareManualReference{FengLiZhang2022Jordan}{%
Z. Feng, Z. Li, and J. Zhang,
{\it Jordan decomposition for weights and the blockwise Alperin weight
conjecture},
Adv.\ Math.\ \textbf{408} (2022), Art.~108594.
}

\DeclareManualReference{FengLiZhang2023Morita}{%
Z. Feng, Z. Li, and J. Zhang,
{\it Morita equivalences and the inductive blockwise Alperin weight
condition for type $\mathsf A$},
Trans.\ Amer.\ Math.\ Soc.\ \textbf{376} (2023), no.~9, 6497--6520.
}

\DeclareManualReference{FengMalle2022}{%
Z. Feng and G. Malle,
{\it The inductive blockwise Alperin weight condition for type $\mathsf C$ and the
prime 2},
J.\ Aust.\ Math.\ Soc.\ \textbf{113} (2022), 1--20.
}

\DeclareManualReference{FengMalleZhang2026}{%
Z. Feng, G. Malle, and J. Zhang,
{\it Generic weights for finite reductive groups},
Math.\ Ann.\ \textbf{394} (2026), Art.~90.
}

\DeclareManualReference{FengSpath2023}{%
Z. Feng and B. Sp\"ath,
{\it Unitriangular basic sets, Brauer characters and coprime actions},
Represent.\ Theory \textbf{27} (2023), 115--148.
}

\DeclareManualReference{FengYuZhang2024}{%
Z. Feng, J. Yu, and J. Zhang,
{\it Radical subgroups of finite reductive groups},
arXiv:2401.00156v4, 5 November 2024.
}

\DeclareManualReference{FongSrinivasan1989}{%
P. Fong and B. Srinivasan,
{\it The blocks of finite classical groups},
J.\ Reine Angew.\ Math.\ \textbf{396} (1989), 122--191.
}

\DeclareManualReference{GAP2026}{%
The \GAP~Group,
{\it \textup{\GAP} -- Groups, Algorithms, and Programming, Version 4.16.0},
2026, available at
\href{https://www.gap-system.org}{www.gap-system.org}.
}

\DeclareManualReference{Geck1993}{%
M. Geck,
{\it Basic sets of Brauer characters of finite groups of Lie type II},
J.\ London Math.\ Soc.\ (2) \textbf{47} (1993), 255--268.
}

\DeclareManualReference{Geck1994BasicIII}{%
M. Geck,
{\it Basic sets of Brauer characters of finite groups of Lie type, III},
Manuscripta Math.\ \textbf{85} (1994), 195--216.
}

\DeclareManualReference{Geck1999GGGR}{%
M. Geck,
{\it Character sheaves and generalized Gelfand--Graev characters},
Proc.\ London Math.\ Soc.\ (3) \textbf{78} (1999), 139--166.
}

\DeclareManualReference{GeckHezard2008}{%
M. Geck and D. H\'ezard,
{\it On the unipotent support of character sheaves},
Osaka J. Math.\ \textbf{45} (2008), 819--831.
}

\DeclareManualReference{GeckHiss1991}{%
M. Geck and G. Hiss,
{\it Basic sets of Brauer characters of finite groups of Lie type},
J.\ Reine Angew.\ Math.\ \textbf{418} (1991), 173--188.
}

\DeclareManualReference{GeckMalle2000}{%
M. Geck and G. Malle,
{\it On the existence of a unipotent support for the irreducible characters
of a finite group of Lie type},
Trans.\ Amer.\ Math.\ Soc.\ \textbf{352} (2000), 429--456.
}

\DeclareManualReference{GeckMalle2020}{%
M. Geck and G. Malle,
{\it The Character Theory of Finite Groups of Lie Type: A Guided Tour},
Cambridge Stud.\ Adv.\ Math., vol.~187,
Cambridge Univ.\ Press, Cambridge, 2020.
}

\DeclareManualReference{GruberHiss1997}{%
J. Gruber and G. Hiss,
{\it Decomposition numbers of finite classical groups for linear primes},
J.\ Reine Angew.\ Math.\ \textbf{485} (1997), 55--91.
}

\DeclareManualReference{Isaacs1976}{%
I.~M. Isaacs,
{\it Character Theory of Finite Groups},
Pure Appl. Math., vol.~69,
Academic Press, New York, 1976.
}

\DeclareManualReference{KessarMalle2013}{%
R. Kessar and G. Malle,
{\it Quasi-isolated blocks and Brauer's height zero conjecture},
Ann. of Math. (2) \textbf{178} (2013), 321--384.
}

\DeclareManualReference{KessarMalle2015}{%
R. Kessar and G. Malle,
{\it Lusztig induction and $\ell$-blocks of finite reductive groups},
Pacific J.\ Math.\ \textbf{279} (2015), 269--298.
}

\DeclareManualReference{KoshitaniSpath2016}{%
S. Koshitani and B. Sp\"ath,
{\it The inductive Alperin--McKay and blockwise Alperin weight conditions
for blocks with cyclic defect groups and odd primes},
J.\ Group Theory \textbf{19} (2016), 777--813.
}

\DeclareManualReference{Letellier2005}{%
E. Letellier,
{\it Fourier Transforms of Invariant Functions on Finite Reductive Lie Algebras},
Lecture Notes in Math., vol.~1859,
Springer, Berlin, 2005.
}

\DeclareManualReference{Li2021}{%
C. Li,
{\it The inductive blockwise Alperin weight condition for
$\operatorname{PSp}_{2n}(q)$ and odd primes},
J.\ Algebra \textbf{567} (2021), 582--612.
}

\DeclareManualReference{LiLi2019}{%
C. Li and Z. Li,
{\it The inductive blockwise Alperin weight condition for
$\operatorname{PSp}_4(q)$},
Algebra Colloq.\ \textbf{26} (2019), 361--386.
}

\DeclareManualReference{Lusztig1992Support}{%
G. Lusztig,
{\it A unipotent support for irreducible representations},
Adv.\ Math.\ \textbf{94} (1992), 139--179.
}

\DeclareManualReference{Macgregor2022}{%
N. Macgregor,
{\it Morita equivalence classes of tame blocks of finite groups},
J.\ Algebra \textbf{608} (2022), 719--754.
}

\DeclareManualReference{MalleTesterman2011}{%
G. Malle and D. Testerman,
{\it Linear Algebraic Groups and Finite Groups of Lie Type},
Cambridge Stud.\ Adv.\ Math., vol.~133,
Cambridge Univ.\ Press, Cambridge, 2011.
}

\DeclareManualReference{MartinezRizoRossi2026}{%
J.~M. Mart\'{\i}nez, N. Rizo, and D. Rossi,
{\it The Alperin weight conjecture and the Glauberman correspondence via
character triples},
Algebra Number Theory \textbf{20} (2026), 333--382.
}

\DeclareManualReference{Navarro1998}{%
G. Navarro,
{\it Characters and Blocks of Finite Groups},
London Math.\ Soc.\ Lecture Note Ser., vol.~250,
Cambridge Univ.\ Press, Cambridge, 1998.
}

\DeclareManualReference{NavarroTiep2011}{%
G. Navarro and P.~H. Tiep,
{\it A reduction theorem for the Alperin weight conjecture},
Invent.\ Math.\ \textbf{184} (2011), 529--565.
}

\DeclareManualReference{Ruhstorfer2022Derived}{%
L. Ruhstorfer,
{\it Derived equivalences and equivariant Jordan decomposition},
Represent. Theory \textbf{26} (2022), 542--584.
}

\DeclareManualReference{Sambale2012}{%
B. Sambale,
{\it Blocks with defect group $D_{2^n}\times C_{2^m}$},
J. Pure Appl. Algebra \textbf{216} (2012), 119--125.
}

\DeclareManualReference{Sambale2014}{%
B. Sambale,
{\it Blocks of Finite Groups and Their Invariants},
Lecture Notes in Math., vol.~2127, Springer, Cham, 2014.
Online version: \url{https://benjaminsambale.github.io/pdfs/lnm.pdf} (accessed 19 September 2026).
}

\DeclareManualReference{SchaefferFry2014}{%
A.~A. Schaeffer Fry,
{\it $\operatorname{Sp}_6(2^a)$ is ``good'' for the McKay, Alperin weight,
and related local-global conjectures},
J.\ Algebra \textbf{401} (2014), 13--47.
}

\DeclareManualReference{Shoji1995II}{%
T. Shoji,
{\it Character sheaves and almost characters of reductive groups, II},
Adv.\ Math.\ \textbf{111} (1995), 314--354.
}

\DeclareManualReference{Spath2013}{%
B. Sp\"ath,
{\it A reduction theorem for the blockwise Alperin weight conjecture},
J.\ Group Theory \textbf{16} (2013), 159--220.
}

\DeclareManualReference{Spath2017Triples}{%
B. Sp\"ath,
{\it Inductive conditions for counting conjectures via character triples},
in {\it Representation Theory---Current Trends and Perspectives},
EMS Ser. Congr. Rep., Eur. Math. Soc., Z\"urich, 2017, 665--680.
}

\DeclareManualReference{Spath2025Extensions}{%
B. Sp\"ath,
{\it Extensions of characters in type $\mathsf D$ and the inductive McKay
condition, II},
Invent.\ Math.\ \textbf{242} (2025), 45--122.
}

\DeclareManualReference{Taylor2013}{%
J. Taylor,
{\it On unipotent supports of reductive groups with a disconnected centre},
J.\ Algebra \textbf{391} (2013), 41--61.
}

\DeclareManualReference{Taylor2014Maximal}{%
J. Taylor,
{\it Finding characters satisfying a maximal condition for their
unipotent support},
J.\ Pure Appl.\ Algebra \textbf{218} (2014), 474--496.
}

\DeclareManualReference{Taylor2016}{%
J. Taylor,
{\it Generalized Gelfand--Graev representations in small characteristics},
Nagoya Math.\ J.\ \textbf{224} (2016), 93--167.
}

\DeclareManualReference{Taylor2018Automorphisms}{%
J. Taylor,
{\it Action of automorphisms on irreducible characters of symplectic
groups},
J.\ Algebra \textbf{505} (2018), 211--246.
}

\DeclareManualReference{AtlasRep2026}{%
R.~A. Wilson, R.~A. Parker, S. Nickerson, J.~N. Bray, and T. Breuer,
{\it AtlasRep, a \textup{\GAP} interface to the Atlas of group
representations, Version 2.1.11},
2026, available at
\href{https://www.math.rwth-aachen.de/~Thomas.Breuer/atlasrep}
{the AtlasRep website}.
}

\DeclareManualReference{Benson1998}{%
D.~J. Benson,
{\it Representations and Cohomology, I: Basic Representation Theory of
Finite Groups and Associative Algebras}, paperback ed.,
Cambridge Studies in Advanced Mathematics, vol.~30,
Cambridge Univ. Press, Cambridge, 1998.
}

\DeclareManualReference{Serre1977}{%
J.-P. Serre,
{\it Linear Representations of Finite Groups},
translated from the French by L.~L. Scott,
Graduate Texts in Mathematics, vol.~42,
Springer-Verlag, New York--Heidelberg, 1977.
}



\hypersetup{pdftitle={ModularRep: Mathematical Library Manual}}
\begin{document}
\frontmatter
\begin{titlepage}
\centering
\vspace*{28mm}
{\Large\sffamily ModularRep\par}
\vspace{12mm}
{\Huge Mathematical Library Manual\par}
\vspace{10mm}
{\large Representation theory in Lean\par}
\vfill
\begin{minipage}{0.82\textwidth}
This manual develops the theory of finite group representations, ordinary
and Brauer characters, reduction, blocks and weights. It explains the mathematical
constructions and their proofs, with references to the corresponding
Lean definitions and theorems.
\end{minipage}
\vspace{18mm}
\end{titlepage}
\tableofcontents
\chapter*{Using this manual}
\addcontentsline{toc}{chapter}{Using this manual}
The chapters follow the constructions needed to pass from representations
to local weights. Normal subgroups and quotient groups come first, followed
by ordinary characters, Brauer characters and exact Grothendieck groups.
These provide the setting for reduction, blocks and weights. The remaining
chapters study equivariant bijections, character counts from finite
matrices, and several general arguments involving bases, stabilisers and
rational series.

Each mathematical statement specifies its coefficient fields and any
required assumptions on splitting fields, roots of unity and compatibility. The exposition
includes standard results used by the library. Implementation paragraphs
distinguish results imported from Mathlib, definitions and deductions in
this library, and results supplied as explicit hypotheses to the Lean code.
References support the statements at the indicated locations, with version
distinctions where these affect the formulation.

\texttt{ModularRep.Library} imports a selection of basic constructions.
For a particular result, use the module named in its implementation paragraph. The accompanying
documentation index lists the source modules and their declarations.
Instructions for compiling the library and checking the proofs are supplied
with the package.
Unless explicitly qualified, theorem and equation numbers refer to this
manual.

The general library and the manuscript applications have separate
hypotheses. The formalisation report describes the assumptions used by
the applications. The Lean code does not construct an instance satisfying all the hypotheses
of the manuscript applications.

\mainmatter
\chapter{Conventions and constructions}
\label{chap:library}

The manual treats ordinary characters, Brauer characters, reduction and
blocks, and the weights defined from ordinary characters of normaliser
quotients.

Three kinds of construction recur in these subjects: passage to a normal
subgroup or quotient, passage from representations to their characters or
Grothendieck classes, and transport under a group action. The chapters
develop these constructions together with the compatibility statements
needed to combine them.

Throughout the manual, groups are finite unless a statement says otherwise.
The modular prime is $\ell$, $k$ has characteristic $\ell$, and $K$ has
characteristic zero. A discrete valuation ring is a local principal ideal
domain that is not a field. Coefficient rings, splitting conditions,
algebraic closure and choices of roots are specified where needed. The
source files often use the letter $p$ for the modular prime. In the Conlon
argument, $p$ denotes a potentially different prime used for permutation
lattices.

\section{From a result to its source}
\label{chap:library:use}

Mathlib supplies the representation types \lean{Representation} and
\lean{FDRep}, the ordinary trace character \lean{Representation.character},
and the underlying group theory, linear algebra and category theory.
ModularRep defines the finite root correspondences, Brauer characters and
exact Grothendieck group used here. The implementation paragraphs explain
which steps use Mathlib results and identify the definitions and proofs
in the modules of this package.

The accompanying module index, linked from the package reading guide,
lists the library and formalisation modules, their declarations and stated
assumptions. Its entries are classified by role and status.
Examples in this manual illustrate the stated mathematical
results.

For local quotient arguments, one can begin with
\begin{samepage}
\begin{verbatim}
import ModularRep.PCore
import ModularRep.NormalCoreTransport
\end{verbatim}
\end{samepage}
The normaliser quotient results apply to a surjective homomorphism whose
kernel lies in the relevant subgroup. The induced isomorphism of
normaliser quotients determines the pullback of representations.
For reduction, the imports
\begin{samepage}
\begin{verbatim}
import ModularRep.BrauerDecompositionMap
import ModularRep.KZeroTwist
\end{verbatim}
\end{samepage}
provide the decomposition homomorphism for the chosen modular system and
root correspondence. Its hypotheses relate reduction of a stable lattice
to Brauer characters. The naturality theorems use the same choices. The modules
\begin{samepage}
\begin{verbatim}
import Formalisation.FibreTransport
import Formalisation.IndexedAssembly
\end{verbatim}
\end{samepage}
construct equivariant bijections of finite sets from equivariant
bijections of their fibres. The chapters state the hypotheses of each
construction.

\section{Hypotheses and proofs}

The mathematical exposition includes standard results used by the library.
An implementation paragraph says whether the corresponding result is
proved in Lean or supplied as an input. In particular,
Hypotheses~\ref{lib:blocks:catalogue} and
\ref{lib:blocks:defect-hypotheses} collect standard published facts
used as inputs to the Lean arguments. Other hypotheses express
compatibility of chosen data, as in
Hypothesis~\ref{lib:reduction:character-hypothesis}.
For example, multiplicativity of
the central Brauer map follows from cancellation of nontrivial
$\ell$-group orbits. The construction of the decomposition homomorphism
uses an explicit hypothesis relating stable reduction to Brauer
characters. The existence of a stable lattice and the deductions from
that compatibility hypothesis are proved separately.

An equivariant bijection of finite sets does not by itself preserve a
specified relation. To construct a bijection preserving a relation from
bijections of the fibres, the latter must preserve the relation and the
relation must be invariant under the acting group. Character and block
relations require their corresponding compatibility hypotheses.

Finite computations also require a mathematical interpretation. A matrix
rank calculation, for example, counts irreducible Brauer characters only
after its rows have been identified with character restrictions and the
required spanning and independence statements have been established.
Chapter~\ref{chap:library-computations} states these hypotheses and proves
the resulting rank formulas.

\chapter{Normal subgroups and local quotients}
\label{chap:library:groups}

Throughout this chapter, $G$ and $H$ are finite groups and $\ell$ is a prime.
For a subgroup $Q\leq G$, write $N_G(Q)$ for its normaliser. Since $Q$ is
normal in $N_G(Q)$, the quotient $N_G(Q)/Q$ is defined. We study these
quotients and radical subgroups under isomorphisms and surjective
homomorphisms. They will occur in the definition of a weight in
Chapter~\ref{chap:library:weights}.

\section{The largest normal \texorpdfstring{$\ell$}{ell}-subgroup}

\begin{definition}
The \emph{$\ell$-core} $O_\ell(G)$ is the largest normal $\ell$-subgroup
of $G$. A subgroup $Q\leq G$ is \emph{$\ell$-radical} if
$Q=O_\ell(N_G(Q))$.
\end{definition}

To see that the core exists, take the subgroup generated by all normal
$\ell$-subgroups of $G$. Indeed,
the product of two normal $\ell$-subgroups is normal and has order a power
of $\ell$. Finiteness reduces the generated subgroup to such a product.
It contains every normal $\ell$-subgroup and is characterised by this
property. An automorphism permutes the normal $\ell$-subgroups, so it
preserves the core. More generally, a group isomorphism $\phi:G\to H$ satisfies
$\phi(O_\ell(G))=O_\ell(H)$.

In particular, a radical subgroup is an $\ell$-group. If $\phi:G\to H$ is
an isomorphism, then $\phi(N_G(Q))=N_H(\phi(Q))$, so the equality of cores
shows that $Q$ is radical if and only if $\phi(Q)$ is radical.
\begin{implementation}
\module{ModularRep.PCore} defines the core using Mathlib's subgroup lattice
and its theorem on suprema of normal $\ell$-subgroups.
The basic core results include
\lean{ModularRep.pCore_isPGroup}, \lean{ModularRep.pCore_normal} and
\lean{ModularRep.normal_pSubgroup_le_pCore}.
\module{ModularRep.RadicalSubgroup} defines radicality by mapping the core
of $N_G(Q)$ into $G$ along the normaliser inclusion.
\module{ModularRep.RadicalTransport} proves invariance under isomorphisms.
\end{implementation}

\begin{proposition}[The core lies in every radical subgroup]
\label{lib:groups:core-contained}
Every $\ell$-radical subgroup of $G$ contains $O_\ell(G)$.
\end{proposition}
\begin{proof}
Let $P=O_\ell(G)$ and let $Q$ be radical. Both are $\ell$-groups, and
normality of $P$ makes $PQ$ an $\ell$-group. Suppose that $Q<PQ$.
The normaliser condition in a finite
$\ell$-group gives $x\in N_{PQ}(Q)\setminus Q$. Write $x=yz$, with
$y\in P$ and $z\in Q$. Both $x$ and $z$ normalise $Q$, hence
$y\in P\cap N_G(Q)$. This intersection is a normal $\ell$-subgroup of
$N_G(Q)$ and therefore lies in $O_\ell(N_G(Q))=Q$. Thus $x=yz\in Q$,
a contradiction. It follows that $PQ=Q$ and $P\leq Q$.
\end{proof}
\begin{implementation}
The theorem is \lean{ModularRep.pCore_le_of_isRadicalSubgroup} in
\module{ModularRep.NormalCoreTransport}. It uses Mathlib's nilpotence theorem
for finite $\ell$-groups and its normaliser condition for nilpotent groups.
\end{implementation}

\section{Surjections and normaliser quotients}

\begin{proposition}[Normaliser quotients under a surjection]
\label{lib:groups:normaliser-quotient}
Let $f:G\twoheadrightarrow H$ be surjective, and let $Q\leq G$ contain
$\ker f$. Then $f(N_G(Q))=N_H(f(Q))$, and the map $gQ\mapsto f(g)f(Q)$
induces an isomorphism
\[
 N_G(Q)/Q\;\simeq\;N_H(f(Q))/f(Q).
\]
\end{proposition}
\begin{proof}
The inclusion $f(N_G(Q))\leq N_H(f(Q))$ follows by applying $f$ to a
conjugate of $Q$. Conversely, choose $h\in N_H(f(Q))$ and lift it to
$g\in G$. Since $Q$ contains the kernel, $Q=f^{-1}(f(Q))$. Conjugation
by $g$ therefore preserves $Q$, proving normaliser surjectivity.

Compose this normaliser map with the quotient by $f(Q)$. An element $g$
of $N_G(Q)$ is in its kernel exactly when $f(g)\in f(Q)$, equivalently
when $g\in Q$. The first isomorphism theorem gives the assertion, with
the stated formula on cosets.
\end{proof}

Thus containment of $\ker f$ in $Q$ suffices to identify the normaliser
quotients. To preserve radicality as well, the following result assumes
that $\ker f$ is an $\ell$-group.

\begin{proposition}[Radical subgroups under a surjection]
\label{lib:groups:radical-quotient}
Suppose that $f:G\twoheadrightarrow H$ has $\ell$-group kernel. Then
$f(O_\ell(G))=O_\ell(H)$. For every $Q\leq G$ containing $\ker f$,
$Q$ is $\ell$-radical if and only if $f(Q)$ is $\ell$-radical.
Consequently image and inverse image give inverse correspondences between
the radical subgroups of $G$ and those of $H$.
\end{proposition}
\begin{proof}
The image of the core is a normal $\ell$-subgroup, so
$f(O_\ell(G))\leq O_\ell(H)$. The inverse image of $O_\ell(H)$ is
an extension of two $\ell$-groups and is normal in $G$. It lies in
$O_\ell(G)$, proving equality.

The restriction of $f$ to $N_G(Q)$ is surjective onto $N_H(f(Q))$ by
Proposition~\ref{lib:groups:normaliser-quotient}, and its kernel is an
$\ell$-group. Applying the core equality to this restriction identifies
the images of the two normaliser cores. Both $Q$ and $O_\ell(N_G(Q))$
contain $\ker f$. Image under $f$ is injective on subgroups containing
the kernel, so equality of their images is equivalent to equality of these
subgroups. This proves the radical equivalence.

Finally, $\ker f\leq O_\ell(G)$, so every radical subgroup of $G$
contains $\ker f$ by Proposition~\ref{lib:groups:core-contained}.
The subgroup correspondence theorem now gives the two inverse maps.
\end{proof}
\begin{implementation}
\module{ModularRep.NormalCoreTransport} proves the core image equality,
normaliser surjectivity and the radical equivalence. Its declarations are
\lean{pCore_map_surjective_of_ker_isPGroup},
\lean{normalizerMap_surjective},
\lean{isRadicalSubgroup_iff_map_surjective_of_ker_le} and
\lean{normalizerQuotientEquivOfSurjectiveOfKerLE}, all in
\lean{ModularRep}. These constructions use Mathlib's subgroup
correspondence and quotient homomorphisms.
\end{implementation}

\begin{example}[Quotient by the core]
Set $P=O_\ell(G)$. The quotient $G/P$ has trivial $\ell$-core, and its
radical subgroups correspond to $Q/P$ for radical subgroups $Q$ of $G$.
For each such $Q$, the normaliser quotient is unchanged up to the isomorphism
of Proposition~\ref{lib:groups:normaliser-quotient}. A representation of
this local quotient can therefore be transported without changing its
dimension or irreducibility. Chapter~\ref{chap:library:weights} explains
the resulting construction on weights. If $G$ is itself an $\ell$-group,
Proposition~\ref{lib:groups:core-contained} shows that its only radical
subgroup is $G$.
\end{example}

\section{Central quotients}

The preceding quotient results apply to subgroups containing the kernel.
A central kernel of order prime to $\ell$ gives a different subgroup
correspondence.

\begin{proposition}[Quotients by central $\ell'$-subgroups]
\label{lib:groups:central-prime-to}
Let $f:G\to H$ be a surjective homomorphism of finite groups.
Suppose that $Z=\ker f$ is central and $\ell\nmid|Z|$.
Then $Q\mapsto f(Q)$ is a bijection between the $\ell$-subgroups
of $G$ and those of $H$. For each $\ell$-subgroup $R\leq H$, the unique
$\ell$-subgroup mapping onto $R$ is $O_\ell(f^{-1}(R))$.
For every $\ell$-subgroup $Q\leq G$,
\[
 f^{-1}\bigl(N_H(f(Q))\bigr)=N_G(Q),\qquad
 Q\text{ is radical}\ \Longleftrightarrow\ f(Q)\text{ is radical}.
\]
\end{proposition}
\begin{proof}
Let $R\leq H$ be an $\ell$-subgroup. A Sylow $\ell$-subgroup $Q$ of
$f^{-1}(R)$ has order $|R|$, since the kernel has order prime to
$\ell$. Since $Q\cap Z=1$, its image is $R$, and
$f^{-1}(R)=QZ$. Centrality of $Z$ makes $Q$ normal in this
inverse image, so $Q$ lies in its $\ell$-core. That core also
maps onto $R$.

The restriction of $f$ to any $\ell$-subgroup is injective,
since its kernel has order dividing both an $\ell$-power and
$|Z|$. If two $\ell$-subgroups have image $R$, they both lie
in the $\ell$-core of $f^{-1}(R)$. Injectivity on that core
then shows that they are equal. This proves the correspondence
and the formula for its inverse.

An element $g\in G$ normalises $Q$ precisely when $f(g)$
normalises $f(Q)$. For the reverse implication, use uniqueness
of the lift on $Q$ and $gQg^{-1}$. This proves the normaliser
formula and surjectivity of $N_G(Q)\to N_H(f(Q))$.

Uniqueness also implies that the lift of a normal
$\ell$-subgroup is normal. Consequently
$f(O_\ell(G))=O_\ell(H)$. Apply the same argument to the
surjection of normalisers. The image of $O_\ell(N_G(Q))$ is
$O_\ell(N_H(f(Q)))$. Injectivity on $\ell$-subgroups now
proves the radicality equivalence.
\end{proof}

In particular, the induced map on local quotients is
\[
 N_G(Q)/Q\longrightarrow N_H(f(Q))/f(Q).
\]
Its kernel is $QZ/Q\simeq Z$, since $Q\cap Z=1$. Quotienting by this
kernel gives
\[
 N_G(Q)/(QZ)\cong N_H(f(Q))/f(Q).
\]
The two local quotient orders therefore have equal $\ell$-parts.
Inflation along this map gives characters of $N_G(Q)/Q$ that are
trivial on $QZ/Q$.

Taking $f$ to be the quotient map by a central $\ell'$-subgroup $Z$
gives, for every $\ell$-radical subgroup $Q$,
\[
 QZ/Z\text{ is radical in }G/Z,
 \qquad N_{G/Z}(QZ/Z)=N_G(Q)/Z.
\]
These conclusions also follow from
\cite[Lemma~2.3(c) and its proof]{NavarroTiep2011}.
Proposition~\ref{lib:groups:normaliser-quotient} does not apply directly
to $Q$, because $Z$ need not be contained in $Q$.

\begin{implementation}
The subgroup correspondence is proved in
\module{ModularRep.PaperProofs.SporadicFi24P3Definition44NamedCarrierCentralPrimeToPSubgroups},
including \lean{existsUnique_pSubgroup_lift} and
\lean{pSubgroupLift_map}. The module whose name ends in
\lean{CentralPrimeToRadicals} proves
\lean{normalizer_comap_of_central_primeTo} and
\lean{radical_iff_map_of_central_primeTo}. The theorem
\lean{fixedCentralQuotientSource} gives the quotient specialisation
as the structure
\lean{ModularRep.NavarroTiep23cFixedCentralQuotientSource}, defined in
\module{ModularRep.PaperProofs.NavarroTiep23cFixedCentralQuotientSource}.
\module{ModularRep.PaperProofs.CentralEllPrimeWeightLocalQuotient}
defines the homomorphism \lean{qW} on normaliser quotients and computes
its kernel. Its theorems on surjectivity and the $\ell$-parts of the
quotient orders are stated for radical $Q$ and use the quotient
specialisation above. The general subgroup correspondence also allows
the trivial subgroup.
\end{implementation}

\section{Cyclic quotients and Sylow subgroups}

\begin{proposition}[Cyclic quotients by the core]
\label{lib:groups:hypoelementary}
If $G/O_\ell(G)$ is cyclic, then $H/O_\ell(H)$ is cyclic for every
subgroup $H\leq G$. If $G$ is cyclic, $G/O_\ell(G)$ has order prime
to $\ell$.
\end{proposition}
\begin{proof}
Let $I=H\cap O_\ell(G)$. This is a normal $\ell$-subgroup of $H$,
so $I\leq O_\ell(H)$. Restriction of the quotient map embeds $H/I$
in $G/O_\ell(G)$. Thus $H/I$ is cyclic, as is its quotient
$H/O_\ell(H)$. For the second assertion, the unique Sylow $\ell$-subgroup
of a finite cyclic group is normal and equals its $\ell$-core. Its index
is prime to $\ell$.
\end{proof}

For a finite group the condition that $G/O_\ell(G)$ be cyclic is
equivalent to $\ell$-hypoelementarity, that is, the existence of a normal
$\ell$-subgroup with cyclic quotient of order prime to $\ell$.
For this terminology, see, for example,
\cite[Definition~5.5.3]{Benson1998}. Indeed, the inverse image of the
Sylow $\ell$-subgroup of a cyclic quotient would be a normal
$\ell$-subgroup of $G$. Maximality of the core forces that quotient
Sylow subgroup to be trivial. Thus the cyclic quotient has $\ell'$-order.
The prime defining hypoelementarity is specified independently of any
coefficient field.
\begin{implementation}
\lean{ModularRep.subgroup_quotient_pCore_isCyclic} and
\lean{ModularRep.not_dvd_card_quotient_pCore_of_isCyclic}.
The latter belongs to \module{ModularRep.CyclicPCore}.
\end{implementation}

\begin{proposition}[A criterion for cyclic Sylow subgroups]
\label{lib:groups:cyclic-sylow}
Let $T\leq G$ be cyclic. If $|G|_\ell$ divides $|T|$, then $G$ has
a cyclic Sylow $\ell$-subgroup, and every $\ell$-subgroup of $G$ is cyclic.
Here $|G|_\ell$ denotes the largest power of $\ell$ dividing $|G|$.
\end{proposition}
\begin{proof}
The cyclic group $T$ contains a subgroup of order $|G|_\ell$. It is a
Sylow subgroup of $G$ by its order. Every $\ell$-subgroup of $G$ is
contained in a Sylow subgroup, and all Sylow subgroups are conjugate.
Hence each is contained in a cyclic group and is cyclic.
\end{proof}
\begin{implementation}
\module{ModularRep.CyclicSylow}, in particular
\lean{exists_cyclic_sylow_of_full_part_dvd_card_cyclic_subgroup} and
\lean{isCyclic_of_isPGroup_of_full_part_dvd_card_cyclic_subgroup}
in that module's namespace. They combine Mathlib's Sylow existence,
containment and conjugacy theorems with transport of cyclicity.
\end{implementation}

\begin{example}[Applying the order criterion]
Suppose $G$ has a cyclic subgroup $T$ and an order
factorisation $|G|=ab$ with $\ell\nmid a$ and $b\mid |T|$.
Then $|G|_\ell\mid b$, so the preceding proposition gives cyclicity of
every $\ell$-subgroup of $G$. This divisibility is proved by
\lean{ModularRep.CyclicSylow.ordProj_mul_dvd_right_of_not_dvd_left}.
\end{example}

\chapter{Ordinary representations and characters}
\label{chap:library:characters}

Pullback, inflation and tensoring with a linear character relate
representations of different groups or give new representations of the
same group. Their character formulas follow from the representing
operators. We first show how characters determine irreducible
representations, then study these operations and extension across a
cyclic quotient.

Throughout, $K$ is a field of characteristic zero and $G$ is a finite group.
A representation $\rho$ of $G$ on a finite dimensional $K$-vector space $V$
has character $\chi_\rho(g)=\tr_K(\rho(g))$. We state algebraic closure
explicitly where it is used. Character separation for finite groups
also holds over an arbitrary field of characteristic zero, as the
averaging argument below shows. Modular systems and reduction are
introduced in Chapter~\ref{chap:library:reduction}.

\section{Representations and their characters}
\label{lib:ordinary:realisations}

\begin{definition}
A representation of $G$ on $V$ is a homomorphism
$\rho:G\to\GL_K(V)$. It is irreducible if $V\ne0$ and its only invariant
subspaces are $0$ and $V$. An equivalence $\rho\simeq\sigma$ is a linear
isomorphism $T:V\to W$ satisfying
\[
 T\rho(g)=\sigma(g)T\qquad(g\in G).
\]
The set $\Irr_K(G)$ consists of the functions $\chi:G\to K$ afforded by
finite dimensional irreducible $K$-representations.
\end{definition}

Thus equality in $\Irr_K(G)$ is equality at every group element. Choosing
different bases or isomorphic realisations does not create new characters:
conjugate operators have equal traces. The value at the identity is
$\chi_\rho(1)=\dim_K V$, regarded as an element of $K$. Every character is
constant on conjugacy classes, since $\rho(xgx^{-1})$ is conjugate to
$\rho(g)$.

\begin{implementation}
Mathlib supplies \lean{Representation},
\lean{Representation.IsIrreducible} and the trace function
\lean{Representation.character}, together with the identities
\lean{Representation.char_one}, \lean{Representation.char_conj} and
\lean{Representation.char_iso}.
The additional type \lean{ModularRep.OrdinaryIrreducibleCharacter.Irr}
collects the functions admitting an irreducible realisation on some
$K^n$, expressed by
\lean{ModularRep.OrdinaryIrreducibleCharacter.Realisation} using these
Mathlib definitions. Equality of these character functions is pointwise,
as expressed by
\lean{ModularRep.OrdinaryIrreducibleCharacter.ext}.
\end{implementation}

\begin{proposition}[Characters determine irreducible representations]
\label{lib:ordinary:separation}
Assume that $K$ is algebraically closed. If two finite dimensional
irreducible $K$-representations $\rho$ and $\sigma$ have equal characters,
then $\rho\simeq\sigma$.
\end{proposition}
For the corresponding statement about traces of simple modules over a
finite dimensional algebra, see, for example,
\cite[Theorem~1.19]{Navarro1998}.
\begin{proof}
Regard the two representation spaces as simple $KG$-modules. Equality of
their traces on $G$ implies equality on every element of $KG$ by linearity.
Suppose that the simple modules are nonisomorphic. Schur's lemma says that
their mutual homomorphism spaces are zero and their endomorphisms are
scalars. The density theorem applied to their direct sum therefore supplies
an element of $KG$ acting as an operator of trace one on the first module
and as zero on the second. This contradicts equality of traces.

An operator of trace one can be obtained by choosing a basis and taking the
projection onto one basis vector. Using this operator also makes the same
separation argument valid in positive characteristic, where the trace of
the identity may be zero.
\end{proof}

\begin{implementation}
The theorem is
\lean{ModularRep.ModularTraceSeparation.representation_equiv_of_character_eq}.
Mathlib supplies the Schur and density statements used in the proof,
\lean{IsSimpleModule.algebraMap_end_bijective_of_isAlgClosed} and
\lean{Module.Finite.toModuleEnd_moduleEnd_surjective}.
The local lemmas
\lean{exists_trace_action_ne_of_not_linearEquiv} and
\lean{linearMap_trace_surjective}, in the same namespace, make the
separating operator explicit. The theorem itself applies to a group over
an algebraically closed field of any characteristic and does not require
the group to be finite. Its use here is in characteristic zero.
\end{implementation}

For a finite group, equality of irreducible characters also implies
equivalence over an arbitrary field $K$ of characteristic zero. The
character averaging formula (see, for example,
\cite[Section~12.1, proof of Proposition~32]{Serre1977}) gives
\[
 \dim_K\operatorname{Hom}_{KG}(V,W)
   =\frac{1}{|G|}\sum_{g\in G}\chi_V(g^{-1})\chi_W(g),
\]
where the integer on the left is viewed in $K$. If $\chi_V=\chi_W$,
the right hand side is also $\dim_K\operatorname{End}_{KG}(V)$.
This dimension is positive because $V\ne0$ and its identity map is
nonzero. Thus there is a nonzero homomorphism $V\to W$, which is an
isomorphism by irreducibility. The proof does not require that the
endomorphism ring consist of scalars.

\begin{implementation}
The theorem
\lean{ModularRep.PaperProofs.TypeBOrdinaryTraceSeparationSplitting.nonempty_equiv_of_character_eq}
proves this implication using Mathlib's character averaging formula. It applies
to every finite group over every field of characteristic zero and also
identifies simple module classes with equal characters. In particular,
it applies when $K$ is the fraction field of a discrete valuation ring.
\end{implementation}

\section{Pullback and quotient descent}
\label{lib:ordinary:pullback}

For a homomorphism $f:H\to G$, the pullback $f^*\rho$ acts on the same
vector space by $(f^*\rho)(h)=\rho(f(h))$. Its character is
$\chi_\rho\circ f$. An intertwiner between two representations of $G$
also intertwines their pullbacks.

\begin{proposition}[Inflation and descent of representations]
\label{lib:ordinary:surjective-pullback}
If $f:H\twoheadrightarrow G$ is surjective, then $f^*\rho$ is irreducible
if and only if $\rho$ is irreducible. If $\pi:G\twoheadrightarrow H$ and
$\ker \pi\leq\ker\rho$, there is a unique representation $\bar\rho$ of $H$
on $V$ such that $\pi^*\bar\rho=\rho$. It is irreducible whenever $\rho$ is.
\end{proposition}
\begin{proof}
A subspace invariant under $\rho$ is invariant under its pullback. Conversely,
surjectivity lets every $g\in G$ be written as $f(h)$, so an invariant
subspace for the pullback is invariant under $\rho$. The lattices of
invariant subspaces are identical.

For descent, define $\bar\rho(\pi(g))=\rho(g)$. If $\pi(g)=\pi(g')$, then
$g'^{-1}g\in\ker \pi\leq\ker\rho$, which makes the value independent of the
lift. This defines a representation and proves the pullback identity.
Surjectivity proves uniqueness, and the first assertion proves
irreducibility.
\end{proof}

\begin{implementation}
\module{ModularRep.Twist} constructs the order isomorphism
\lean{Representation.subrepresentationPullbackOrderIso}, from which
\lean{Representation.isIrreducible_pullback_iff} follows. These results
hold over any field, and the order isomorphism is available more generally
over a semiring. In
\module{ModularRep.RepresentationSurjectiveDescent},
\lean{descend} factors through $G/\ker \pi$ and uses Mathlib's quotient
isomorphism $G/\ker \pi\simeq H$. The declarations
\lean{descend_apply}, \lean{descend_pullback} and
\lean{descend_irreducible} give its defining identities and irreducibility.
Uniqueness on the fixed vector space is the direct surjectivity argument
in the proof above.
\end{implementation}

To descend a character, it remains to show that the kernel of the quotient
map acts trivially in an affording representation. Over $\mathbb C$, the
standard identity
$\ker\rho=\{g\in G:\chi_\rho(g)=\chi_\rho(1)\}$ gives this directly
(see, for example, \cite[Lemma~2.19]{Isaacs1976}). The averaging argument below proves
the needed implication over any field of characteristic zero, without
assuming irreducibility.

\begin{lemma}[Kernel containment by averaging]
\label{lib:ordinary:averaging}
Let $S$ be a finite subgroup of $G$. If
$\chi_\rho(s)=\chi_\rho(1)$ for every $s\in S$, then $S\leq\ker\rho$.
\end{lemma}
\begin{proof}
The averaging operator
\[
 P_S=\frac1{|S|}\sum_{s\in S}\rho(s)
\]
is a projection onto the invariant subspace $V^S$. Its trace is both
$\dim_K V^S$ and, by the hypothesis, $\dim_K V$. These integers are equal
because their images in a field of characteristic zero are equal. Hence
$V^S=V$, and every element of $S$ acts trivially.
\end{proof}

\begin{implementation}
\lean{Representation.subgroup_le_ker_of_character_eq_one}, in
\module{ModularRep.OrdinaryCharacterLiesOverKernel}, proves this using Mathlib's
\lean{Representation.card_inv_mul_sum_char_eq_finrank} for the restricted
representation. Only $S$ must be finite. No irreducibility, normality or
centrality hypothesis occurs in this lemma.
\end{implementation}

\begin{theorem}[Descent of an ordinary irreducible character]
\label{lib:ordinary:character-descent}
Let $\pi:G\twoheadrightarrow H$ be a surjective homomorphism and let
$\chi\in\Irr_K(G)$. Suppose that $\chi(x)=\chi(1)$ for all
$x\in\ker \pi$. Then there is a unique $\bar\chi\in\Irr_K(H)$ with
\[
 \chi(g)=\bar\chi(\pi(g))\qquad(g\in G).
\]
\end{theorem}
\begin{proof}
Choose an irreducible realisation of $\chi$. Lemma~\ref{lib:ordinary:averaging}
places $\ker \pi$ in its kernel. Descend this representation by
Proposition~\ref{lib:ordinary:surjective-pullback} and take its trace.
The resulting representation remains irreducible. Any two functions with
the required pullback agree at every element of $H$, since $\pi$ is surjective.
\end{proof}

\begin{implementation}
The descended character, its factorisation and its uniqueness are given by
\lean{ModularRep.OrdinaryIrreducibleCharacterSurjectiveDescent.descendCharacter},
\lean{descendCharacter_factorisation} and \lean{descendCharacter_unique}.
The construction applies averaging to an irreducible realisation and
then descends that representation.
\end{implementation}

For a positive integer $a$, write $a_\ell$ for the largest power of the
prime $\ell$ dividing $a$. An ordinary irreducible character $\chi$ has
\emph{defect zero at $\ell$} if $\chi(1)_\ell=|G|_\ell$, where
$\chi(1)$ denotes its positive integral degree.
Over a splitting field this is the usual notion of defect zero. Over an
arbitrary field of characteristic zero, we use the same numerical
condition on the dimension of an irreducible representation in the
following statement.

\begin{corollary}[Preservation of defect zero]
\label{lib:ordinary:defect-descent}
Let $\ell$ be a prime. Under the hypotheses of
Theorem~\ref{lib:ordinary:character-descent}, assume that $\ell\nmid|\ker \pi|$.
If $\chi$ has defect zero at $\ell$, then $\bar\chi$ has defect zero.
\end{corollary}
\begin{proof}
An affording representation $V$ has dimension $\chi(1)$. Descent retains
this vector space and hence the degree.
The equality $|G|=|H|\,|\ker \pi|$ gives $|G|_\ell=|H|_\ell$, so the same
dimension satisfies the equality defining defect zero for $H$.
\end{proof}

\begin{implementation}
\lean{ModularRep.IsDefectZeroOrdinaryCharacter} expresses the dimension
equality for an irreducible representation in Mathlib's \lean{FDRep}
affording the character.
\lean{descendCharacter_defectZero} in the descent namespace applies
averaging to the affording representation and descends it without changing
its dimension.
\end{implementation}

\begin{example}[Projection from a direct product]
Let $G=H\times C_r$, with $\ell\nmid r$, and let $\pi$ be the first
projection. To apply the descent theorem to an ordinary irreducible
character $\chi$, it is enough to verify
$\chi(1,c)=\chi(1,1)$ for every $c\in C_r$. The theorem then constructs
$\bar\chi\in\Irr_K(H)$ and proves
$\chi(h,c)=\bar\chi(h)$ for all $(h,c)$. If the original character has
defect zero at $\ell$, the corollary preserves that property.
\end{example}

\section{Linear characters and automorphisms}
\label{lib:ordinary:actions}

Write $L_K(G)=\operatorname{Hom}(G,K^\times)$ for the group of linear
characters. For $\lambda\in L_K(G)$ define
$(\rho\otimes\lambda)(g)=\lambda(g)\rho(g)$ on $V$. The usual tensor
product with a one dimensional representation is realised on $V$ by
this formula. For an automorphism $\alpha$ of $G$, set
$\rho^\alpha=\rho\circ\alpha$, $\lambda^\alpha=\lambda\circ\alpha$
and $\chi^\alpha=\chi\circ\alpha$.

\begin{proposition}[Tensoring with a linear character]
\label{lib:ordinary:linear-twist}
Tensoring with a linear character preserves and reflects irreducibility,
and
\[
 \chi_{\rho\otimes\lambda}(g)=\lambda(g)\chi_\rho(g).
\]
For $\alpha\in\Aut(G)$, one also has
\[
 (\rho\otimes\lambda)^\alpha
   =\rho^\alpha\otimes\lambda^\alpha,
 \qquad
 (\chi\lambda)^\alpha=\chi^\alpha\lambda^\alpha.
\]
\end{proposition}
\begin{proof}
Multiplication by the nonzero scalar $\lambda(g)$ does not change which
subspaces are invariant under the action of $g$. Multiplication by
$\lambda(g)^{-1}$ gives the converse. The trace formula follows from
linearity of trace. The compatibility with $\alpha$ follows by evaluating
both sides at $g$.
\end{proof}

\begin{implementation}
\module{ModularRep.LinearCharacterTensorAction} defines
\lean{Representation.linearCharacterTwist} and proves
\lean{Representation.isIrreducible_linearCharacterTwist_iff},
\lean{Representation.character_linearCharacterTwist} and
\lean{Representation.linearCharacterTwist_twist}.
The induced operation on ordinary irreducible characters is
\lean{ModularRep.OrdinaryIrreducibleCharacter.linearTwist}.
\end{implementation}

The convention for automorphisms is a right action:
$(\chi^\alpha)^\beta=\chi^{\alpha\beta}$, with
$(\alpha\beta)(g)=\alpha(\beta(g))$. The combined tensor and automorphism
action below is instead written as a left action, with the indicated inverses.

More precisely, let $A$ act on $G$ through $a\mapsto\alpha_a$. Its left
action on $L_K(G)$ is $a\cdot\lambda=\lambda\circ\alpha_{a^{-1}}$.
The resulting semidirect product acts on $\Irr_K(G)$ by
\[
 (\lambda,a)\cdot\chi
   =\lambda^{-1}\,\chi^{\alpha_{a^{-1}}}.
\]
The compatibility in Proposition~\ref{lib:ordinary:linear-twist} proves
the action law. The inverses affect formulas for the action, although they
do not change the orbits or stabilisers of either factor acting alone.

\begin{implementation}
The declarations
\lean{ModularRep.OrdinaryIrreducibleCharacter.linearCharacterFieldAction},
\lean{tensorField_semidirectCompatible},
\lean{tensorFieldSemidirectAction} and
\lean{tensorFieldSemidirectAction_apply} specify this convention.
They use the semidirect product action in
\module{Formalisation.SemidirectStabilizer}. The right action by
automorphisms is represented as a left action of the opposite group.
These declarations apply to any group $A$ acting on $G$.
\end{implementation}

\begin{example}[A stabiliser calculation]
Suppose $A$ and a subgroup $L\leq L_K(G)$ are given, and $A$ preserves $L$.
Under the combined action, $(\lambda,a)$ fixes $\chi$ if and only if
$\chi(\alpha_{a^{-1}}(g))=\lambda(g)\chi(g)$ for every $g\in G$.
Pullback followed by tensoring with a linear character gives an irreducible
representation affording the transformed character. To restrict this
action to the characters in a block, one must also prove that the block
is preserved.
\end{example}

\section{Extension across a cyclic quotient}
\label{lib:ordinary:cyclic-extension}

Let $N\trianglelefteq H$ with $H$ finite. An extension of a representation
$\rho$ of $N$ is a representation $\widehat\rho$ of $H$ on the same vector
space such that
$\widehat\rho|_N\simeq\rho$. An extension makes $\rho$ invariant
under conjugation by $H$: the action of $h^{-1}$ intertwines the
conjugation twist of $\widehat\rho|_N$ with $\widehat\rho|_N$.
Also, an extension of an irreducible representation is irreducible, since
every $H$-invariant subspace is $N$-invariant.

\begin{theorem}[Cyclic extension]
\label{lib:ordinary:cyclic-principle}
Assume that $K$ is algebraically closed. Let $H$ be finite, let
$N\trianglelefteq H$, and let $\rho$ be a finite
dimensional irreducible $K$-representation of $N$. If $H/N$ is cyclic
and $\rho\circ c_h\simeq\rho$ for every $h\in H$, where
$c_h(n)=hnh^{-1}$, then $\rho$ extends to $H$.
\end{theorem}
\begin{proof}
Choose $h\in H$ whose coset generates $H/N$, and set $r=|H/N|$.
Conjugation invariance gives an invertible operator $T$ such that
$T\rho(n)T^{-1}=\rho(hnh^{-1})$ for every $n\in N$.
Consequently $\rho(h^r)^{-1}T^r$ commutes with $\rho(N)$ and is
a nonzero scalar $c$ by Schur's lemma. Choose $u\in K$ with
$u^r=c^{-1}$. Replacing $T$ by $uT$ preserves the conjugation
identity and gives $T^r=\rho(h^r)$.
These identities ensure that $\widehat\rho(nh^i)=\rho(n)T^i$,
for $n\in N$ and $0\leq i<r$, defines a representation of $H$
extending $\rho$.
\end{proof}

For the complex character formulation, see
\cite[Corollary~11.22]{Isaacs1976}. The proof above uses
algebraic closure to apply Schur's lemma and take the required root.
The modular analogue is discussed in
Section~\ref{lib:brauer:actions-extensions}.

\begin{proposition}[Extension of an invariant ordinary character]
\label{lib:ordinary:cyclic-bridge}
Assume that $K$ is algebraically closed. Let $H$ be finite, let
$N\trianglelefteq H$, let $H/N$ be cyclic, and let $\chi\in\Irr_K(N)$. If
\[
 \chi(hnh^{-1})=\chi(n)\qquad(h\in H,\ n\in N),
\]
then there exists $\widehat\chi\in\Irr_K(H)$ such that
$\widehat\chi|_N=\chi$.
\end{proposition}
\begin{proof}
Choose an irreducible realisation $\rho$ of $\chi$. Each conjugation twist
of $\rho$ has character $\chi$, so
Proposition~\ref{lib:ordinary:separation} supplies an equivalence with
$\rho$. The cyclic extension theorem gives $\widehat\rho$ and its
restriction equivalence. The extension is irreducible, and taking traces
of the intertwining identity proves the required restriction equality.
\end{proof}

\begin{implementation}
\module{ModularRep.CyclicExtension} defines
\lean{Representation.ConjugationInvariant} and
\lean{Representation.Extension}.
The Lean construction assumes
\lean{Representation.CyclicExtensionPrinciple}, which states the
representation extension property over the chosen coefficient field.
This principle is an explicit hypothesis of the Lean construction,
including when $K$ is algebraically closed. Theorem~\ref{lib:ordinary:cyclic-principle}
proves the mathematical assertion under algebraic closure.
Over an arbitrary field of characteristic zero the assertion can fail. In
\module{ModularRep.OrdinaryCharacterCyclicExtensionBridge},
\lean{ModularRep.OrdinaryIrreducibleCharacter.realisation_conjugationInvariant_of_fixed}
proves invariance of an irreducible realisation, and
\lean{ModularRep.OrdinaryIrreducibleCharacter.exists_extension_of_fixed_cyclic_quotient}
constructs the extending character under the stated extension hypothesis.
An \lean{ExtensionWitness} consists of
an irreducible character of $H$ together with its restriction identity
on $N$.
\end{implementation}

\begin{example}[Extension to a semidirect product]
Let $D$ and $A$ be finite, with $A$ cyclic. For $H=D\rtimes A$, the normal
copy of $D$ has quotient $H/D\simeq A$. A conjugation invariant irreducible
representation of $D$ therefore extends to $H$ whenever the cyclic
extension property holds over the coefficient field.
\lean{Representation.exists_extension_to_semidirect_of_cyclic_outer}
constructs this quotient isomorphism and applies the principle.
\end{example}

\chapter{Brauer characters}
\label{chap:library:brauer}

Let $\ell$ be a prime, let $G$ be finite, let $k$ be an algebraically closed
field of characteristic $\ell$, and let $K$ be a field of characteristic
zero. Fix a multiplicative correspondence between the roots of unity in
$k$ and $K$ of orders dividing $|G|_{\ell'}$. We use it to lift the
eigenvalues of a modular representation and define its Brauer character.
For the usual complex valued construction, see, for example,
\cite[pp.~17--18]{Navarro1998}. The construction here uses only a finite
group of roots of unity in $K$, which need not be algebraically closed.

We prove additivity on short exact sequences, separation of irreducible
representations and linear independence of their Brauer characters.
These results will identify composition multiplicities from character
values in Chapter~\ref{chap:library:kzero}.

\section{Regular elements and roots of unity}
\label{lib:brauer:roots}

An element is $\ell$-regular if its order is prime to $\ell$. Write
$G_{\ell'}$ for these elements and
$\operatorname{CF}_{\ell'}(G,K)$ for the $K$-valued functions on
$G_{\ell'}$ invariant under conjugation by $G$. Set
\[
 m_G=|G|_{\ell'},\qquad
 \mu_{m_G}(k)=\{a\in k^\times:a^{m_G}=1\}.
\]
Here $|G|_{\ell'}$ is the largest divisor of $|G|$ prime to $\ell$.
It is a common exponent for the regular elements and need not be their
least common exponent.

\begin{definition}
\label{lib:brauer:root-correspondence}
A root correspondence for $G$ is a multiplicative isomorphism
$\iota_G:\mu_{m_G}(k)\to\mu_{m_G}(K)$.
\end{definition}

Since $k$ is algebraically closed and $\ell\nmid m_G$, its group of
$m_G$th roots is cyclic of order $m_G$. The existence of $\iota_G$
therefore requires $K$ to contain all $m_G$th roots of unity.
The correspondence $\iota_G$ is specified separately from the modular system in
Chapter~\ref{chap:library:reduction}. When restricting characters or
comparing character tables, the chosen correspondences must agree on the
roots involved.

\begin{implementation}
\module{ModularRep.BrauerRootConvention} defines
\lean{PrimeRegularRootEmbedding.AgreesOnRoots}: two correspondences agree
on every root whose order divides the exponent for the group being restricted
to. Agreement is reflexive and is transitive when the required exponents
divide one another. The theorem \lean{agreesOnRoots_of_commonRoot}
proves agreement for correspondences induced by one common root
identification. Separate choices of roots do not establish this agreement.
\end{implementation}

\begin{lemma}[Orders of eigenvalues at regular elements]
\label{lib:brauer:eigenvalues}
Let $\rho$ be a finite dimensional $k$-representation of $G$.
If $g\in G_{\ell'}$ and $a$ is a root of the characteristic polynomial of
$\rho(g)$, then $a^{m_G}=1$.
\end{lemma}
\begin{proof}
The order of $g$ divides $|G|$ and is prime to $\ell$, so it divides $m_G$.
Hence $g^{m_G}=1$. A root of the characteristic polynomial is an eigenvalue, so
there is $v\ne0$ with $\rho(g)v=av$. Applying the $m_G$th power gives
$v=a^{m_G}v$, and cancellation proves the assertion.
\end{proof}

\begin{implementation}
\lean{ModularRep.primeRegularExponent} is $m_G$.
\lean{ModularRep.PrimeRegularRootEmbedding} specifies a prime
\lean{p} and an equivalence between the groups defined by Mathlib's
\lean{rootsOfUnity}. The exponent identity and the
eigenvalue lemma are
\lean{ModularRep.PrimeRegularElement.pow_primeRegularExponent_eq_one} and
\lean{Representation.charpolyRoot_pow_primeRegularExponent_eq_one} in
\module{ModularRep.PrimeRegularEigenvalues}.
These two elementary results only need a field and a finite dimensional
representation. The field characteristic and algebraic closure enter the
subsequent construction of Brauer characters.
\end{implementation}

\section{The character of a modular representation}
\label{lib:brauer:construction}

For a finite dimensional $k$-representation $\rho$, let
$\mathcal R_\rho(g)$ be the multiset of roots of the characteristic
polynomial of $\rho(g)$, with algebraic multiplicities.

\begin{definition}
\label{lib:brauer:character-definition}
The Brauer character afforded by $\rho$, relative to $\iota_G$, is
\[
 \varphi_\rho(g)=\sum_{a\in\mathcal R_\rho(g)}\iota_G(a)
 \qquad(g\in G_{\ell'}).
\]
An irreducible Brauer character is one afforded by an irreducible finite
dimensional $k$-representation. Their set is denoted
$\IBr_{\iota_G}(G)$.
\end{definition}

The preceding lemma places every root in the domain of $\iota_G$.
Algebraic closure ensures that the characteristic polynomial splits over
$k$. The eigenvalues are lifted to $K$ before they are summed.
Their sum in $k$ is the trace $\tr_k(\rho(g))$.

\begin{proposition}[Invariance of Brauer characters]
\label{lib:brauer:invariance}
The function $\varphi_\rho$ is constant on conjugacy classes in
$G_{\ell'}$. Equivalent representations have equal Brauer characters.
For $\alpha\in\Aut(G)$,
\[
 \varphi_{\rho\circ\alpha}(g)=\varphi_\rho(\alpha(g)).
\]
\end{proposition}
\begin{proof}
Conjugating a group element conjugates its representing operator.
An equivalence of representations also conjugates the corresponding
operators. Conjugate operators have the same characteristic polynomial,
hence the same root multiset. The definition as a sum of lifted roots gives the first two
assertions. For the last assertion the operator on the left at $g$ is
exactly the operator on the right at $\alpha(g)$.
\end{proof}

\begin{implementation}
\module{ModularRep.BrauerCharacter} defines
\lean{Representation.brauerCharacterOfRootEmbedding} using Mathlib's
characteristic polynomials and their root multisets. It has values in
\lean{ModularRep.PrimeRegularClassFunction}. It uses a total function
\lean{PrimeRegularRootEmbedding.lift}, extended by zero outside the root
group, to map the root multiset. The theorem
\lean{Representation.lift_charpolyRootAsRootOfUnity} proves agreement with
the specified root equivalence on every root that occurs.
The invariance declarations are
\lean{Representation.charpoly_conj},
\lean{Representation.brauerCharacterOfRootEmbedding_iso} and
\lean{Representation.brauerCharacterOfRootEmbedding_twist}.
The conjugacy lemma specialises Mathlib's
\lean{LinearEquiv.charpoly_conj} to representation operators.
\lean{ModularRep.IBr} consists of class functions admitting an
irreducible \lean{FDRep} realisation. Equality of its elements is
pointwise equality of functions.
\end{implementation}

\begin{example}[Diagonal representations]
Suppose $g\in G_{\ell'}$ acts diagonally with entries
$a_1,\ldots,a_d\in\mu_{m_G}(k)$. Its Brauer value is
$\iota_G(a_1)+\cdots+\iota_G(a_d)$. If several entries coincide, every
occurrence contributes to the sum. At the identity all entries are one,
so $\varphi_\rho(1)=d$ in $K$. This remains the integer $d$ even when the
trace of the identity in $k$ vanishes because $\ell\mid d$.
\end{example}

\section{Additivity on exact sequences}
\label{lib:brauer:exactness}

\begin{lemma}[Characteristic polynomial on a subspace and quotient]
\label{lib:brauer:charpoly-factorisation}
Let $T$ be an endomorphism of a finite dimensional vector space $V$ over a
field, and let $U\leq V$ be $T$-invariant. If $\bar T$ is the induced
endomorphism of $V/U$, then
\[
 \det(XI-T)=\det(XI-T|_U)\det(XI-\bar T).
\]
\end{lemma}
\begin{proof}
Choose a basis of $U$, a basis of $V/U$ and lifts of the quotient basis.
Together they give a basis of $V$ in which $T$ has the form
\[
 \begin{pmatrix}A&B\\0&D\end{pmatrix}.
\]
The diagonal blocks represent $T|_U$ and $\bar T$. The determinant of the
resulting block upper triangular polynomial matrix is the product of the
two diagonal determinants.
\end{proof}

\begin{implementation}
\lean{LinearMap.charpoly_eq_restrict_mul_mapQ} in
\module{ModularRep.CharpolyInvariantSubmodule} constructs this basis and
identifies all four blocks. It uses Mathlib's basis construction
\lean{Module.Basis.sumQuot} and block matrix identity
\lean{Matrix.charpoly_fromBlocks_zero}\textsubscript{\texttt{21}}.
\end{implementation}

\begin{theorem}[Additivity on short exact sequences]
\label{lib:brauer:short-exact}
For a short exact sequence of finite dimensional $kG$-representations
\[
 0\longrightarrow U\overset{f}{\longrightarrow}V
   \overset{\pi}{\longrightarrow}W\longrightarrow0,
\]
the Brauer characters formed using $\iota_G$ satisfy
$\varphi_V=\varphi_U+\varphi_W$. In particular,
$\varphi_{U\oplus W}=\varphi_U+\varphi_W$.
\end{theorem}
\begin{proof}
First take an invariant subspace of $V$. By
Lemma~\ref{lib:brauer:charpoly-factorisation}, the root multiset of the
operator on $V$ is the sum of the root multisets on the subspace and its
quotient. Applying the same lift to each occurrence and summing proves
additivity. This step uses no additive property of $\iota_G$.

For a short exact sequence, $\im f=\ker \pi$ is invariant. Injectivity of
$f$ identifies $U$ with $\im f$, and surjectivity of $\pi$ identifies
$V/\im f$ with $W$. Both identifications intertwine the actions. The
first part and Proposition~\ref{lib:brauer:invariance} give the assertion.
The case of a direct sum follows in the same way, or from the characteristic
polynomial of a block diagonal operator.
\end{proof}

\begin{implementation}
\module{ModularRep.BrauerCharacterExact} proves
\lean{Representation.brauerCharacterOfRootEmbedding_subrepresentation_quotient}
and \lean{Representation.brauerCharacterOfRootEmbedding_prod}.
The statement for the category of finite dimensional representations is
\lean{FDRep.brauerCharacterOfRootEmbedding_shortExact}.
\lean{FDRep.shortExact_underlying_functions} gives injectivity,
surjectivity and exactness of the underlying linear maps. Intertwining
identities identify the subrepresentation and quotient with the first
and last terms of the sequence.
\end{implementation}

\begin{example}[An invariant filtration]
If $0=V_0\leq V_1\leq\cdots\leq V_r=V$ is a finite invariant filtration,
repeated exact additivity gives
$\varphi_V=\sum_{i=1}^r\varphi_{V_i/V_{i-1}}$.
For a representation given by block upper triangular matrices, the
diagonal blocks therefore suffice to calculate its Brauer character.
The blocks above the diagonal can distinguish a nonsplit extension from
a direct sum, even though their Brauer characters agree.
\end{example}

\section{Separation and linear independence}
\label{lib:brauer:traces}

The modular traces of nonisomorphic simple representations are linearly
independent. To pass from this fact to Brauer characters, we compare the
same sums of roots of unity in the two fields. Cyclotomic polynomials
provide the comparison. For linear independence of complex valued Brauer
characters, see \cite[Theorem~2.6]{Navarro1998}.

\begin{lemma}[Traces on regular parts]
\label{lib:brauer:regular-traces}
Let $\rho$ be a finite dimensional $k$-representation of $G$.
For every $g\in G$, write $g=g_\ell g_{\ell'}$ as commuting factors of
$\ell$-power order and order prime to $\ell$. Then
\[
 \tr_k(\rho(g))=\tr_k(\rho(g_{\ell'})).
\]
\end{lemma}
This trace identity also appears in \cite[Lemma~2.4]{Navarro1998}.
\begin{proof}
If $g_\ell$ has order dividing $\ell^a$, then in characteristic $\ell$
the operator $P=\rho(g_\ell)$ satisfies
$(P-I)^{\ell^a}=P^{\ell^a}-I=0$. Let $R=\rho(g_{\ell'})$.
The operators $P$ and $R$ commute. Hence $R(P-I)$ is nilpotent and has
trace zero. Since $RP=R+R(P-I)$, taking traces proves the claim.
\end{proof}

\begin{implementation}
\lean{Representation.character_eq_character_primeRegularPart} in
\module{ModularRep.ModularTraceRegularPart} uses the commuting primary
decomposition from \module{ModularRep.PrimeRegularPart} and Mathlib's trace
identity for a commuting nilpotent term.
\end{implementation}

\begin{proposition}[Brauer characters determine irreducible representations]
\label{lib:brauer:separation}
If two finite dimensional modular representations have equal Brauer
characters relative to $\iota_G$, then their $k$-valued trace functions
on $G$ are equal. If both representations are irreducible, they are
equivalent.
\end{proposition}
\begin{proof}
Choose a primitive $m_G$th root $\zeta\in k$. Its image
$\widehat\zeta=\iota_G(\zeta)$ is primitive in $K$. At a regular element,
write the two eigenvalue multisets as powers of $\zeta$. The difference
of the sums of these powers is the value of some polynomial
$F\in\Z[X]$ at $\zeta$. Equality of the Brauer values says
$F(\widehat\zeta)=0$. The minimal polynomial of $\widehat\zeta$ over
$\Q$ is the cyclotomic polynomial $\Phi_{m_G}$, so
$\Phi_{m_G}$ divides $F$ in $\Z[X]$. Since $m_G$ is prime to
$\ell$, the primitive root $\zeta$ is also a root of the corresponding
cyclotomic polynomial in $k$. Thus $F(\zeta)=0$.

The sum of the roots of the characteristic polynomial is the trace, so the traces
agree at regular elements. Lemma~\ref{lib:brauer:regular-traces} gives
agreement on all of $G$. The argument separating simple modules by their traces in
Proposition~\ref{lib:ordinary:separation}, applied over the algebraically
closed field $k$, proves the last assertion.
\end{proof}

\begin{implementation}
\module{ModularRep.CyclotomicRootSumReflection} proves that
equality of the lifted sums in $K$ implies equality of the sums in $k$.
\module{ModularRep.BrauerTraceRecovery} identifies these sums with the
root multisets used by the character definition.
\lean{ModularRep.FDRepSimpleClassKZero.brauerCharacterDeterminesModularTrace_of_rootEmbedding}
and
\lean{ModularRep.FDRepSimpleClassKZero.irreducibleBrauerCharacterInjectivity_of_rootEmbedding}
are the resulting theorems in \module{ModularRep.BrauerCharacterSeparation}.
The latter proves that simple modules with equal Brauer characters are
isomorphic. Trace recovery is proved from the root correspondence.
\end{implementation}

\begin{theorem}[Linear independence of irreducible Brauer characters]
\label{lib:brauer:independence}
The functions in $\IBr_{\iota_G}(G)$ are linearly independent over $K$
on $G_{\ell'}$.
\end{theorem}
\begin{proof}
Choose one representative of every simple $kG$-module class. There are
finitely many such classes: each is a composition factor of the finite
dimensional regular module. Denote the representatives by
$S_1,\ldots,S_s$ and their trace functions by $t_1,\ldots,t_s$.

First, the $t_i$ are linearly independent over $k$. For each $i$, Schur's
lemma and density on $\bigoplus_jS_j$ give an element of $kG$ acting as
an operator of trace one on $S_i$ and as zero on every other summand.
Evaluating any linear relation of trace functionals at this element
extracts its $i$th coefficient. Restriction of a linear functional on
$kG$ to the basis $G$ is injective, so the trace functions on $G$ are
independent too. Lemma~\ref{lib:brauer:regular-traces} then gives
independence of their restrictions to $G_{\ell'}$.

Independent functions admit a nonsingular square evaluation matrix.
Indeed, their evaluation vectors span $k^s$, since a linear functional
annihilating that span would give a relation between the functions.
Choose a basis of evaluation vectors, at regular elements
$x_1,\ldots,x_s$, and write
\[
 M_k=(t_i(x_j))_{i,j},\qquad \det M_k\ne0.
\]
Express each root multiset occurring in this matrix as powers of the
primitive root $\zeta$. Replacing each $\zeta$ by the class of $X$
defines a matrix $M$ over the common ring
$R=\Z[X]/(\Phi_{m_G})$. Evaluation at $\zeta$ gives $M_k$.
Evaluation at $\widehat\zeta$ gives the Brauer evaluation matrix
$M_K=(\varphi_{S_i}(x_j))_{i,j}$.

The map $R\to K$ is injective by the theorem on the minimal polynomial of a primitive root of unity. If $\det M_K=0$, then $\det M=0$ in $R$, and evaluating in
$k$ would give $\det M_k=0$. This contradiction shows
$\det M_K\ne0$. Evaluating a $K$-linear relation among the Brauer
characters at $x_1,\ldots,x_s$ now forces all coefficients to vanish.
Finally, separation and realisation identify this family indexed by
simple classes with the set of functions $\IBr_{\iota_G}(G)$.
\end{proof}

\begin{implementation}
\module{ModularRep.SimpleTraceLinearIndependence} proves independence of
the modular simple trace functions, and
\module{ModularRep.PrimeRegularEvaluationMatrix} obtains the regular
evaluation points. The coefficient argument is
\lean{ModularRep.CyclotomicDeterminantSpecialization.rootMultisetMatrix_det_ne_zero_transfer},
using \lean{CyclotomicOrder} and
\lean{cyclotomicSpecialization_injective} in the same namespace.
\lean{ModularRep.FDRepSimpleClassKZero.exists_simpleClass_primeRegular_brauer_det_ne_zero}
constructs the nonsingular Brauer evaluation matrix.
The final theorem is
\lean{ModularRep.FDRepSimpleClassKZero.irreducibleBrauerCharacters_linearIndependent}.
The determinant is first computed over the cyclotomic integer ring.
Its nonvanishing after evaluation in $K$ permits the final linear
algebra argument over $K$.
\end{implementation}

\begin{example}[Equality of composition multiplicities]
Suppose exact additivity gives
$\varphi_V=\sum_i a_i\varphi_i$ and
$\varphi_W=\sum_i b_i\varphi_i$, where $a_i,b_i$ are nonnegative integers
and the $\varphi_i$ are distinct irreducible Brauer characters.
If $\varphi_V=\varphi_W$, linear independence gives
$a_i=b_i$ for every $i$, using characteristic zero to recover equality
of integers from equality in $K$. This is the step used to prove
injectivity of the Brauer character map on the exact Grothendieck group.
It determines composition multiplicities without asserting that
$V\simeq W$.
\end{example}

\section{Compatible choices under pullback and restriction}
\label{lib:brauer:compatible-roots}

Let $f:H\to G$ be a homomorphism of finite groups with root
correspondences $\iota_H$ and $\iota_G$. A regular element of $H$ maps
to a regular element of $G$, so class functions can always be pulled back
by $f$. To compare this pullback with the character of $f^*\rho$, the
two lifts must agree on the roots of $\rho(f(h))$ for every
$h\in H_{\ell'}$.

\begin{proposition}[Pullback with compatible root correspondences]
\label{lib:brauer:compatible-pullback}
Let $f:H\to G$ be a homomorphism and let $\rho$ be a finite dimensional
$k$-representation of $G$. Suppose that $\iota_H(a)=\iota_G(a)$
for every eigenvalue $a$ of $\rho(f(h))$, for every $h\in H_{\ell'}$.
Then
\[
 \varphi_{f^*\rho}^{\,\iota_H}(h)
     =\varphi_\rho^{\,\iota_G}(f(h))
       \qquad(h\in H_{\ell'}).
\]
Such agreement holds if the two correspondences are restrictions of one
multiplicative equivalence $\mu_M(k)\simeq\mu_M(K)$, where both
$m_H$ and $m_G$ divide $M$.
\end{proposition}
\begin{proof}
The representing operator on both sides is $\rho(f(h))$. Its root
multiset is therefore the same on both sides, and the agreement of
the lifts gives equality term by term. A common root correspondence
gives that agreement on every root in the intersection of the two
root groups.
\end{proof}

\begin{implementation}
Compatibility on these eigenvalues is expressed by
\lean{Representation.BrauerRootLiftCompatibleAlong}.
\lean{Representation.brauerCharacterOfRootEmbedding_pullback_of_compatible}
proves the formula.
\module{ModularRep.BrauerCharacterCommonRootCompatibility} restricts one
common equivalence and proves the condition for any homomorphism.
Equality of the total functions $k\to K$ is a sufficient stronger
condition, but is unnecessary away from the occurring roots.
\end{implementation}

A group isomorphism $\phi:G\to H$ preserves $|G|_{\ell'}$ and therefore
transports a chosen root correspondence with the same function lifting
roots. Pullback along $\phi^{-1}$ then transports an irreducible Brauer
character together with an affording representation. The constructions
\lean{PrimeRegularRootEmbedding.alongMulEquiv} and
\lean{IrreducibleBrauerCharacter.alongMulEquiv} in
\module{ModularRep.BrauerCharacterEquivTransport} give this transport.
If an extension is also transported, the ambient and base group
isomorphisms must commute with the inclusion maps. A separately specified
target character must be identified with the transported base character.
These are the hypotheses of
\lean{Representation.Extension.BrauerCharacterExtensionWitness.alongMulEquivOfBase}
in \module{ModularRep.BrauerCharacterExtensionWitnessTransport}.

\begin{proposition}[Extending a root correspondence from a subgroup]
\label{lib:brauer:root-reselection}
Let $H\leq G$, and prescribe a root correspondence for $H$. If at least
one root correspondence for $G$ exists, one can choose a correspondence
for $G$ whose restriction agrees with the prescribed correspondence
on $\mu_{m_H}(k)$.
\end{proposition}
\begin{proof}
The divisibility $m_H\mid m_G$ permits restriction of any chosen ambient
correspondence to the smaller root group. If $j$ is this restriction, then
$j^{-1}\iota_H$ is an automorphism of the finite cyclic group
$\mu_{m_H}(k)$. An automorphism of a finite cyclic subgroup extends
to the ambient finite cyclic group: in power coordinates this is
surjectivity of reduction on the corresponding groups of units.
Compose the ambient correspondence with such an extension. Its
restriction is now the prescribed one.
\end{proof}

\begin{implementation}
\lean{ModularRep.PrimeRegularRootEmbedding.exists_ambient_agreeing_on_subgroup}
assumes existence of an ambient root correspondence, expressed by
\lean{Nonempty}, and uses
\lean{exists_mulAut_extension} in the same namespace. It adjusts an
ambient correspondence while retaining the prescribed subgroup correspondence.
\end{implementation}

There is a particularly direct choice for a normal $\ell$-subgroup
$N\trianglelefteq G$. Since $|G|=|G/N|\,|N|$, one has
$m_G=m_{G/N}$. A prescribed quotient correspondence can therefore be
transported to $G$ with the same root lifting map.
\lean{ModularRep.PrimeRegularRootEmbeddingPQuotient.exponent_quotient_eq},
\lean{ofPQuotient} and \lean{ofPQuotient_lift} implement this observation.

\section{Inflation and descent of irreducible Brauer characters}
\label{lib:brauer:descent}

For $\varphi\in\IBr_{\iota_G}(G)$, define $\ker\varphi=\ker\rho$,
where $\rho$ is any irreducible representation affording $\varphi$.
Proposition~\ref{lib:brauer:separation} makes this independent of the
choice of $\rho$, since equivalent representations have the same kernel.
Let $\pi:G\twoheadrightarrow H$ be a surjection.

\begin{theorem}[Inflation of irreducible Brauer characters]
\label{lib:brauer:quotient-equivalence}
Suppose that $\ell\nmid|\ker \pi|$, and that the root correspondences
are compatible along $\pi$ for every finite dimensional representation
of $H$. Inflation gives a bijection between
$\IBr_{\iota_H}(H)$ and
$\{\varphi\in\IBr_{\iota_G}(G):\ker \pi\leq\ker\varphi\}$.
If $\bar\varphi$ corresponds to $\varphi$, then
\[
 \varphi(g)=\bar\varphi(\pi(g))\qquad(g\in G_{\ell'}).
\]
\end{theorem}
\begin{proof}
Pullback along $\pi$ preserves irreducibility, and the pulled back
representation has $\ker \pi$ in its kernel. Root compatibility identifies
its Brauer character with the inflated function. Conversely, descend a
representation affording $\varphi$ on the same vector space, using
$\ker \pi\leq\ker\varphi$.
Representation descent preserves irreducibility, and compatible
pullback gives the asserted identity.

It remains to prove that pullback is injective on functions on $H_{\ell'}$. Every lift
$g$ of an $\ell$-regular $h$ is itself $\ell$-regular in this case.
Indeed, if $h$ has order $d$, then $g^d\in\ker \pi$, so the order of $g$ divides
$d\,|\ker \pi|$, which is prime to $\ell$. Hence $\pi$ is surjective on
regular elements and pullback of their class functions is injective.
The two character identities now show that inflation and descent are inverse.
\end{proof}

\begin{implementation}
\lean{ModularRep.IrreducibleBrauerCharacterSurjectiveDescent.KernelTrivialIBrAlong}
consists of irreducible Brauer characters admitting a realisation on
which $\ker \pi$ acts trivially. By character separation, this is the
condition $\ker \pi\leq\ker\varphi$. In the same namespace,
\lean{inflateToKernelTrivialIBrAlong}, \lean{descendIBrAlong} and
\lean{quotientIBrEquivKernelTrivialIBrAlong} construct the maps and the
bijection. The representation descent itself only requires surjectivity
and an irreducible realisation on which $\ker \pi$ acts trivially. Root
compatibility proves its character identity. The hypothesis that $\ker \pi$
has order prime to $\ell$ proves injectivity of pullback on class functions.
This hypothesis differs from the normal $\ell$-subgroup hypothesis in
the preceding construction of a root correspondence.
The classical inflation correspondence allows an arbitrary normal kernel
\cite[p.~39]{Navarro1998}. The theorem stated here uses the stronger
hypothesis that the kernel has order prime to $\ell$, as does the cited
Lean equivalence.
\end{implementation}

\section{Tensor actions and cyclic extensions}
\label{lib:brauer:actions-extensions}

For a modular linear character $\lambda:G\to k^\times$, define
$\widehat\lambda(g)=\iota_G(\lambda(g))$ on $G_{\ell'}$.
The value $\lambda(g)$ belongs to the required root group because
$g^{m_G}=1$. Multiplicativity of the root correspondence proves
$\widehat{\lambda\mu}=\widehat\lambda\widehat\mu$ and
$(\widehat\lambda)^\alpha=\widehat{\lambda\circ\alpha}$.

\begin{theorem}[Tensoring with a linear character]
\label{lib:brauer:tensor-source}
For every finite dimensional modular representation $V$ and every
$\lambda:G\to k^\times$, one has
\[
 \varphi_{V\otimes\lambda}=\widehat\lambda\,\varphi_V.
\]
\end{theorem}
\begin{proof}
At $g\in G_{\ell'}$, the representing operator on $V\otimes\lambda$
is $\lambda(g)\rho_V(g)$. Its eigenvalues, counted with multiplicity,
are $\lambda(g)a$ as $a$ runs through the eigenvalues of $\rho_V(g)$.
Multiplicativity of $\iota_G$ gives
$\iota_G(\lambda(g)a)=\widehat\lambda(g)\iota_G(a)$.
Summing proves the formula.
\end{proof}

This is the case of a one dimensional factor in the tensor product
formula for Brauer characters. See, for example,
\cite[proof of Theorem~2.23]{Navarro1998}.
The argument with invariant subspaces in
Proposition~\ref{lib:ordinary:linear-twist} also applies over $k$.
Thus $V\otimes\lambda$ is irreducible when $V$ is irreducible, and
$\widehat\lambda\varphi_V$ is then an irreducible Brauer character.

\begin{implementation}
The Lean construction uses
\lean{ModularRep.BrauerLinearTensorProductFormula} as a hypothesis.
\module{ModularRep.BrauerLinearCharacterTensorAction} defines
\lean{ModularRep.PrimeRegularRootEmbedding.liftedLinearCharacter} and proves
its multiplicative and automorphism identities. Given the tensor formula,
\lean{ModularRep.IrreducibleBrauerCharacter.linearTwist} constructs the
irreducible Brauer character. Its \lean{linearTwist_mul},
\lean{twist_linearTwist} and \lean{tensorFieldSemidirectAction} give the
action laws, with the inverse convention explained in
Section~\ref{lib:ordinary:actions}.
\end{implementation}

\begin{theorem}[Cyclic extension of a Brauer character]
Let $H$ be finite, let $N\trianglelefteq H$, and suppose that $H/N$
is cyclic. Choose root correspondences $\iota_N$ and $\iota_H$
compatible with restriction to $N$ for every finite dimensional
$k$-representation of $H$. If
$\varphi\in\IBr_{\iota_N}(N)$ is invariant under conjugation by $H$,
then there is $\widehat\varphi\in\IBr_{\iota_H}(H)$ with
$\widehat\varphi|_{N_{\ell'}}=\varphi$.
\end{theorem}
\begin{proof}
Choose an irreducible realisation $\rho$ of $\varphi$. Its conjugation
twists have the same Brauer character, so
Proposition~\ref{lib:brauer:separation} makes them equivalent to $\rho$.
The proof of Theorem~\ref{lib:ordinary:cyclic-principle} applies over
the algebraically closed field $k$ in any characteristic. It therefore
gives an extension $\widehat\rho$, which is irreducible because
$\rho$ is irreducible. Compatible roots identify the restriction of
its Brauer character with $\varphi$ by
Proposition~\ref{lib:brauer:compatible-pullback}.
\end{proof}

For the usual complex valued formulation and the representation extension
argument, see \cite[Theorem~8.12]{Navarro1998}.
There is no restriction on the order of the cyclic quotient relative to
$\ell$. The value field $K$ need not be algebraically closed, but the
root correspondences must satisfy the stated compatibility.

\begin{implementation}
The Lean construction assumes
\lean{Representation.BrauerCyclicExtensionPrinciple} over the
algebraically closed modular field $k$. It does not prove this
extension property from algebraic closure.
\module{ModularRep.BrauerCharacterExtensionBridge} proves
\lean{Representation.conjugationInvariant_of_brauerCharacter_fixed} and
\lean{Representation.exists_extension_realisation_of_ibr_fixed_cyclic_quotient}.
The latter gives an affording representation and a representation
extension. In \module{ModularRep.BrauerCharacterHomPullback},
\lean{Representation.Extension.brauerCharacterExtensionWitnessOfCompatible}
uses compatible roots to construct the extending irreducible Brauer character
and its restriction equality.
\end{implementation}

\chapter{Grothendieck groups and composition factors}
\label{chap:library:kzero}

The Grothendieck group assigns a class $[V]$ to each representation, with
$[V]=[U]+[W]$ for every short exact sequence $0\to U\to V\to W\to0$.
For finite dimensional representations, the Jordan--H\"older theorem shows
that the classes of simple modules form a basis. Thus composition
multiplicities determine the class of a representation.

In the parts of this chapter concerning representations, $G$ is a finite
group and $k$ is a field. No splitting or characteristic assumption is needed
until the final section on Brauer characters. Write $kG$ for the group algebra,
$\operatorname{FDRep}_k(G)$ for the category of finite dimensional
$k$-representations, and $K_0(kG)$ for its Grothendieck group defined by short exact sequences.

\section{The Grothendieck group and its universal property}
\label{lib:kzero:presentation}

\begin{definition}
Let $\mathcal C$ be an abelian category whose isomorphism classes of
objects form a set.
Let $F(\mathcal C)$ be the free abelian group on the isomorphism classes
of its objects.
The Grothendieck group defined by short exact sequences is
\[
 K_0(\mathcal C)=F(\mathcal C)/R(\mathcal C),
\]
where
\[
 R(\mathcal C)=
 \left\langle e_V-e_U-e_W:
 0\longrightarrow U\longrightarrow V\longrightarrow W\longrightarrow0
 \text{ is exact}\right\rangle.
\]
The symbol $e_X$ denotes the generator corresponding to the isomorphism
class of $X$. Its image in the quotient is denoted by $[X]$.
\end{definition}

For categories of modules, this definition is recalled in
\cite[Appendix, p.~163]{Serre1977}.
Isomorphic objects have the same class. Each defining relation gives
$[V]=[U]+[W]$. The short exact
sequence with all three terms zero implies
$[0]=[0]+[0]$, hence
$[0]=0$. The split sequence for $U\oplus W$ gives
$[U\oplus W]=[U]+[W]$. Exact relations also apply to extensions that do not
split, which is the feature needed for modular representations.

\begin{proposition}[Universal property]
\label{lib:kzero:universal}
Let $A$ be an abelian group. An isomorphism invariant assignment $X\mapsto a_X$
defines a homomorphism $K_0(\mathcal C)\to A$ with $[X]\mapsto a_X$ if and
only if
\[
 a_V=a_U+a_W
\]
for every short exact sequence $0\to U\to V\to W\to0$. Such a homomorphism is
unique. Consequently two homomorphisms from $K_0(\mathcal C)$ agree if they
agree on the classes $[X]$.
\end{proposition}
\begin{proof}
The assignment defines a unique homomorphism $F(\mathcal C)\to A$ by the
universal property of the free abelian group. It vanishes on
$e_V-e_U-e_W$ precisely when the stated identity holds. Vanishing on the
generators of $R(\mathcal C)$ is equivalent to vanishing on that subgroup,
so the homomorphism descends to the quotient. Conversely every homomorphism
on the quotient respects its defining relations. Object classes generate
the quotient, which proves uniqueness.
\end{proof}

\begin{implementation}
\module{ModularRep.ExactGrothendieckGroup} constructs this presentation
from Mathlib's skeleton, free abelian group and short exact complexes. Here
\lean{FreeClassGroup} denotes the free abelian group,
\lean{shortExactRelationSubgroup} its subgroup of relations, and
\lean{KZero} the quotient. The identities for isomorphic objects and
short exact sequences are \lean{classOf_iso} and
\lean{classOf_middle_eq}. The universal property is \lean{lift}, with
\lean{lift_classOf} giving its value on an object class and
\lean{hom_ext} expressing uniqueness. The elements of {\small\texttt{Skeleton C}}
are isomorphism classes of objects.
These definitions also apply to categories with zero morphisms, using
short exact complexes $U\to V\to W$. The statements above concern
abelian categories.
\end{implementation}

\begin{proposition}[Maps induced by exact functors]
\label{lib:kzero:functor}
Let $F:\mathcal C\to\mathcal D$ be an exact functor between abelian
categories, that is, an additive functor that preserves short exact
sequences. Then $F$ induces a homomorphism
\[
 K_0(F):K_0(\mathcal C)\longrightarrow K_0(\mathcal D),
 \qquad [X]\longmapsto[F(X)].
\]
\end{proposition}
\begin{proof}
The hypotheses ensure that the image of a short exact sequence is short
exact. Thus the free homomorphism induced by $F$ sends every defining
relation to a defining relation for $K_0(\mathcal D)$. It therefore induces a
homomorphism on the quotients.
\end{proof}

\begin{implementation}
The induced homomorphism is \lean{ExactGrothendieckGroup.map}.
Its hypotheses require preservation of zero morphisms, finite limits and
finite colimits. For abelian categories, these are equivalent to exactness.
The construction uses
\lean{freeMap_shortExactRelation} and
\lean{map_shortExactRelationSubgroup}. The formula on object classes is
\lean{map_classOf}. No choice of bases is involved.
\end{implementation}

\section{Composition factors and the Jordan--H\"older theorem}
\label{lib:kzero:finite-length}

Let $R$ be a ring and $M$ a left $R$-module of finite length. A composition series of $M$ is a chain
\[
 0=M_0<M_1<\cdots<M_r=M
\]
whose successive quotients $M_{i+1}/M_i$ are simple. We also consider
composition series between two fixed submodules $A\leq B$ of $M$, with
first term $A$ and last term $B$. For an account of composition series and
the Jordan--H\"older theorem, see, for example,
\cite[\S1.1, especially Theorem~1.1.4]{Benson1998}.

\begin{definition}
Write $\operatorname{Mod}R$ for the category of all left $R$-modules in
the chosen universe. Recall that $F(\operatorname{Mod}R)$ is the free
abelian group on their isomorphism classes.
For a composition series $s$ of $M$, its formal factor sum is
\[
 J(s)=\sum_{i=0}^{r-1} e_{M_{i+1}/M_i}
 \in F(\operatorname{Mod}R),
\]
and the coefficient $m_X(s)$ at an isomorphism class $X$ counts the factors
isomorphic to $X$.
\end{definition}

\begin{proposition}[Independence of composition multiplicities]
\label{lib:kzero:series-independent}
Any two composition series $s$ and $t$ between the same two submodules
of $M$ have the same simple factors up to isomorphism, counted with
multiplicity.
Thus $J(s)=J(t)$ and $m_X(s)=m_X(t)$ for every module
isomorphism class $X$. In particular, a composition series from $0$ to a
module of finite length defines $J(M)$ and $m_X(M)$ independently of
the choice of series.
\end{proposition}
\begin{proof}
The Jordan--H\"older theorem gives a bijection between the factor indices of
the two series and an isomorphism between corresponding factors. After
passing to isomorphism classes, reindexing the finite sum gives the result.
This proves equality in the free abelian group, before imposing exact relations.
\end{proof}

\begin{implementation}
\module{ModularRep.JordanHolderCoordinates} defines
\lean{compositionFactor}, \lean{compositionFactorIsoSum}, and
\lean{compositionCoefficient}. Each quotient is formed inside its
numerator: the denominator in
\lean{compositionFactor} is the inverse image of the preceding submodule
under the numerator's inclusion into $M$. Simplicity follows from the
covering relation between adjacent submodules in
\lean{compositionFactor_isSimple}.
The proof of \lean{compositionFactorIsoSum_eq} applies Mathlib's
\lean{CompositionSeries.jordan_holder} and the corresponding module
isomorphisms. The formula
\lean{compositionCoefficient_eq_sum_ite} counts the factors in each
isomorphism class.
\end{implementation}

Every finite dimensional representation has finite length as a module over
its monoid algebra, even when the monoid is infinite. Indeed it is both
Noetherian and Artinian as a vector space. A chain of submodules over the
monoid algebra is also a chain of vector subspaces, so both chain conditions
pass to that module structure. In
\module{ModularRep.FDRepFiniteLength},
\lean{asModule_isFiniteLength} applies Mathlib's criteria for finite length
and its results on restriction of scalars. The theorem
\lean{exists_compositionSeries} then applies Mathlib's equivalence between
finite length and the existence of a composition series from $0$ to the whole module.
For finite $G$, all simple $kG$-modules are finite dimensional, as shown
in Section~\ref{lib:kzero:simple-basis}.

\section{Splicing series in a short exact sequence}
\label{lib:kzero:splicing}

\begin{proposition}[Additivity of composition factors]
\label{lib:kzero:factor-additivity}
For a short exact sequence of finite length left $R$-modules
\[
 0\longrightarrow U\xrightarrow{f}V\xrightarrow{\pi}W\longrightarrow0,
\]
the formal factor sums satisfy $J(V)=J(U)+J(W)$. Consequently
$m_X(V)=m_X(U)+m_X(W)$ for every module isomorphism class $X$.
\end{proposition}
\begin{proof}
Choose composition series
$0=U_0<\cdots<U_a=U$ and $0=W_0<\cdots<W_b=W$.
The images $f(U_i)$ form a composition series from $0$ to $\im f$, since
$f$ is injective. The inverse images $\pi^{-1}(W_j)$ form one from
$\ker \pi$ to $V$, since $\pi$ is surjective. Exactness identifies the meeting
terms $\im f=\ker \pi$, so the two chains concatenate to
\[
 0=f(U_0)<\cdots<f(U_a)=\pi^{-1}(W_0)
 <\cdots<\pi^{-1}(W_b)=V.
\]
The factor isomorphisms are
\[
 U_{i+1}/U_i\simeq f(U_{i+1})/f(U_i),
 \qquad
 \pi^{-1}(W_{j+1})/\pi^{-1}(W_j)\simeq W_{j+1}/W_j.
\]
The first is induced by the injective map $f$. For the second,
restrict $\pi$ to $\pi^{-1}(W_{j+1})$. The induced map on quotients is
surjective and has zero kernel, since the inverse image of $W_j$ is
$\pi^{-1}(W_j)$. Thus the concatenated series has
factor sum $J(U)+J(W)$. Jordan--H\"older compares it with any chosen full
series of $V$.
\end{proof}

\begin{implementation}
The proof in \module{ModularRep.FDRepJordanHolderKZero} constructs these
chains as \lean{mapCompositionSeries} and \lean{comapCompositionSeries}.
The factor isomorphisms are
\lean{quotientMapEquivOfInjective} and
\lean{quotientComapEquivOfSurjective}. Concatenation uses Mathlib's
composition series operation \lean{smash}, and the local theorem
\lean{compositionFactorIsoSum_smash} expresses its factor sum as the sum
of the two factor sums. This gives
\lean{compositionFactorIsoSum_eq_add_of_exact}.

For representations, \lean{fdRepJordanHolderSum_shortExact} applies the
exact underlying $kG$-module functor and uses the injection, surjection,
and exactness of the resulting linear maps. The universal property gives
\lean{fdRepJordanHolderKZeroHom}, with values in the free abelian group
on module isomorphism classes.
\end{implementation}

\section{The basis of simple modules}
\label{lib:kzero:simple-basis}

Let $\mathcal S_k(G)$ denote the set of isomorphism classes of simple left
$kG$-modules. Every such module is cyclic: any nonzero vector generates a
nonzero submodule and hence the whole module. Since $G$ is finite, $kG$ is
finite dimensional over $k$, so every simple module belongs to
$\operatorname{FDRep}_k(G)$.

For an account of the following standard description of the Grothendieck
group, see \cite[\S14.1, Proposition~40]{Serre1977}.

\begin{theorem}[Basis of simple module classes]
\label{lib:kzero:basis-theorem}
The map
\[
 K_0(kG)\longrightarrow\bigoplus_{S\in\mathcal S_k(G)}\Z e_S,
 \qquad [V]\longmapsto\sum_S m_S(V)e_S
\]
is an isomorphism of abelian groups. Its inverse sends $e_S$ to $[S]$.
\end{theorem}
\begin{proof}
All factors in $J(V)$ are simple, so its sum lies in the subgroup freely
generated by $\mathcal S_k(G)$. The additivity just proved gives the
homomorphism in the statement. In the other direction, the free abelian
group admits a homomorphism $e_S\mapsto[S]$.

For a composition series from $0$ to $V$, each step has a short exact sequence
\[
 0\longrightarrow V_i\longrightarrow V_{i+1}
 \longrightarrow V_{i+1}/V_i\longrightarrow0.
\]
Starting from $[0]=0$ and successively applying the defining relations gives
\[
 [V]=\sum_{i=0}^{r-1}[V_{i+1}/V_i].
\]
Thus the composite back to $K_0(kG)$ fixes every object class and hence the
whole group. For a simple module $S$, the series $0<S$ has the single
factor $S$. The other composite therefore
fixes every generator $e_S$. The maps are inverse.
\end{proof}

\begin{implementation}
In \module{ModularRep.FDRepSimpleClassKZero},
\lean{SimpleModuleClass} is the subtype of the module skeleton represented
by simple modules. The homomorphism \lean{simpleClassProjection} fixes the
generators represented by simple modules and sends the others to zero.
Composing it with the homomorphism defined by composition multiplicities gives
\lean{fdRepSimpleJordanHolderHom}.

The reverse map is \lean{simpleClassToFDRepKZero}.
The proof first sums the step relations in the category
{\small\texttt{FGModuleCat k[G]}} of finitely generated $kG$-modules using
\lean{compositionSeriesFGKZero_decomposition}, and then applies
\lean{FDRepGroupAlgebraEquivalence.equivalenceFGModule}.
This is a restriction of Mathlib's equivalence between representations
and modules over the monoid algebra. Finiteness of $G$ ensures that finite
generation over $kG$ implies finite dimension over $k$. The resulting
isomorphism is
\lean{fdRepKZeroEquivSimpleModuleClassGroup}.
\end{implementation}

\begin{example}[Classes of extensions of simple modules]
\label{lib:kzero:extension-example}
Suppose $0\to S\to V\to T\to0$ is short exact, where $S$ and $T$ are
simple. Then $[V]=[S]+[T]$. If $S\not\simeq T$, the coefficients of
$[S]$ and $[T]$ are one and all other coefficients are zero. If $S\simeq T$,
the coefficient of $[S]$ is two and all other coefficients are zero.
Thus any two such extensions have the same class in $K_0(kG)$, even when
their middle terms are not isomorphic.
\end{example}

For virtual classes, coefficients may be negative. Equality in $K_0(kG)$
means equality of all these integer coefficients. For a representation
$V$, the coefficient of $[S]$ is $m_S(V)$, its composition multiplicity.
This is the identity
\lean{FDRepJordanHolderKZero.fdRepJordanHolderCoordinate_classOf}.

\section{Finiteness of the set of simple classes}
\label{lib:kzero:simple-finite}

\begin{proposition}[Finiteness of the simple module classes]
\label{lib:kzero:finite-simples-theorem}
If the left regular module of a ring $R$ has a full finite composition
series, every simple left $R$-module occurs among its factors. In particular
$\mathcal S_k(G)$ is finite for a finite group $G$ and any field $k$.
\end{proposition}
\begin{proof}
A simple module $S$ is isomorphic to $R/I$ for a maximal left ideal $I$:
the map $R\to S$ sending $r$ to $rs$ is onto for any nonzero $s\in S$.
Since $R$ has finite length, so does $I$. Append $R$ to a
composition series of $I$. The resulting series of $R$ has final
factor $R/I\simeq S$. The Jordan--H\"older theorem identifies this
factor with a factor of any prescribed composition series from $0$ to $R$.

The finite index set of the factors therefore maps onto all simple module
classes. For $R=kG$, finite dimension over $k$ makes the regular module
both Noetherian and Artinian, so it has a full finite composition series.
\end{proof}

\begin{implementation}
\module{ModularRep.SimpleModuleClassFinite} proves
\lean{simpleClass_appears_in_compositionSeries} using a series whose
last factor is the quotient by a maximal left ideal.
\lean{finite_simpleModuleClass_of_compositionSeries} applies finiteness
to the resulting surjection. The result
\lean{finiteSimpleModuleClassGroupAlgebra} uses the chain conditions
that follow from finite dimension of $kG$.
\end{implementation}

Proposition~\ref{lib:kzero:finite-simples-theorem}, together with
Theorem~\ref{lib:kzero:basis-theorem}, proves that $K_0(kG)$ is free
of finite rank over $\Z$, for every field $k$.

\section{The Grothendieck group of all modules}
\label{lib:kzero:swindle}

\begin{proposition}[Eilenberg swindle]
\label{lib:kzero:swindle-theorem}
For a ring $R$, the exact Grothendieck group of the category of all left
$R$-modules in the chosen universe is zero.
\end{proposition}
\begin{proof}
For any $R$-module $M$, let $T=\prod_{n\geq0}M$. The shift isomorphism
\[
 M\oplus T\longrightarrow T,
 \qquad (x,f)\longmapsto(x,f(0),f(1),\ldots)
\]
is $R$-linear. The split short exact sequence
$0\to M\to M\oplus T\to T\to0$ gives $[M\oplus T]=[M]+[T]$.
The isomorphism gives $[M\oplus T]=[T]$, so cancellation yields $[M]=0$.
Since object classes generate the group, every element is zero.
\end{proof}

\begin{implementation}
\module{ModularRep.ModuleCategoryKZeroTrivial} uses the countable product
{\small\texttt{Nat -> M}}. Its results are \lean{classOf_eq_zero} and
\lean{kZero_subsingleton}. The homomorphism of composition factors
\lean{fdRepJordanHolderKZeroHom} takes values in
{\small\texttt{FreeClassGroup (ModuleCat k[G])}}, the free abelian group on
isomorphism classes of all $kG$-modules. Imposing all short exact relations would give
{\small\texttt{KZero (ModuleCat k[G])}}, which is zero by this argument.
\end{implementation}

For a nonzero $kG$-module $M$, the product $\prod_{n\geq0}M$ is not finite
dimensional over $k$, so the argument does not apply in
$\operatorname{FDRep}_k(G)$.
This is why the restriction to finite dimensional representations is
necessary in Theorem~\ref{lib:kzero:basis-theorem}.

\section{The Brauer character homomorphism}
\label{lib:kzero:characters}

Now let $\ell$ be prime, let $k$ be algebraically closed of characteristic
$\ell$, and let $K$ be a field of characteristic zero. Fix a finite root
correspondence $\iota$ for $G$ as in Chapter~\ref{chap:library:brauer}.
Let $\operatorname{CF}_{\ell'}(G,K)$ be the additive group of $K$-valued
class functions on the $\ell$-regular elements of $G$.

\begin{proposition}[Injectivity of the Brauer character homomorphism]
\label{lib:kzero:brauer-injective}
Taking Brauer characters induces an injective homomorphism
\[
 \operatorname{br}_{\iota}:K_0(kG)\longrightarrow
 \operatorname{CF}_{\ell'}(G,K),
 \qquad [V]\longmapsto\varphi_V.
\]
\end{proposition}
\begin{proof}
Isomorphism invariance and short exact additivity of Brauer characters
allow the universal property to define the homomorphism. Under the basis
of Theorem~\ref{lib:kzero:basis-theorem}, it sends $e_S$ to the irreducible
Brauer character $\varphi_S$. The functions $\varphi_S$ are pairwise distinct by
Proposition~\ref{lib:brauer:separation}, and linearly independent over
$K$ by Theorem~\ref{lib:brauer:independence}. Linear independence forces
each coefficient in an integer relation to vanish in $K$. Since
$\operatorname{char}K=0$, the map $\Z\to K$ is injective, so the
coefficients vanish in $\Z$.
\end{proof}

\begin{implementation}
\module{ModularRep.BrauerCharacterKZero} constructs
\lean{FDRepSimpleClassKZero.brauerCharacterKZeroHom} using exact
additivity. The theorem \lean{brauerCharacterKZeroHom_injective}
uses the basis of simple module classes and
\lean{irreducibleBrauerCharacterLinearIndependence_of_rootEmbedding}.
The proof restricts scalars from $K$ to $\Z$ and uses the injectivity of
the corresponding map taking a finite integer linear combination of characters.
\end{implementation}

Thus representations with equal Brauer characters have the same
composition multiplicities. In Chapter~\ref{chap:library:reduction},
character compatibility then identifies the Grothendieck classes of
different lattice reductions and proves independence of the chosen lattice.

\chapter{Stable lattices and reduction}
\label{chap:library:reduction}

A stable lattice in an ordinary representation gives a modular
representation by reduction modulo the maximal ideal. Different lattices
may have nonisomorphic reductions, but the multiplicities of their simple
composition factors agree. The decomposition homomorphism assigns this
common Grothendieck class to the ordinary representation.

The construction below identifies tensor reduction with reduction modulo
the maximal ideal. Lattice independence and the decomposition homomorphism
then follow from the ordinary and Brauer character identity in
Hypothesis~\ref{lib:reduction:character-hypothesis}.

Throughout the chapter, $G$ is a finite group and $\ell$ is prime.
The coefficient fields for ordinary and modular representations are
$K$ and $k$, respectively. The integral coefficient ring is denoted by
$\mathcal O$. Each section states the coefficient hypotheses it uses.

\section{Integral coefficients and the residue map}
\label{lib:reduction:modular-system}

\begin{definition}
A complete $\ell$-modular system $(K,\mathcal O,k)$ consists of a complete
discrete valuation ring $\mathcal O$, its fraction field $K$ of
characteristic zero, a field $k$ of characteristic $\ell$, and a specified
isomorphism
\[
 \mathcal O/\mathfrak m\simeq k,
\]
where $\mathfrak m$ is the maximal ideal of $\mathcal O$. Completeness is
with respect to the $\mathfrak m$-adic topology. The residue map
$\pi:\mathcal O\to k$ is the composite of the quotient map with this
isomorphism.
\end{definition}

The map $\pi$ is onto and has kernel $\mathfrak m$. Its choice specifies
the $\mathcal O$-algebra structure on $k$ used in tensor reduction.
Splitting hypotheses for $G$ and a correspondence between roots of unity
in $k$ and $K$ are additional conditions. Algebraic closure of the
fraction field $K$ is not required.

\begin{implementation}
The definition in \module{ModularRep.ModularSystem} includes the prime,
the complete discrete valuation ring, its fraction field, the two
characteristics, and the residue field isomorphism \lean{residueEquiv}.
The definitions \lean{ModularSystem.residue} and
\lean{ModularSystem.residueAlgebra} give $\pi$ and the induced
$\mathcal O$-algebra structure on $k$.
The theorems \lean{residue_surjective} and \lean{ker_residue} prove the
two properties of $\pi$ used below. All reductions associated with this
modular system use \lean{residueAlgebra}.
\end{implementation}

\section{Constructing a stable lattice}
\label{lib:reduction:stable-lattice}

\begin{samepage}
For this section it is enough that $\mathcal O$ is a commutative ring,
$K$ is a field with an $\mathcal O$-algebra structure, and $V$ is a finite
dimensional $K$-space, with its compatible $\mathcal O$-module structure.
Let $\rho:G\to\GL_K(V)$ be a representation.

\end{samepage}

\begin{definition}
An $\mathcal O$-lattice in $V$ is a finitely generated
$\mathcal O$-submodule $L\leq V$ whose $K$-linear span is $V$.
It is stable if $\rho(g)L\subseteq L$ for every $g\in G$.
\end{definition}

When $\mathcal O$ is a discrete valuation ring with fraction field $K$,
this agrees with the usual notion of a full lattice. Only finite
generation and full span are needed for the construction below. For the
usual construction over a discrete valuation ring, see, for example,
\cite[\S15.2]{Serre1977}.

\begin{proposition}[Stable closure of a lattice]
\label{lib:reduction:existence}
If $L_0$ is an $\mathcal O$-lattice in $V$, then
\[
 L=\sum_{g\in G}\rho(g)L_0
\]
is a stable $\mathcal O$-lattice containing $L_0$. In particular every
finite dimensional $K$-representation of a finite group has a stable
$\mathcal O$-lattice.
\end{proposition}
\begin{proof}
Each image $\rho(g)L_0$ is finitely generated over $\mathcal O$.
There are only finitely many such submodules, so their sum is finitely
generated. The term $g=1$ shows $L_0\subseteq L$, which implies that the
$K$-span of $L$ is all of $V$. For $h\in G$,
$\rho(h)\rho(g)L_0=\rho(hg)L_0$ is another summand. Thus $\rho(h)$
preserves $L$.

For existence, choose a finite $K$-basis $b_1,\ldots,b_n$ of $V$ and
take $L_0=\sum_i\mathcal O b_i$. This submodule is finitely generated
and spans $V$ over $K$, so the construction applies.
\end{proof}

\begin{implementation}
\module{ModularRep.StableLattice} defines \lean{StableLattice} as a stable
subrepresentation satisfying Mathlib's \lean{IsLattice} predicate. The orbit sum is
\lean{StableLattice.orbitSpan}, represented by the supremum of the image
submodules. Its containment, stability, finite generation, and full span
are proved in \lean{le_orbitSpan}, \lean{orbitSpan_stable},
\lean{orbitSpan_fg}, and \lean{orbitSpan_isLattice}. The resulting lattice
is \lean{stableClosure}, and the basis construction is \lean{ofBasis}.
The construction \lean{ofBasis} also applies when $K$ is a commutative
ring and a finite basis is given. It requires neither a discrete valuation
nor completeness.
\end{implementation}

The integral representation on $L$ is obtained by restricting $\rho(g)$
to $L$ for each $g$. Once an $\mathcal O$-algebra structure on $k$ is
specified, define its reduction by
\[
 \overline L=k\otimes_{\mathcal O}L,
 \qquad
 g(a\otimes x)=a\otimes\rho(g)x.
\]
If $x_1,\ldots,x_t$ generate $L$ over $\mathcal O$, the tensors
$1\otimes x_1,\ldots,1\otimes x_t$ generate $\overline L$ over $k$.
Indeed expansion in the generators and the balancing relation move all
$\mathcal O$-coefficients to the first tensor factor. Hence the reduction
is finite dimensional whenever $k$ is a field.

\begin{implementation}
The declarations \lean{StableLattice.integralRepresentation} and
\lean{StableLattice.reduction} construct these actions.
\lean{reduction_apply_tmul} gives the action on pure tensors, and
\lean{reduction_moduleFinite} proves finite generation by the tensors of
an arbitrary finite generating set. In
\module{ModularRep.StableLatticeKZero}, \lean{chosenStableLattice}
uses a finite basis of an \lean{FDRep} object, and
\lean{stableLatticeReduction} is its finite dimensional reduction.
\end{implementation}

\section{Tensor reduction as an equivariant quotient}
\label{lib:reduction:quotient}

The quotient description of reduction is valid under weaker hypotheses
than those of a modular system. Let $\pi:\mathcal O\to k$ be a
surjective homomorphism of commutative rings, and write $I=\ker\pi$.
For an $\mathcal O$-module $M$, the notation $IM$ means the submodule
generated by the products $rx$ with $r\in I$ and $x\in M$.
The quotient $M/IM$ is a $k$-module through
$\mathcal O/I\simeq k$. Write $\overline x=x+IM$ for the class of $x\in M$.

\begin{proposition}[Reduction modulo an ideal]
\label{lib:reduction:tensor-quotient-theorem}
There is a $k$-linear isomorphism
\[
 \Phi:k\otimes_{\mathcal O}M\xrightarrow{\sim}M/IM,
 \qquad a\otimes x\longmapsto a\overline x,
 \qquad \Phi^{-1}(\overline x)=1\otimes x.
\]
If $M$ carries an $\mathcal O$-linear action of $G$, then $IM$ is stable,
the action descends by $g\overline x=\overline{gx}$, and $\Phi$ intertwines
the tensor and quotient actions.
\end{proposition}
\begin{proof}
Surjectivity gives the isomorphism $\mathcal O/I\simeq k$.
The standard tensor quotient isomorphism identifies
$(\mathcal O/I)\otimes_{\mathcal O}M$ with $M/IM$. Using
$\mathcal O/I\simeq k$ in the first tensor factor gives $\Phi$ with the stated
formula. The quotient relation makes $\overline x\mapsto1\otimes x$
well defined: if $r\in I$, then
$1\otimes rx=\pi(r)\otimes x=0$.

Every group operator is $\mathcal O$-linear, so it sends $rx$ to $r(gx)$
and preserves $IM$. This defines the quotient action directly. For a pure
tensor,
\[
 \Phi\bigl(g(a\otimes x)\bigr)=a\overline{gx}
   =g(a\overline x)=g\Phi(a\otimes x).
\]
Additivity extends the identity to the whole tensor product.
\end{proof}

\begin{implementation}
\module{ModularRep.ReductionModulo} defines $IM$ as
\lean{ReductionModulo.kernelSMul}. Its \lean{tensorQuotientEquiv}
combines Mathlib's \lean{TensorProduct.quotTensorEquivQuotSMul} with
its first isomorphism theorem for algebras. The local construction
transports the quotient scalar action and uses surjectivity to obtain
the $k$-linear equivalence. The versions for an explicit ring map and a
specified ideal are \lean{tensorQuotientEquivOfSurjectiveRingHom} and
\lean{tensorQuotientEquivOfKernelEq}.

\module{ModularRep.ReductionModuloEquivariance} first constructs
\lean{ReductionModulo.quotientRepresentation} by descending the integral
action. The intertwining identity is
\lean{tensorQuotientEquiv_isIntertwining}, and the resulting equivalence
of representations is \lean{tensorQuotientRepresentationEquiv}.
\end{implementation}

For a modular system, $I=\mathfrak m$ and the proposition gives
$\overline L\simeq L/\mathfrak mL$ as $kG$-modules. This is
\lean{stableLatticeReductionQuotientEquiv}. Its underlying linear
equivalence agrees with
\lean{stableLatticeReductionCarrierQuotientEquiv}.
For a fixed stable lattice, this isomorphism does not require $G$ to be finite.

\section{Character compatibility and lattice independence}
\label{lib:reduction:compatibility}

Fix a complete $\ell$-modular system. For the character construction in
this section, additionally assume that $k$ is algebraically closed and
choose a finite root correspondence $\iota$ for $G$. The field $K$ has
characteristic zero and is the fraction field of $\mathcal O$. Write
\[
 \operatorname{br}_{\iota}:K_0(kG)\longrightarrow
     \operatorname{CF}_{\ell'}(G,K),
 \qquad
 c:K_0(KG)\longrightarrow\operatorname{CF}_{\ell'}(G,K)
\]
for the Brauer character homomorphism and the homomorphism taking an
ordinary representation class to its trace character restricted to
$\ell$-regular elements. The first is injective by
Proposition~\ref{lib:kzero:brauer-injective}. The second is defined by
ordinary trace additivity and the universal property of the Grothendieck group.

\begin{hypothesis}[Reduction compatibility]
\label{lib:reduction:character-hypothesis}
Write $G_{\ell'}$ for the set of $\ell$-regular elements. For every finite
dimensional $K$-representation $V$, with ordinary trace character $\chi_V$,
and every stable $\mathcal O$-lattice $L$ in $V$, the fixed root
correspondence satisfies
\[
 \varphi_{k\otimes_{\mathcal O}L}
      =\chi_V\big|_{G_{\ell'}}.
\]
\end{hypothesis}

The reduction identity with compatible roots is standard
\cite[Corollary~2.9]{Navarro1998}.
Here the residue map and the root correspondence are specified
separately. Their compatibility is required: the identity is assumed
for every stable lattice with these same choices.

\begin{implementation}
This hypothesis is \lean{StableReductionBrauerCharacterCompatibility}.
It is an additional condition on \lean{ModularSystem} and
\lean{PrimeRegularRootEmbedding}.
\end{implementation}

Lattice independence and the decomposition homomorphism are treated in
\cite[\S15.2, Theorem~32 and the following discussion]{Serre1977}. That account proves lattice independence by
comparing lattices directly. It also notes that completeness is unnecessary
for that argument. Under the coefficient hypotheses and character identity
stated above, lattice independence also follows from injectivity of the
Brauer character homomorphism, as in the proof below.

\begin{theorem}[Decomposition homomorphism]
\label{lib:reduction:decomposition-theorem}
Under Hypothesis~\ref{lib:reduction:character-hypothesis}, there is a unique
homomorphism of abelian groups
\[
 d:K_0(KG)\longrightarrow K_0(kG)
 \quad\text{such that}\quad
 \operatorname{br}_{\iota}\circ d=c.
\]
For every stable lattice $L$ in $V$, it satisfies
$d([V])=[k\otimes_{\mathcal O}L]$. Consequently two stable lattices in
the same ordinary representation have equal modular Grothendieck classes.
\end{theorem}
\begin{proof}
Choose one stable lattice for every ordinary representation, using
Proposition~\ref{lib:reduction:existence}, and assign its reduction class
to the representation. If $L_1,L_2$ lie in the same $V$, compatibility gives
\[
 \operatorname{br}_{\iota}([\overline L_1])
  =c([V])
  =\operatorname{br}_{\iota}([\overline L_2]).
\]
Injectivity implies $[\overline L_1]=[\overline L_2]$. The same argument
using invariance of $c$ under isomorphism compares choices for isomorphic
ordinary representations.

For a short exact sequence with terms $U,V,W$, let $a_U,a_V,a_W$ be the
chosen reduction classes. Ordinary trace additivity gives
\[
 \operatorname{br}_{\iota}(a_V)
   =c([V])=c([U])+c([W])
   =\operatorname{br}_{\iota}(a_U+a_W).
\]
Injectivity gives $a_V=a_U+a_W$.
The universal property therefore defines $d$. On every object class,
compatibility gives $\operatorname{br}_{\iota}d=c$, so the two
homomorphisms agree everywhere. Comparing $d([V])$ with any other stable
lattice reduction proves the asserted formula for every $L$.
Finally, if $d'$ has the same character identity, then
$\operatorname{br}_{\iota}(d'(x))=\operatorname{br}_{\iota}(d(x))$
for every $x$. Injectivity gives $d'=d$.
\end{proof}

\begin{implementation}
\module{ModularRep.ExactCharacterBridge} proves the same assertion for
homomorphisms from the two Grothendieck groups to an abelian group $A$,
provided the modular homomorphism is injective and the character identity
holds for every stable lattice. The resulting definitions are
\lean{ExactCharacterBridge.toExactDecompositionMapData} and
\lean{ExactCharacterBridge.decompositionMap}.

\module{ModularRep.BrauerDecompositionMap} specialises this argument to
ordinary and Brauer characters. Its
\lean{exactCharacterBridgeOfStableReduction} uses
\lean{brauerCharacterKZeroHom_injective} for injectivity on $K_0(kG)$.
The resulting homomorphism is \lean{decompositionMapOfStableReduction},
its character identity is
\lean{decompositionMapOfStableReduction_characterSquare}, and its
uniqueness is \lean{existsUnique_decompositionMapOfStableReduction}.
The declarations
\lean{decompositionMap_classOf_eq_stableLatticeReduction} and
\lean{stableLatticeReductionClass_eq} give the two lattice statements.
\end{implementation}

By Theorem~\ref{lib:kzero:basis-theorem}, equality of these
classes is equivalent to equality of their composition multiplicities.

\begin{example}[Decomposition numbers]
\label{lib:reduction:coefficient-example}
For a simple $kG$-module $S$, consider the coefficient of $[S]$ in
$d([V])$ with respect to the basis of simple module classes. If $L$ is any
stable lattice in $V$, this coefficient is $m_S(\overline L)$ and is a nonnegative integer.
For a short exact sequence $0\to U\to V\to W\to0$, it is the sum of
the coefficients for $U$ and $W$. For irreducible $V$, these coefficients
are the decomposition numbers.
This definition makes sense over the coefficient fields fixed here. To
identify them with the usual character decomposition numbers for $G$, one
also chooses $K$ to be a splitting field for $G$.
\end{example}

\section{Twists and tensor products}
\label{lib:reduction:naturality}

For $\alpha\in\Aut(G)$, define the twist of $V$ by precomposition:
$\rho^{\alpha}(g)=\rho(\alpha(g))$. The induced operation on exact
$K_0$ is denoted by $T_{\alpha}$. It comes from an exact autoequivalence,
and the same notation is used over $K$ and $k$.

\begin{proposition}[Compatibility with automorphisms]
\label{lib:reduction:twist-theorem}
Under the hypotheses of Theorem~\ref{lib:reduction:decomposition-theorem},
\[
 d(T_{\alpha}x)=T_{\alpha}d(x)
 \qquad (\alpha\in\Aut(G),\ x\in K_0(KG)).
\]
\end{proposition}
\begin{proof}
Both character homomorphisms commute with precomposition by $\alpha$.
Applying $\operatorname{br}_{\iota}$ to the two sides therefore gives
the same function, namely $c(x)\circ\alpha$ on $G_{\ell'}$.
Injectivity proves the equality in $K_0(kG)$.
\end{proof}

\begin{implementation}
\module{ModularRep.KZeroTwist} defines \lean{FDRep.twistEquivalence}
by applying Mathlib's \lean{Action.resEquiv} to finite dimensional
representations. It constructs the induced map \lean{twistKZero} and proves
\lean{ordinaryCharacterKZero_twist} and
\lean{brauerCharacterKZeroHom_twist}. The consequence above is
\lean{decompositionMapOfStableReduction_twist}.
The convention is a right action: with automorphism multiplication given
by composition, \lean{twistKZero_comp} states
$T_{\alpha}\circ T_{\beta}=T_{\beta\alpha}$.
\end{implementation}

For a linear character $\lambda:G\to k^{\times}$, write $L_{\lambda}V$
for the representation on the same vector space with action
$g\cdot v=\lambda(g)\rho(g)v$. This is tensoring by the corresponding
one dimensional representation. Tensoring by $\lambda^{-1}$ is its
inverse, so it is an exact autoequivalence. Use $L_{\lambda}$ also for
the induced operation on $K_0$.

Twisting and tensoring satisfy
\[
 T_\alpha L_\lambda=L_{\lambda\circ\alpha}T_\alpha.
\]
The two composite exact functors are naturally isomorphic: the identity
on the representation space intertwines their equal actions. This
argument applies over either coefficient field and requires neither a
modular system nor a formula for Brauer characters of tensor products.
The natural isomorphism is
\lean{ModularRep.FDRep.linearCharacterTwistTwistIso}, and its consequence
on $K_0$ is \lean{ModularRep.twistKZero_linearCharacterTwistKZero}, both
in \module{ModularRep.KZeroLinearCharacterSemidirect}.

For a modular linear character $\lambda$, write $\widehat\lambda(g)$
for the lift of $\lambda(g)$ through $\iota$, for $g\in G_{\ell'}$.
Let $\lambda_K:G\to K^{\times}$ and $\lambda_k:G\to k^{\times}$
be linear characters, and assume that
$\lambda_K\big|_{G_{\ell'}}=\widehat{\lambda_k}$.
By Theorem~\ref{lib:brauer:tensor-source},
\[
 \varphi_{L_{\lambda}V}=\widehat\lambda\,\varphi_V
 \qquad(V\in\operatorname{FDRep}_k(G),\ \lambda:G\to k^{\times})
\]
on $G_{\ell'}$, with pointwise multiplication on the right.
For the general tensor product formula, see
\cite[proof of Theorem~2.23]{Navarro1998}.

\begin{proposition}[Compatibility with linear characters]
\label{lib:reduction:tensor-theorem}
If $d$ is the homomorphism of
Theorem~\ref{lib:reduction:decomposition-theorem} and $\lambda_K$ and
$\lambda_k$ satisfy the compatibility condition above, then, for every
$x\in K_0(KG)$,
\[
 d(L_{\lambda_K}x)=L_{\lambda_k}d(x).
\]
Consequently the combined operation also satisfies
$d(L_{\lambda_K}T_{\alpha}x)=L_{\lambda_k}T_{\alpha}d(x)$.
\end{proposition}
\begin{proof}
Ordinary trace on $L_{\lambda_K}V$ is the product of $\lambda_K$ with
the trace on $V$, since scalar multiplication by $\lambda_K(g)$ scales
the trace by the same factor. Theorem~\ref{lib:brauer:tensor-source}
gives the analogous statement over $k$. Compatibility of the two linear characters
therefore identifies the images under $\operatorname{br}_{\iota}$ of
$d(L_{\lambda_K}x)$ and $L_{\lambda_k}d(x)$. Injectivity gives the first
identity. Combining it with Proposition~\ref{lib:reduction:twist-theorem}
gives the second.
\end{proof}

\begin{implementation}
\module{ModularRep.KZeroLinearCharacterTensor} constructs the exact
autoequivalence and \lean{linearCharacterTwistKZero}, and proves the
ordinary trace formula. In
\module{ModularRep.KZeroLinearCharacterDecomposition},
\lean{LinearCharacterReductionCompatible} states the equality
$\lambda_K\big|_{G_{\ell'}}=\widehat{\lambda_k}$. The conclusions are
\lean{decompositionMapOfStableReduction_linearCharacterTwist} and
\lean{decompositionMapOfStableReduction_linearCharacterTwist_twist}.
Both the compatibility of the linear characters and
\lean{BrauerLinearTensorProductFormula} are hypotheses of these results.
\end{implementation}

\section{Specifying a decomposition homomorphism}
\label{lib:reduction:interfaces}

An assignment $[V]\mapsto a_V\in K_0(kG)$ that is additive on short
exact sequences induces a homomorphism $K_0(KG)\to K_0(kG)$. If
$a_V=[\overline L]$ for every stable lattice $L$ in $V$, this is the
decomposition homomorphism.
Theorem~\ref{lib:reduction:decomposition-theorem} proves these properties
from character compatibility.

\begin{implementation}
\lean{ExactDecompositionMapData} describes an assignment on ordinary
isomorphism classes that is additive on short exact sequences. Its
\lean{toAddMonoidHom} is the homomorphism given by the universal property.
The condition \lean{RealisesStableLatticeReduction} says that every
stable lattice has the assigned reduction class. For
\lean{decompositionMapOfStableReduction}, this condition is proved by
\lean{exactDecompositionMapDataOfStableReduction_realises}.

\lean{DecompositionMapInterface} describes a homomorphism between
specified ordinary and modular class groups that agrees with reduction
for every stable lattice. Equality of the classes from
two lattices then follows by expressing both as the image of the same
ordinary class.
\end{implementation}

\chapter{Central characters and blocks}
\label{chap:library:blocks}

A central idempotent gives a direct product decomposition of a ring.
For a group algebra, primitive central idempotents define the blocks.
Central subgroup idempotents relate the scalar action on a simple module
to its block. The Brauer homomorphism and block induction relate blocks
of different groups. They give Brauer's first main theorem under the
stated facts about defect groups. We also bound the group of linear
characters stabilising a block.

Throughout, $G$ is a finite group. Unless a stronger assumption is stated,
$k$ is a field. In arguments of positive characteristic, $\ell$ denotes a
prime and $\operatorname{char}k=\ell$. For
$z=\sum_{g\in G}z_g g\in kG$, the symbol $z_g$ denotes its coefficient at
$g$. A block means a primitive central idempotent of $kG$.
When blocks are indexed by a finite set $\mathcal B(G)$, their idempotents
are denoted by $b_B$, for $B\in\mathcal B(G)$. A complete primitive
decomposition and the central characters of its blocks will be assumed
where they are needed.

\section{Central idempotents and simple modules}
\label{lib:blocks:support}

The initial arguments apply to an arbitrary unital ring $R$.

\begin{definition}
A \emph{complete family of orthogonal central idempotents} is a finite family
$(e_i)_{i\in I}$ in $R$ such that each $e_i$ is central,
$e_i^2=e_i$, $e_i e_j=0$ for $i\ne j$, and $\sum_i e_i=1$.
A nonzero central idempotent $b$ is \emph{primitive} if every central
idempotent $c$ with $cb=c$ is either $0$ or $b$.
A module \emph{belongs to the block} $b$ if $b$ acts as the identity.
\end{definition}

\begin{proposition}[Central idempotents on simple modules]
\label{lib:blocks:unique-support}
Let $(e_i)_{i\in I}$ be a complete family of orthogonal central
idempotents in $R$.
Every simple $R$-module $V$ has a unique index $i$ such that $e_i$ acts
as the identity. Every other $e_j$ acts as zero.
Every primitive central idempotent $b$ likewise has a unique index $i$
such that
\[
 be_i=b,\qquad be_j=0\quad(j\ne i).
\]
If $b$ acts as the identity on $V$, the two indices coincide.
\end{proposition}
\begin{proof}
Centrality makes $v\mapsto e_i v$ an $R$-linear endomorphism. On a
simple module it is either zero or bijective. In the latter case its
idempotence implies that it is the identity: injectivity applied to
$e_i(e_i v)=e_i v$ gives $e_i v=v$. If every $e_i$ acted as zero, then
$1=\sum_i e_i$ would act as zero on the nonzero module $V$. Thus at least
one acts as the identity. If $e_i$ does, then
$e_jv=e_je_iv=0$ for $j\ne i$, proving uniqueness.

For the assertion about $b$, each $be_i$ is a central idempotent with
$(be_i)b=be_i$. Primitivity makes it either zero or $b$. The equation
$\sum_i be_i=b\ne0$ gives a nonzero term. If $be_i=b$, multiplication
by $e_j$ and orthogonality give $be_j=0$ for $j\ne i$.
Finally, if $b$ acts as the identity, then
$e_jv=b(e_jv)=(be_j)v$, so the two indices coincide.
\end{proof}

\begin{implementation}
The structure for a complete family of central idempotents extends Mathlib's
\lean{CompleteOrthogonalIdempotents}. The primitivity predicate is local,
and the proof for simple modules uses Mathlib's Schur lemma.
The two existence and uniqueness theorems are
\lean{ModularRep.CompleteOrthogonalCentralIdempotents.existsUnique_isSupport}
and \lean{ModularRep.IsPrimitiveCentralIdempotent.existsUnique_isSupport}.
The agreement of the two indices is
\lean{ModularRep.IsPrimitiveCentralIdempotent.IsSupport.toModuleSupport}.
Their modules are \module{ModularRep.CentralIdempotentSupport} and
\module{ModularRep.PrimitiveCentralIdempotent}.
\end{implementation}

A \emph{block decomposition} is a complete family of orthogonal
primitive central idempotents. For the expansion of central idempotents in
the group algebra case, see, for example,
\cite[Theorem~3.11]{Navarro1998}.

\begin{proposition}[Expansion in primitive idempotents]
\label{lib:blocks:idempotent-expansion}
Let $(b_B)_{B\in\mathcal B}$ be a block decomposition of $R$, and let
$e$ be a central idempotent. Set $S(e)=\{B\in\mathcal B:b_Be\ne0\}$. Then
\[
 b_Be\in\{0,b_B\},\qquad
 e=\sum_{B\in S(e)}b_B.
\]
The set $S(e)$ is the unique subset of $\mathcal B$ with this sum.
In particular, $e\ne0$ if and only if $S(e)$ is nonempty.
\end{proposition}
\begin{proof}
The product $b_Be$ is a central idempotent below $b_B$, so primitivity
gives the two alternatives. Consequently
\[
 e=1e=\sum_B b_Be=\sum_{B\in S(e)}b_B.
\]
Multiplication of a proposed sum by $b_C$ gives $b_C$ when $C$ is in
its index set and zero otherwise. Since $b_C\ne0$, this recovers the
index set from the sum and proves uniqueness. An empty sum is zero,
whereas a nonempty sum has a nonzero product with any of its members.
\end{proof}

\begin{implementation}
\module{ModularRep.CentralIdempotentBlockExpansion} proves the expansion
and its uniqueness in
\lean{ModularRep.CompleteOrthogonalCentralIdempotents.sum_centralIdempotentSupport_eq}
and \lean{ModularRep.CompleteOrthogonalCentralIdempotents.eq_centralIdempotentSupport_of_sum_eq}.
For a given block decomposition,
\module{ModularRep.BlockIdempotentDecomposition} defines
\lean{moduleBlock} and the equivalence \lean{primitiveBlockEquiv}
between the indexing set and the primitive central idempotents.
\end{implementation}

\begin{example}[Using a decomposition]
Given a simple representation and a block decomposition, the unique
idempotent acting as the identity determines its block. An isomorphism of modules
preserves this index because it preserves each equation $b_Bv=v$.
This is proved in
\lean{ModularRep.BlockIdempotentDecomposition.moduleBlock_eq_of_linearEquiv}.
\end{example}

\section{Idempotents of central subgroups}
\label{lib:blocks:central-sectors}

Let $Z\leq Z(G)$ have order invertible in $k$. For a linear character
$\nu:Z\to k^\times$, define
\[
 s_\nu=\sum_{z\in Z}\nu(z)^{-1}z,
 \qquad e_\nu=|Z|^{-1}s_\nu\in kG.
\]
The inverse in the coefficient ensures that $e_\nu$ acts as the identity
on a module with scalar action $zv=\nu(z)v$.

\begin{proposition}[Idempotents associated with central characters]
\label{lib:blocks:character-idempotents}
Each $e_\nu$ is a central idempotent, and $e_\nu e_\mu=0$ for
$\nu\ne\mu$. If $k$ is algebraically closed, then
\[
 \sum_{\nu\in\operatorname{Hom}(Z,k^\times)}e_\nu=1.
\]
Thus the idempotents form a complete family of orthogonal central idempotents.
\end{proposition}
\begin{proof}
For a nontrivial linear character $\chi$ of a finite group, choose $a$
with $\chi(a)\ne1$. Translation by $a$ permutes the group, so
\[
 \chi(a)\sum_z\chi(z)=\sum_z\chi(az)=\sum_z\chi(z).
\]
Cancellation in the field gives $\sum_z\chi(z)=0$. For the trivial
character the sum is the group order. Reindexing the product of the
weighted sums by $t=xy$ gives
\[
 s_\nu s_\mu
 =\left(\sum_{x\in Z}\nu(x)^{-1}\mu(x)\right)s_\mu.
\]
This is $|Z|s_\nu$ when $\nu=\mu$ and zero otherwise. Division by
$|Z|^2$ proves idempotence and orthogonality. Each group element in
the sum is central in $G$, which proves centrality in $kG$.

For completeness, finite abelian duality applies because $k$ has enough
roots of unity for the exponent of $Z$. Its character group has order
$|Z|$ and separates the elements of $Z$. If $z\ne1$, evaluation at
$z^{-1}$ is therefore a nontrivial character of this dual group. The
same translation argument gives
\[
 \sum_\nu\nu(z)^{-1}=0\quad(z\ne1),
 \qquad \sum_\nu\nu(1)^{-1}=|Z|.
\]
These are exactly the coefficient equations for $\sum_\nu e_\nu=1$.
Algebraic closure and invertibility of $|Z|$ provide the required roots.
\end{proof}

\begin{implementation}
The weighted product is
\lean{ModularRep.linearCharacterWeightedSum_mul} in
\module{ModularRep.LinearCharacterOrthogonality}.
The completeness proof is
\lean{ModularRep.sum_centralCharacterIdempotent_eq_one} in
\module{ModularRep.LinearCharacterCompleteness}. It uses Mathlib's finite
abelian duality to count the characters and separate group elements,
with the required roots of unity specified. The algebraically
closed version is proved in
\lean{ModularRep.completeCentralCharacterIdempotentsOfIsAlgClosed}.
\end{implementation}

For a finite dimensional irreducible representation $\rho$ over an
algebraically closed field, Schur's lemma gives a unique linear character
$\nu_\rho:Z\to k^\times$ such that
$\rho(z)=\nu_\rho(z)\id$. The block central character $\omega_B$
introduced below has domain $Z(kG)$, whereas $\nu_\rho$ has domain $Z$.

\begin{proposition}[The central character of a block]
\label{lib:blocks:sector}
Assume $k$ is algebraically closed and $|Z|$ is invertible in $k$.
Every primitive central idempotent $b$ lies above a unique character $\nu_b$
such that
\[
 be_{\nu_b}=b,\qquad be_\mu=0\quad(\mu\ne\nu_b).
\]
If $\rho$ is finite dimensional and irreducible, and $b$ acts as the
identity on $\rho$, then $\nu_\rho=\nu_b$. Hence any two such
representations have the same scalar action of $Z$.
\end{proposition}
\begin{proof}
Apply Proposition~\ref{lib:blocks:unique-support} to the complete
family of Proposition~\ref{lib:blocks:character-idempotents}.
On $\rho$, the action of $e_\mu$ is multiplication by
\[
 |Z|^{-1}\sum_{z\in Z}\mu(z)^{-1}\nu_\rho(z),
\]
which is one for $\mu=\nu_\rho$ and zero otherwise. Since $b$ acts as
the identity on $\rho$, Proposition~\ref{lib:blocks:unique-support}
gives $\nu_b=\nu_\rho$.
\end{proof}

\begin{implementation}
\lean{ModularRep.IsPrimitiveCentralIdempotent.existsUnique_centralCharacterSector}
and \lean{Representation.primitiveCentralIdempotentSector_eq_centralCharacter}
are in \module{ModularRep.CentralCharacterBlockSector}.
The scalar calculation is in \module{ModularRep.CentralCharacterAction}.
The equality of central characters assumes an irreducible representation
belonging to the block.
\end{implementation}

\begin{example}[Automorphisms and central characters]
Assume $k$ is algebraically closed. A ring isomorphism preserves
primitivity: pull back a central idempotent
below the image of $b$, apply primitivity of $b$, then map forward.
For a group automorphism $\alpha$, transport the basis of the group algebra by
$\alpha^{-1}$. If the order of $Z(G)$ is invertible in $k$, coefficient
comparison gives
\[
 e_\nu\longmapsto e_{\nu\circ\alpha|_{Z(G)}}.
\]
Thus a block lying above $\nu$ is sent to a block lying above
$\nu\circ\alpha|_{Z(G)}$. The corresponding theorem is
\lean{ModularRep.IsPrimitiveCentralIdempotent.centralCharacterSector_mapDomainRingEquiv_center}.
This example concerns the whole centre under the stated invertibility
hypothesis, as in \module{ModularRep.PrimitiveBlockAutomorphism}.
\end{example}

\section{Block central characters and induction}
\label{lib:blocks:induction}

Assume in this section that $k$ is algebraically closed, and fix a block
decomposition. The following standard description of block central
characters is used as an explicit input to the Lean argument.
In positive characteristic it is given by
\cite[Theorem~3.11]{Navarro1998}.

\begin{hypothesis}[Central characters of blocks]
\label{lib:blocks:catalogue}
For each $B\in\mathcal B(G)$ there is a $k$-algebra homomorphism
$\omega_B:Z(kG)\to k$ satisfying
\[
 \omega_B(b_C)=
 \begin{cases}1& B=C,\\0& B\ne C.\end{cases}
\]
Every $k$-algebra homomorphism $Z(kG)\to k$ is one of the $\omega_B$.
\end{hypothesis}

The evaluation equations imply that distinct blocks have distinct
central characters: evaluate an equality $\omega_B=\omega_C$ at $b_B$.
Thus $B\mapsto\omega_B$ is a bijection between the blocks and the
$k$-algebra homomorphisms $Z(kG)\to k$. The central character $\omega_B$
does not generally describe multiplication on the block algebra:
$zb_B=\omega_B(z)b_B$ need not hold for $z\in Z(kG)$.
\begin{implementation}
The hypotheses are expressed by
\lean{ModularRep.BlockCentralCharacterCatalogue}. The bijection is
\lean{ModularRep.BlockCentralCharacterCatalogue.centralCharacterEquiv}.
\end{implementation}

For $H\leq G$, coefficient restriction is the linear map
\[
 r_H:kG\longrightarrow kH,
 \qquad r_H\left(\sum_{g\in G}z_gg\right)=\sum_{h\in H}z_hh.
\]
It carries $Z(kG)$ into $Z(kH)$: the coefficients of a central
element are constant on $G$-conjugacy classes, so the restricted
coefficients are constant on $H$-conjugacy classes. Thus $r_H(z)$
commutes with every basis element of $kH$. Also $r_H(1)=1$.
Multiplicativity does not hold for a general subgroup. Even for an
abelian group and $H=1$, an element $g\ne1$ has
$r_H(g)=r_H(g^{-1})=0$ but $r_H(gg^{-1})=1$.

\begin{definition}
We use the definition of block induction by central characters
(see, for example, \cite[p.~87]{Navarro1998}). Given the block central characters for
$H$ and $G$, block induction from a block
$b\in\mathcal B(H)$ is \emph{defined} when the linear function
$\omega_b\circ r_H:Z(kG)\to k$ is multiplicative.
A block $B\in\mathcal B(G)$ is \emph{induced from} $b$ if this
condition holds and $\omega_B=\omega_b\circ r_H$.
\end{definition}

\begin{proposition}[Existence and uniqueness of an induced block]
\label{lib:blocks:induced-block}
If block induction from $b$ is defined, there is exactly one induced
block $B$.
\end{proposition}
\begin{proof}
The composite is $k$-linear, multiplicative and sends one to one. It is
therefore a $k$-algebra homomorphism. By
Hypothesis~\ref{lib:blocks:catalogue}, it equals $\omega_B$ for some $B$.
Uniqueness follows from injectivity of
$B\mapsto\omega_B$, proved by evaluation at block idempotents.
\end{proof}

\begin{implementation}
\module{ModularRep.GroupAlgebraCentralFunctions} uses Mathlib's group algebra
and its centre. It proves centrality of coefficient restriction and
constructs \lean{centerCoeffRestrict}.
\module{ModularRep.BlockInduction} defines
\lean{IsBlockInductionDefined}, proves
\lean{existsUnique_blockInducesTo}, and defines \lean{inducedBlock}.
The induced block is defined from multiplicativity and the given block
central characters.
The external source interfaces in
\module{ModularRep.SpathNavarroSourceAdapter} use an algebraically closed
field $k$ of characteristic $\ell$ that is algebraic over $\F_\ell$.
This additional hypothesis is recorded by
\lean{ModularRep.SpathCoefficientField}.
\end{implementation}

\begin{proposition}[Preservation on a central subgroup]
\label{lib:blocks:induction-sector}
Let $Z\leq Z(G)\cap H$, with $|Z|$ invertible in $k$.
Suppose $b\in\mathcal B(H)$ induces to $B\in\mathcal B(G)$.
Let $V$ and $W$ be finite dimensional irreducible representations of
$H$ and $G$, belonging to the blocks $b$ and $B$, respectively.
After identifying $Z$ with its copy in $H$, their scalar characters
on $Z$ agree.
\end{proposition}
\begin{proof}
For a central idempotent $e\in kG$, primitivity gives
$b_Be\in\{0,b_B\}$. Applying $\omega_B$ shows
\[
 \omega_B(e)=1\quad\Longleftrightarrow\quad b_Be=b_B.
\]
The reverse implication uses $\omega_B(b_B)=1$. In the forward
implication the alternative $b_Be=0$ would give $1=0$.

Choose $\nu$ so that $b_B$ lies above $\nu$ on $Z$. Since $Z\leq H$,
$e_\nu$ is supported on $H$, so $r_H(e_\nu)$ is the corresponding
character idempotent in
$kH$. In particular it is still idempotent. Induction gives
\[
 \omega_b(r_H(e_\nu))=\omega_B(e_\nu)=1.
\]
The same criterion in $kH$ gives
$b_b r_H(e_\nu)=b_b$. Hence $b_b$ also lies above $\nu$.
Proposition~\ref{lib:blocks:sector} identifies $\nu$ with
the scalar characters of $V$ and $W$.
\end{proof}

\begin{implementation}
The criterion for the value on a central idempotent is
\lean{ModularRep.BlockCentralCharacterCatalogue.centralIdempotent_value_one_iff}.
The conclusion is
\lean{Representation.centralCharacter_eq_comp_of_blockInducesTo}
in \module{ModularRep.CentralCharacterBlockInductionRestriction}.
The proof identifies $Z$ with its copy in $H$ and uses the resulting
equality of orders. It does not require $V$ to occur in the restriction
of $W$.
\end{implementation}

\section{The central Brauer map}
\label{lib:blocks:brauer-map}

Now assume $\operatorname{char}k=\ell$ and let $P\leq G$ be an
$\ell$-subgroup. Write $C=C_G(P)$. The \emph{central Brauer map}
is coefficient restriction to $C$ on the centre of $kG$:
\[
 \operatorname{Br}_P:Z(kG)\longrightarrow Z(kC),
 \qquad \operatorname{Br}_P(z)=\sum_{c\in C}z_cc.
\]
The target is the centre of $kC$, by the restriction argument above.
Its multiplicativity follows from an orbit calculation. For this case
$H=C_G(P)$, see, for example, \cite[Theorem~4.9]{Navarro1998}.

\begin{lemma}[Summing over fixed points]
\label{lib:blocks:orbit-sum}
Let an $\ell$-group $Q$ act on a finite set $X$. If $f:X\to k$ is
constant on each $Q$-orbit, then
\[
 \sum_{x\in X}f(x)=\sum_{x\in X^Q}f(x).
\]
\end{lemma}
\begin{proof}
Partition $X$ into orbits. On an orbit $\mathcal O$, the sum equals
$|\mathcal O|f(x)$ for any $x\in\mathcal O$. Orbit sizes are powers
of $\ell$, so an orbit of size greater than one contributes zero in
characteristic $\ell$. An orbit of size one is exactly a fixed point
and contributes its own value.
\end{proof}

\begin{theorem}[The Brauer map on the centre]
\label{lib:blocks:brauer-multiplicative}
For a finite group $G$, an $\ell$-subgroup $P$ and a field $k$ of
characteristic $\ell$, the map $\operatorname{Br}_P$ is a
$k$-algebra homomorphism. Algebraic closure is not required.
\end{theorem}
\begin{proof}
Fix $c\in C$. The coefficient of $zw$ at $c$ is
\[
 (zw)_c=\sum_{a\in G}z_a w_{a^{-1}c}.
\]
Let $P$ act on $G$ by conjugation. If $q\in P$, then $qc=cq$, so
\[
 (qaq^{-1})^{-1}c=q(a^{-1}c)q^{-1}.
\]
Centrality of $z$ and $w$ makes the function
$a\mapsto z_aw_{a^{-1}c}$ constant on these orbits. Its fixed points
are precisely $a\in C$. Lemma~\ref{lib:blocks:orbit-sum} gives
\[
 (zw)_c=\sum_{a\in C}z_aw_{a^{-1}c}
       =\bigl(\operatorname{Br}_P(z)\operatorname{Br}_P(w)\bigr)_c.
\]
Here $a^{-1}c\in C$ whenever $a,c\in C$, so the final expression is
the product coefficient in $kC$. Equality at every $c\in C$ proves
multiplicativity. Linearity and preservation of one were already proved
for coefficient restriction.
\end{proof}

\begin{implementation}
The proof uses Mathlib's theorem that orbit sizes for an $\ell$-group
are powers of $\ell$. The resulting sum formula is
\lean{ModularRep.IsPGroup.sum_eq_sum_fixedPoints_of_invariant}
in \module{ModularRep.PGroupInvariantSum}.
Multiplicativity and the algebra homomorphism are
\lean{ModularRep.centralBrauerRestriction_mul} and
\lean{ModularRep.centralBrauerMap} in
\module{ModularRep.GroupAlgebraCentralBrauerMap}.
These declarations allow any commutative coefficient semiring of
characteristic $\ell$.
\end{implementation}

\begin{proposition}[The Brauer map to an intermediate subgroup]
\label{lib:blocks:interval-map}
Suppose $C_G(P)\leq H\leq N_G(P)$. Extend the coefficients of
$\operatorname{Br}_P(z)$ by zero from $C_G(P)$ to $H$.
The resulting map
\[
 \operatorname{Br}_P^H:Z(kG)\longrightarrow Z(kH)
\]
is a $k$-algebra homomorphism. Moreover,
$\operatorname{Br}_P(z)\ne0$ if and only if
$\operatorname{Br}_P^H(z)\ne0$.
\end{proposition}
\begin{proof}
The inclusion $kC_G(P)\hookrightarrow kH$ is an injective algebra
homomorphism. Every $h\in H$ normalises $P$, and therefore preserves
$C_G(P)$ by conjugation. The coefficients $z_c$ are constant
under this conjugation because $z\in Z(kG)$. The extended element
thus lies in $Z(kH)$. Multiplication and the unit are preserved by the
composite. Injectivity gives the asserted equivalence of nonvanishing.
\end{proof}

For this form of the Brauer homomorphism, see, for example,
\cite[Theorem~4.9]{Navarro1998}. Its image is central in $kH$
because the restricted coefficients are invariant under $H$.
An arbitrary element of $Z(kC_G(P))$ need not have this property.

\begin{implementation}
The homomorphism \lean{ModularRep.centralBrauerMapTo} is in
\module{ModularRep.GroupAlgebraCentralBrauerMapTo}.
The equivalence of nonvanishing,
\lean{ModularRep.centralBrauerMap_ne_zero_iff_mapTo_ne_zero},
is in \module{ModularRep.CentralBrauerIntervalSupport}.
For $H=N_G(P)$, the map is named
\lean{ModularRep.normalizerCentralBrauerMap}.
\end{implementation}

\begin{example}[A local sum of blocks]
The image of a primitive central idempotent of $kG$ under
$\operatorname{Br}_P^H$ is a central idempotent, although it can be
zero or a sum of several primitive central idempotents of $kH$. A block
decomposition of $kH$ and
Proposition~\ref{lib:blocks:idempotent-expansion} determine its unique
support. This is implemented by
\lean{ModularRep.sum_centralBrauerMapToLocalSupport_eq} and
\lean{ModularRep.centralBrauerMapToPrimitiveImage_ne_zero_iff_support_nonempty}.
In particular, the Brauer homomorphism need not preserve primitivity.
\end{example}

\section{Brauer's first main theorem}
\label{lib:blocks:normaliser-correspondence}

Brauer's first main theorem identifies the blocks of $G$ with defect
group $P$ with the blocks of $N_G(P)$ having defect group $P$.
See, for example, \cite[Theorem~4.17]{Navarro1998}.
We prove the correspondence from the central character identity and the
five facts about defect groups stated below.
Assume $k$ is algebraically closed of
characteristic $\ell$. Fix an $\ell$-subgroup $P$, write
$N=N_G(P)$, and fix block decompositions and block central characters
for $N$ and $G$ as in Hypothesis~\ref{lib:blocks:catalogue}.

We first use the central character identity
(see, for example, \cite[Theorem~4.14]{Navarro1998}). For
$PC_G(P)\leq H\leq N_G(P)$, it is
\begin{equation}
 \omega_b\circ r_H=\omega_b\circ\operatorname{Br}_P^H
 \qquad(b\in\mathcal B(H)).
 \label{lib:blocks:414-equality}
\end{equation}
The hypothesis includes $P\leq H$, whereas the construction of
$\operatorname{Br}_P^H$ only required $C_G(P)\leq H\leq N_G(P)$.
Since the right side of \eqref{lib:blocks:414-equality} is an algebra
homomorphism, block induction is defined. For $H=N$, denote the induced
block map by
$I:\mathcal B(N)\to\mathcal B(G)$. Thus
\begin{equation}
 \omega_{I(b)}=\omega_b\circ\operatorname{Br}_P^N.
 \label{lib:blocks:induction-brauer-character}
\end{equation}

For the following identity, see, for example, the final assertion of
\cite[Theorem~4.14]{Navarro1998}, with $H=N_G(P)$.

\begin{proposition}[The Brauer image as a sum of blocks]
\label{lib:blocks:fibre}
Under \eqref{lib:blocks:414-equality}, for every block $B$ of $G$,
\[
 \operatorname{Br}_P^N(b_B)=\sum_{b\in\mathcal B(N):\ I(b)=B}b_b.
\]
In particular, if the image is nonzero, some block of $N$ induces to $B$.
\end{proposition}
\begin{proof}
Expand the image as in Proposition~\ref{lib:blocks:idempotent-expansion}.
A local block $b$ lies in its support exactly when its central
character takes value one on the image. Indeed, multiplication by $b_b$
is either zero or $b_b$, and evaluation by $\omega_b$ distinguishes
these alternatives. Equation~\eqref{lib:blocks:induction-brauer-character}
identifies this value with $\omega_{I(b)}(b_B)$. The evaluation equations
for the ambient block central characters say that this is one exactly
when $I(b)=B$.
\end{proof}

\begin{implementation}
\module{ModularRep.Navarro414IntervalCentralCharacterSource} expresses
\eqref{lib:blocks:414-equality}, from which the induced block is defined.
The equivalence between occurrence in this sum and induction to $B$ is
\lean{ModularRep.mem_navarro417PhaseANormalizerSupport_iff_inductionMap_eq}
in \module{ModularRep.Navarro417PhaseASupport}.
\end{implementation}

For the Brauer map characterisation of defect groups, see, for example,
\cite[Theorem~4.11]{Navarro1998}. For a block $B$ of a finite group
$X$, let $\mathcal D_X(B)$ be the set of $\ell$-subgroups $D\leq X$
such that
\begin{equation}
 \operatorname{Br}_Q(b_B)\ne0
 \quad\Longleftrightarrow\quad
 Q\text{ is conjugate to a subgroup of }D
 \quad\text{for every }\ell\text{-subgroup }Q\leq X.
 \label{lib:blocks:defect-support}
\end{equation}
In the following argument, we use this criterion to define
$\mathcal D_X(B)$. Its members are the defect groups of $B$.
Any two members are conjugate: the criterion makes each conjugate to
a subgroup of the other, and equality of their finite orders then gives
conjugacy. Existence will be used explicitly below.

\begin{hypothesis}[Standard facts for the defect correspondence]
\label{lib:blocks:defect-hypotheses}
We use the following five standard facts about the fixed block
decompositions and central characters.
\begin{enumerate}
\item If $P\in\mathcal D_N(b)$ and an $N$-conjugacy class
$\mathcal K$ has nonzero coefficient in $b_b$, then
$\mathcal K=\mathcal L\cap C_G(P)$ for some $G$-conjugacy class
$\mathcal L$, with the intersection viewed inside $N$.
\item The set $\mathcal D_N(b)$ is nonempty for every block $b$ of $N$.
\item If $D\in\mathcal D_N(b)$, then $O_\ell(N)\leq D$.
\item If $b$ induces to $B$ and $D\in\mathcal D_N(b)$, then there is
$E\in\mathcal D_G(B)$ containing $D$, viewed as a subgroup of $G$.
\item If $P\in\mathcal D_N(b)$, then there is
$E\in\mathcal D_G(I(b))$ with $E\leq P$.
\end{enumerate}
\end{hypothesis}

For the first condition, Lemma~4.15 restricts the nonzero class
coefficients to classes with defect group $P$, and Lemma~4.16 writes
each such class as the required intersection
\cite[pp.~88--89]{Navarro1998}. The second condition follows from
existence of defect groups and their Brauer map characterisation
\cite[Theorems~4.3 and~4.11]{Navarro1998}.
The third and fourth conditions follow from
\cite[Theorem~4.8 and Lemma~4.13]{Navarro1998},
respectively.

The fifth condition is the containment proved in the first paragraph
of the proof of \cite[Theorem~4.17]{Navarro1998}. That argument
chooses a defect class of $b$, uses Corollary~4.5 and Lemma~4.16 to
lift it to an ambient class with defect group $P$, and applies
Theorems~4.14 and~4.4 to obtain a defect group of $I(b)$ contained in
$P$. We use this containment, together with the other four conditions
and \eqref{lib:blocks:414-equality}, to prove the correspondence.

\begin{proposition}[Lifting a local block idempotent]
\label{lib:blocks:class-lift}
Assume the first condition of
Hypothesis~\ref{lib:blocks:defect-hypotheses} and
\eqref{lib:blocks:414-equality}. If $P\in\mathcal D_N(b)$, then
$b_b$ is in the image of $\operatorname{Br}_P^N$.
If $I(b)=I(c)$ for any local block $c$, then $b=c$. Consequently,
\[
 \operatorname{Br}_P^N(b_{I(b)})=b_b.
\]
\end{proposition}
\begin{proof}
For an $N$-class $\mathcal K$, write
$\mathcal K^+=\sum_{x\in\mathcal K}x$. Centrality gives a unique
class expansion
\[
 b_b=\sum_{\mathcal K}a_{\mathcal K}\mathcal K^+.
\]
For each nonzero coefficient choose an ambient class
$\mathcal L_{\mathcal K}$ with
$\mathcal L_{\mathcal K}\cap C_G(P)=\mathcal K$.
Coefficient restriction gives
$\operatorname{Br}_P^N(\mathcal L_{\mathcal K}^+)=\mathcal K^+$.
Thus the finite sum
\[
 z=\sum_{a_{\mathcal K}\ne0}
       a_{\mathcal K}\mathcal L_{\mathcal K}^+\in Z(kG)
 \quad\text{satisfies}\quad \operatorname{Br}_P^N(z)=b_b.
\]

If $I(b)=I(c)$, equation~\eqref{lib:blocks:induction-brauer-character}
gives equal composites
$\omega_b\operatorname{Br}_P^N=\omega_c\operatorname{Br}_P^N$.
Evaluation at $z$ gives $\omega_b(b_b)=\omega_c(b_b)$. The left side
is one, and the right side is zero unless $b=c$. This proves
injectivity even when only the first block is known to have defect
$P$. In the sum of Proposition~\ref{lib:blocks:fibre}, the only block
inducing to $I(b)$ is therefore $b$. This gives the final equation.
\end{proof}

\begin{implementation}
The element $z$ is constructed in
\module{ModularRep.Navarro415416ClassSumAdapter}, in the theorem
\lean{ModularRep.Navarro415416RawClassSource.localBlockIdempotent_mem_normalizerCentralBrauerMap_range}.
Injectivity and the equality of block idempotents are
\lean{ModularRep.navarro417PhaseAInductionMap_eq_imp_eq_of_localP_left}
and \lean{ModularRep.normalizerCentralBrauerMap_selectedBlockIdempotent_eq}
in \module{ModularRep.Navarro417RestrictedInjectivity}.
The class intersection property and the central character equality
suffice for all three conclusions.
\end{implementation}

\begin{theorem}[Blocks with a prescribed defect group]
\label{lib:blocks:exact-defect}
Assume \eqref{lib:blocks:414-equality} and
Hypothesis~\ref{lib:blocks:defect-hypotheses}. Write
$\mathcal B(X\mid P)=\{B\in\mathcal B(X):P\in\mathcal D_X(B)\}$,
where $X=G$ or $N_G(P)$ and $P$ is viewed as a subgroup of $X$.
Then block induction is a bijection
\[
 I:\mathcal B(N_G(P)\mid P)\longrightarrow\mathcal B(G\mid P).
\]
\end{theorem}
\begin{proof}
Let $b\in\mathcal B(N\mid P)$. The fourth condition gives
$E_+\in\mathcal D_G(I(b))$ with $P\leq E_+$.
The fifth gives $E_-\in\mathcal D_G(I(b))$ with $E_-\leq P$.
They are conjugate by
\eqref{lib:blocks:defect-support}, so
\[
 |P|\leq|E_+|=|E_-|\leq|P|.
\]
The inclusions are therefore equalities. In particular,
$P\in\mathcal D_G(I(b))$.

Conversely, suppose $I(b)=B$ and $P\in\mathcal D_G(B)$.
By the second condition, choose $D\in\mathcal D_N(b)$. Since
$P$ is a normal $\ell$-subgroup of $N$, the third condition gives
\[
 P\leq O_\ell(N)\leq D.
\]
The fourth condition gives $E\in\mathcal D_G(B)$ with
$D\leq E$ inside $G$. The defect groups $E$ and $P$ are conjugate. Hence
$|P|\leq|D|\leq|E|=|P|$, and the local inclusion is an equality.
Thus $P\in\mathcal D_N(b)$ for every local block inducing to $B$.

Finally, if $P\in\mathcal D_G(B)$, then
$\operatorname{Br}_P(b_B)\ne0$ by
\eqref{lib:blocks:defect-support}. Proposition~\ref{lib:blocks:interval-map}
preserves this nonvanishing in $kN$, and
Proposition~\ref{lib:blocks:fibre} supplies a local block over $B$.
The preceding paragraph puts that block in $\mathcal B(N\mid P)$,
proving surjectivity. Injectivity follows from
Proposition~\ref{lib:blocks:class-lift}.
\end{proof}

\begin{implementation}
\module{ModularRep.Navarro417ExactDefectAdapter} proves the two
containment arguments and the bijection as
\lean{navarro417PhaseAInductionMap_has_ambientP},
\lean{navarro417LocalP_of_blockInducesTo_ambientP}, and
\lean{navarro417ExactDefectMap_bijective}, all in namespace
\lean{ModularRep}. The resulting equivalence is
\lean{ModularRep.navarro417ExactDefectEquiv}.
The construction takes the five conditions above and
\eqref{lib:blocks:414-equality}, for the given block decompositions
and central characters. The third condition, containment of the normal
$\ell$-core in every defect group, is proved by
\lean{navarro408PCoreDefectSource} in
\module{ModularRep.PaperProofs.SporadicFi24P3Definition44NamedCarrierPCoreDefectContainment}.
The remaining source
conditions used by this construction are stated separately.
\end{implementation}

\section{Linear characters stabilising a block}
\label{lib:blocks:linear-stabilisers}

Let $N\trianglelefteq G$ and assume again that $k$ is algebraically
closed of characteristic $\ell$. Write $\pi:G\to G/NZ(G)$ for the
quotient homomorphism. A linear character
$\lambda:G\to k^\times$ twists a representation by
\[
 (\lambda\otimes\rho)(g)=\lambda(g)\rho(g).
\]
If a central subgroup $Z$ acts on $\rho$ by a scalar character
$\nu_\rho$, it acts on $\lambda\otimes\rho$ by
$\lambda|_Z\nu_\rho$. This follows from the scalar action and does
not require division by the degree of $\rho$.

Suppose a group $A$ acts on the primitive central idempotents of $kG$, and there
is an injective homomorphism
$a\mapsto\lambda_a$ from $A$ to the linear characters of $G$.
Assume every $\lambda_a$ is trivial on $N$. Suppose every block contains
a finite dimensional simple representation, and tensoring any such
representation in $b$ by $\lambda_a$ gives a representation in
$a\cdot b$.

\begin{theorem}[Stabilisers under tensoring by linear characters]
\label{lib:blocks:stabiliser-bound}
Under these assumptions, the stabiliser
$A_b=\{a\in A:a\cdot b=b\}$ is finite. For every $a\in A_b$, the
character $\lambda_a$ factors through $G/NZ(G)$, and
\[
 |A_b|\leq [G:NZ(G)].
\]
\end{theorem}
\begin{proof}
The central $\ell$-regular elements form a subgroup
$Z_{\ell'}\leq Z(G)$ of order prime to $\ell$.
Choose a simple representation $\rho$ belonging to $b$. If $a\in A_b$,
then $\lambda_a\otimes\rho$ belongs to the same block.
Proposition~\ref{lib:blocks:sector} gives equality of their scalar characters on
$Z_{\ell'}$. Since twisting multiplies that character by
$\lambda_a|_{Z_{\ell'}}$, cancellation in $k^\times$ proves that
$\lambda_a$ is trivial on $Z_{\ell'}$.

It is also trivial on every $\ell$-element. If $g^{\ell^m}=1$, then
$\lambda_a(g)^{\ell^m}=1$. In characteristic $\ell$,
$X^{\ell^m}-1=(X-1)^{\ell^m}$, so a field has no nontrivial
$\ell$-power roots of unity. Every central element is the product
of its commuting $\ell$-part and $\ell$-regular part, both powers
of the original element. Hence $\lambda_a$ is trivial on all of
$Z(G)$, and therefore on $NZ(G)$.

Thus $\lambda_a$ induces a character
$\overline\lambda_a:G/NZ(G)\to k^\times$ with
$\lambda_a=\overline\lambda_a\circ\pi$. These quotient characters
are distinct, since composing with $\pi$ recovers $\lambda_a$ and
$a\mapsto\lambda_a$ is injective. These quotient characters form a
linearly independent family in the space of all functions on
$G/NZ(G)$, which has dimension $[G:NZ(G)]$. Therefore their indexing
set $A_b$ is finite and $|A_b|\leq [G:NZ(G)]$.
\end{proof}

\begin{implementation}
The assumptions on simple representations and tensoring are expressed by
\lean{ModularRep.LinearBlockStabilizer.SimpleBlockTensorSource}.
It supplies a simple representation in every block and compatibility
between the block action and tensoring. The theorem is
\lean{ModularRep.LinearBlockStabilizer.modular_block_stabilizer_bound}.
The quotient character is defined by \lean{descendCenterQuotient}.
Mathlib's linear independence of multiplicative characters gives
\lean{finite_of_injective_linear_characters} and the bound
\lean{card_le_of_injective_linear_characters}, in the same namespace.
The theorem \lean{block_stabilizer_finite} explicitly supplies finiteness
of the stabiliser for both the ordinary and modular character actions.
\end{implementation}

\begin{example}[All quotient characters]
Take $A$ to be the linear characters of $G$ trivial on $N$, identified
with those of $G/N$. Once their block action and its compatibility with
tensoring are given, every character stabilising a block factors through
$G/NZ(G)$, and the stated bound follows. In the special case $NZ(G)=G$,
the bound is one, so every block has trivial stabiliser for this
action.
\end{example}

For the ordinary form, let $K$ be a field of characteristic zero.
Set $m=|G|_{\ell'}$ and fix an isomorphism
$\mu_m(k)\simeq\mu_m(K)$ between the groups of $m$th roots of unity.
Applied pointwise, it identifies ordinary linear characters of $G/N$
of order prime to $\ell$ with modular linear characters of $G/N$.
Transport the block action through this identification. The stabilisers
are isomorphic, and preservation of kernels gives the factorisation
through $G/NZ(G)$ and the same bound.

\begin{implementation}
\module{ModularRep.LinearBlockStabilizerReduction} constructs
\lean{ordinaryToBrauerQuotient}, \lean{reductionStabilizerEquiv}, and
\lean{ordinary_block_stabilizer_bound} in namespace
\lean{ModularRep.LinearBlockStabilizer}.
The theorem \lean{reduction_commutes_with_tensor} requires a modular system,
\lean{StableReductionBrauerCharacterCompatibility} and
\lean{BrauerLinearTensorProductFormula}.
\end{implementation}

\chapter{Weights and changes of representatives}
\label{chap:library:weights}

Let $G$ be a finite group and let $\ell$ be a prime. A weight is a
local representation associated with an $\ell$-subgroup of $G$.
The modular formulation uses simple modular
representations that are projective over the normaliser quotient
\cite[p.~369]{Alperin1987}. We use the ordinary character formulation.
See, for example, \cite[p.~90]{Navarro1998}.

Throughout this chapter, $K$ is a field of characteristic zero.
Write $\Irr_K(H)$ for the trace characters of irreducible finite
dimensional $K$-representations of a finite group $H$.
The weight definitions below use these trace characters. They agree
with the usual ordinary character definition of weights when $K$
is a splitting field for the local quotients $N_G(Q)/Q$.

\section{Weights and their local representations}

\begin{definition}
A character $\theta\in\Irr_K(H)$ has \emph{defect zero} if it is
afforded by a simple $K H$-module $V$ such that
\[
 (\dim_K V)_\ell=|H|_\ell.
\]
Denote the set of these characters by
$\operatorname{dz}_\ell(H)$. Since $\theta(1)=\dim_K V$, the numerical
condition is equivalent to $\theta(1)_\ell=|H|_\ell$.

In \emph{character form}, a weight of $G$ is a pair $(Q,\theta)$, where $Q$ is
an $\ell$-radical subgroup of $G$ and
\[
 \theta\in\operatorname{dz}_\ell(N_G(Q)/Q).
\]
In \emph{representation form}, a weight is a pair $(Q,V)$ with the same
condition on $Q$, where $V$ is a simple finite dimensional
$K[N_G(Q)/Q]$-module of defect zero.
\end{definition}

At a fixed subgroup $Q$, two weights in representation form are identified
when their local representations are isomorphic. Two weights in character form
are identified when their local character functions are equal.
The condition that $Q$ is radical is included in these definitions.

\begin{proposition}[Comparison by characters]
Taking characters gives a bijection from weights in representation form up to
local isomorphism to weights in character form. It induces an
$\Aut(G)$-equivariant bijection on $G$-conjugacy classes of weights.
\end{proposition}
\begin{proof}
Trace is invariant under isomorphism, so the map is well defined.
Character separation for finite groups over fields of characteristic
zero gives injectivity, as in Section~\ref{lib:ordinary:realisations}.
A simple representation affording a character of defect zero gives
surjectivity. Transport along a group isomorphism preserves dimension
and simplicity, and taking traces commutes with this transport.
The bijection therefore commutes with automorphisms and descends to
conjugacy classes.
\end{proof}

\begin{implementation}
The local representations use Mathlib's \lean{FDRep}, with
\lean{FDRep.char_iso} giving invariance of the trace under isomorphism.
The Lean comparison maps take
\lean{ModularRep.LocalSplittingCharacterInput} in
\module{ModularRep.WeightCharacterBridge} as an explicit argument.
Its fields require that the trace of every simple local representation
has an irreducible realisation and that equal traces determine simple
local representations up to isomorphism. These properties follow
mathematically from the definition of $\Irr_K$ and the character
separation argument in Section~\ref{lib:ordinary:realisations}.
The record does not require $K$ to be a splitting field.
The character map is
\lean{ModularRep.LocalSplittingCharacterInput.characterOfSimple}.
The two bijections are
\lean{ModularRep.CharacterWeight.isoClassEquiv} and
\lean{ModularRep.CharacterWeight.conjugacyClassEquiv}, for
\lean{ModularRep.CharacterWeight}.
\end{implementation}

\section{Automorphisms and conjugacy classes}

An isomorphism $\phi:H\to H'$ transports a representation by composition
with $\phi^{-1}$. Its character is therefore $\theta\circ \phi^{-1}$.
This is also the convention for transport of class functions in
\cite[Problem~2.1]{Navarro1998}.
In particular, $\alpha\in\Aut(G)$ sends a weight at $Q$ to a weight at
$\alpha(Q)$ through the induced isomorphism
\[
 \bar\alpha:N_G(Q)/Q\longrightarrow
 N_G(\alpha(Q))/\alpha(Q).
\]
The transported local character is $\theta\circ\bar\alpha^{-1}$.
Composition of these transports gives the usual left action of
$\Aut(G)$. The corresponding right action uses $\alpha^{-1}$, so its
subgroup coordinate is $\alpha^{-1}(Q)$.
\begin{implementation}
\lean{ModularRep.OrdinaryIrreducibleCharacter.mapEquiv_apply}
gives the formula $\theta(\phi^{-1}h')$.
Representation transport uses Mathlib's restriction functor
\lean{Action.res} and the equivalence \lean{Action.resEquiv}.
ModularRep constructs the induced normaliser quotient maps and proves
that the resulting weight transports preserve radicality and defect zero.
The actions on weights in representation and character form are defined before
passing to their classes.
\end{implementation}

Conjugation in $G$ gives an equivalence relation on weights.
Write $\Alp(G)$ for the set of weights in character form modulo this relation.
Automorphisms act on $\Alp(G)$ because they preserve conjugacy.
Local isomorphism and conjugacy in $G$ have different roles:
local isomorphism keeps $Q$ fixed, whereas conjugation may replace it
by another subgroup.

\begin{proposition}[Weights with a fixed radical subgroup]
\label{lib:weights:representative}
Fix an $\ell$-radical subgroup $Q\leq G$. Let $\Alp(G)_{[Q]}$ be the
set of $G$-conjugacy classes of weights in character form whose radical
subgroup is conjugate to $Q$. Then
\[
 \operatorname{dz}_\ell(N_G(Q)/Q)\longrightarrow\Alp(G)_{[Q]},
 \qquad \theta\longmapsto[(Q,\theta)]
\]
is a bijection.
\end{proposition}
\begin{proof}
Every class in $\Alp(G)_{[Q]}$ has a representative at $Q$, obtained
by conjugating its radical subgroup to $Q$. This proves surjectivity.

If $(Q,\theta)$ and $(Q,\eta)$ are conjugate by $g\in G$, then
$g\in N_G(Q)$. The induced automorphism of $N_G(Q)/Q$ is conjugation
by $gQ$. Characters are constant on conjugacy classes, so it fixes
$\theta$. Hence $\eta=\theta$, proving injectivity. The same argument
shows that the character obtained by transporting to $Q$ is
independent of the chosen conjugating element.
\end{proof}
\begin{implementation}
The bijection is
\lean{ModularRep.CharacterWeight.localDefectZeroEquivWeightRadicalFibre}.
Its restriction to a block is
\lean{representativeDZEquivWeightBlockRadicalFibre} in the same
namespace, with the local block induction hypotheses described below.
\end{implementation}

\section{The block of a weight}

Fix compatible ordinary and modular character conventions for the
local groups, including the splitting and root identifications needed
for their block assignments.
Let $(Q,\theta)$ be a weight. Inflate $\theta$ to $N_G(Q)$ and let
$b$ be its block. The weight belongs to an ambient block $B$ when
Brauer block induction is defined and $b^G=B$
\cite[p.~90]{Navarro1998}.

Recall the definition of this induction. If $H\leq G$ and
$\lambda_b:Z(kH)\to k$ is the central character of a block $b$,
extend it linearly to $Z(kG)$ by
\[
 \lambda_b^G(C^+)=\lambda_b((C\cap H)^+)
\]
for each conjugacy class $C$ of $G$. Here $C^+$ denotes the sum of
the elements of $C$ in the group algebra, and $k$ is algebraically
closed of characteristic $\ell$.
If $\lambda_b^G$ is an algebra homomorphism, it is the central
character of a unique block, denoted by $b^G$.
This is Brauer's definition of induced block
\cite[p.~87]{Navarro1998}.
For an $\ell$-subgroup $P$, the inclusions
\[
 P C_G(P)\leq H\leq N_G(P)
\]
ensure that induction is defined for every block of $H$
\cite[Theorem~4.14]{Navarro1998}.

Write $\operatorname{Bl}_\ell(G)$ for the set of $\ell$-blocks of $G$.
Suppose the local block assignment, inflation and block induction
commute with automorphisms. Suppose also that inner automorphisms fix
ambient blocks. Then the ambient block assigned to a weight depends
only on its $G$-conjugacy class, and the resulting map
$\Alp(G)\to\operatorname{Bl}_\ell(G)$ is equivariant.
The corresponding compatibility of Brauer block induction with group
isomorphisms is stated in \cite[Problem~4.4]{Navarro1998}.

\begin{implementation}
\lean{ModularRep.CharacterWeight.LocalBlockInductionOperations} in
\module{ModularRep.CharacterWeightBlockAssignment} supplies the local
character block and inflation. It also supplies finite complete block
decompositions of $G$ and the relevant normalisers, together with their
block central characters, and multiplicativity of the induced central
function. These data determine the ambient block.
\lean{ModularRep.CharacterWeight.LocalBlockInductionSource} adds
compatibility with automorphisms and invariance under inner
automorphisms. Its maps
\lean{LocalBlockInductionSource.weightBlock} and
\lean{LocalBlockInductionSource.weightBlock_transport}
give the assignment on conjugacy classes and its equivariance.
The identification of the supplied operations with ordinary and
modular block theory is a hypothesis.
\end{implementation}

\begin{example}[Local maps from a global correspondence]
Let $\Omega:X\to\Alp(G)$ be a bijection. For each representative $Q$
of a radical conjugacy class, set
\[
 X_Q=\{x\in X:\Omega(x)\in\Alp(G)_{[Q]}\}.
\]
These subsets partition $X$. Restricting $\Omega$ and applying
Proposition~\ref{lib:weights:representative} gives a bijection
\[
 X_Q\longrightarrow\operatorname{dz}_\ell(N_G(Q)/Q).
\]
If $\Omega$ is equivariant, these local maps commute with transport
of subgroups and characters.
The construction for Brauer characters with a fixed compatible root
embedding is in
\module{ModularRep.PaperProofs.IntrinsicGlobalToRepresentativeMaps}.
\end{example}

\newpage
\section{Quotient by the \texorpdfstring{$\ell$}{ell}-core}

\begin{proposition}[Transport through the normal $\ell$-core]
Let $P=O_\ell(G)$. A weight $(Q,V)$ of $G$ in representation form determines
a weight in representation form of $G/P$ with radical subgroup $Q/P$ and
local representation transported along
\[
 N_G(Q)/Q\simeq N_{G/P}(Q/P)/(Q/P).
\]
Conversely, a weight in representation form of $G/P$ determines one of $G$
by taking the full inverse image of its radical subgroup and
transporting the local representation along the inverse isomorphism.
Both constructions preserve the dimension of the local representation.
\end{proposition}
\begin{proof}
By Proposition~\ref{lib:groups:core-contained}, every radical subgroup
$Q$ contains $P$. Proposition~\ref{lib:groups:radical-quotient}
identifies radicality before and after quotienting by $P$.
The normaliser quotient isomorphism then transports the local
representation. An isomorphism of groups preserves simplicity,
dimension and group order, and hence preserves defect zero.
The same reasoning applies to the full inverse image of a radical
subgroup of $G/P$.
\end{proof}
\begin{implementation}
The two constructions are
\lean{ModularRep.RepresentationWeight.quotientPCore} and
\lean{ModularRep.RepresentationWeight.liftFromQuotientPCore}
in \module{ModularRep.NormalCoreTransport}.
They construct individual weights. An inverse equivalence on
conjugacy classes is a separate assertion.
\end{implementation}

A comparison for weights in a specified block also requires
compatibility of the character conventions and block assignments on
the two local quotients. For applications involving modular character
triples, the required triple relation must hold for the transported
local character, with the specified roots and block operations.

\chapter{Basic sets and equivariant bijections}
\label{chap:library:assembly}

An integral basic set identifies two free abelian groups, one indexed
by ordinary characters and the other by Brauer characters. If a group
acts compatibly with decomposition, this identification is equivariant.
For a hypoelementary group, Conlon's theorem and injectivity of the mark homomorphism
give a bijection of the two sets of characters. Further results
in this chapter extend equivariant maps from stabilisers, cancel
invariant subsets and construct fixed representatives in product orbits.

Write $A$ for an acting group. An $A$-equivariant map satisfies
$f(ax)=af(x)$, and $A_x$ denotes the stabiliser of $x$. If a map $\pi:X\to I$
is given, its fibre over $i$ is the subset $X_i=\{x:\pi(x)=i\}$.

\section{Naturality and an integral basic set}

Let $G$ be finite and let $\ell$ be prime. Let $K$ and $k$ be fields
of characteristic zero and $\ell$, respectively. Consider the Grothendieck groups
$K_0(KG)$ and $K_0(kG)$ of Chapter~\ref{chap:library:kzero}, and a specified
decomposition homomorphism $d:K_0(KG)\to K_0(kG)$. Let $X$ and $Y$
be finite sets. For a set $X$, let
$\Z X$ denote the free abelian group with basis $(e_x)_{x\in X}$.
Suppose $X$ and $Y$ index distinct simple modules over $KG$ and $kG$,
respectively. Sending a basis vector to the corresponding module class
gives injective maps $j_X:\Z X\to K_0(KG)$ and
$j_Y:\Z Y\to K_0(kG)$, by the basis of simple classes in
Theorem~\ref{lib:kzero:basis-theorem}.

\begin{definition}
The family indexed by $X$ is an \emph{integral basic set} for the family
indexed by $Y$ if the restriction of $d$ is an isomorphism
$j_X(\Z X)\to j_Y(\Z Y)$. Equivalently, there is an isomorphism
$D:\Z X\to\Z Y$ such that $j_YD=dj_X$.
If $A$ acts on $X$ and $Y$,
it acts on the free groups by $a e_x=e_{ax}$ and $a e_y=e_{ay}$.
Compatible actions on the two Grothendieck groups are actions for which
$j_X$ and $j_Y$ commute with $A$.
\end{definition}

\begin{proposition}[Equivariance of the restricted decomposition map]
\label{lib:actions:basic-set}
Suppose these actions are compatible and $d$ commutes with $A$. Then $D$
is an isomorphism of integral permutation lattices. If
$D(e_x)=\sum_y d_{y,x}e_y$, its coefficients satisfy
$d_{ay,ax}=d_{y,x}$ for every $a\in A$.
\end{proposition}
\begin{proof}
For $v\in\Z X$, use the defining square and naturality to obtain
\[
 j_Y D(av)=d j_X(av)=d(a j_Xv)
          =a d j_Xv=a j_Y Dv=j_Y(a Dv).
\]
Injectivity of $j_Y$ gives $D(av)=aD(v)$. Taking coefficients on basis
vectors proves the stated matrix identity. Conversely, that identity on
basis vectors implies equivariance by linearity.
\end{proof}

The next section
gives conditions under which the two indexing sets are also
equivariantly bijective. For automorphism twisting, the construction in
Chapter~\ref{chap:library:reduction} gives naturality under its stated
compatibility hypothesis. For tensoring with linear characters, the ordinary and modular actions
must satisfy the corresponding character and root compatibility laws.
\begin{implementation}
The permutation representations are Mathlib's
\lean{Representation.ofMulAction}, with coefficients in $\Z$.
\module{ModularRep.DecompositionBasicSetBridge} defines
\lean{RestrictedIntegralBasicSet} and proves
\lean{matrixEquivariant_of_kZero_naturality}. Its
\lean{matrixEquivariant_of_stableReduction} obtains naturality from the
constructed decomposition map. The equivalence between the matrix and
linear map formulations is
in \module{ModularRep.IntegralBasicSetBridge}.
\end{implementation}

\begin{proposition}[Restriction to an invariant block]
\label{lib:actions:block-restriction}
Suppose $D$ is equivariant and $\pi_X:X\to I$, $\pi_Y:Y\to I$ are
$A$-invariant block assignments. If
\[
 D(\Z X_i)\subseteq\Z Y_i,
 \qquad D^{-1}(\Z Y_i)\subseteq\Z X_i
 \qquad(i\in I),
\]
then $D$ restricts to an isomorphism of $\Z A$-lattices
$\Z X_i\to\Z Y_i$ for every $i\in I$.
\end{proposition}
\begin{proof}
The support conditions restrict $D$ and its inverse to the two submodules
spanned by the indicated fibres. The restrictions are inverse because the
original maps are inverse. The action preserves each submodule, and the
restricted maps still commute with it. Identifying each supported submodule
with the free group on its fibre gives the assertion.
\end{proof}
\begin{implementation}
\module{ModularRep.BlockFibreRestriction}, especially
\lean{restrictBlockLinearEquiv}, \lean{restrictBlockLinearEquiv_equivariant}
and \lean{restrictRestrictedIntegralBasicSet}. Both support conditions
are hypotheses.
\end{implementation}

\section{From permutation lattices to finite sets}

Let $p$ be a prime, and write $\Z_p$ for the ring of $p$-adic
integers. The prime $p$ may differ from the characteristic $\ell$
used for modular characters. A finite group $A$ is
\emph{$p$-hypoelementary} if it has a normal $p$-subgroup with cyclic
quotient of order prime to $p$.
For an $A$-set $X$ and a subgroup $U\leq A$, the \emph{mark} at $U$
is $|X^U|$, where
\[
 X^U=\{x\in X:ux=x\text{ for every }u\in U\}.
\]

\begin{theorem}[Conlon]
Let $A$ be a finite group and let $X,Y$ be finite $A$-sets.
If $\Z_p X$ and $\Z_p Y$ are isomorphic as $\Z_p A$-lattices, then
\[
 |X^U|=|Y^U|
\]
for every $p$-hypoelementary subgroup $U$ of $A$.
\end{theorem}

Conlon's fixed coset calculation gives this implication
\cite[Section~4, especially (4.7)]{Conlon1968}.
Bouc gives a modern formulation, including the converse, in
\cite[Theorem~3.5.5]{Bouc2000}.
In that account the coefficient ring is a complete commutative
Noetherian local ring with residue characteristic $p$, and lattices
are finitely generated and free over that ring. The choice $\Z_p$
satisfies these hypotheses.

\begin{theorem}[Burnside's determination by marks]
Let $A$ be a finite group. Two finite $A$-sets $X$ and $Y$ are
isomorphic if and only if $|X^U|=|Y^U|$ for every subgroup $U\leq A$.
\end{theorem}
\begin{proof}
An equivariant bijection identifies the fixed points of every subgroup.
For the converse, decompose each set into transitive orbits $A/H$,
with $H$ taken from a set of representatives of subgroup conjugacy
classes. The matrix
\[
 m(U,H)=|(A/H)^U|
\]
is triangular when those representatives are ordered by subgroup
order. Its diagonal entries are $|N_A(H):H|$, which are nonzero.
Equality of marks therefore gives equality of the multiplicities of
all transitive orbits.
\end{proof}
For this statement and proof, see, for example,
\cite[Theorem~2.3.2]{Bouc2000}.
The equivalent formulation is injectivity of the mark homomorphism
from the Burnside ring.

\begin{corollary}[Permutation lattices]
\label{lib:actions:conlon}
Let $A$ be a finite $p$-hypoelementary group and let $X,Y$ be finite
$A$-sets. If $\Z_p X\simeq\Z_p Y$ as $\Z_p A$-lattices, then
$X$ and $Y$ are equivariantly bijective. In particular, the same
conclusion holds if $\Z X\simeq\Z Y$ as $\Z A$-lattices.
\end{corollary}
\begin{proof}
Conlon's theorem gives equality of the marks at all
$p$-hypoelementary subgroups. Every subgroup of $A$ is
$p$-hypoelementary by Proposition~\ref{lib:groups:hypoelementary}.
Thus all marks agree, and the preceding theorem gives the bijection.
For the integral version, extend the isomorphism and its inverse
from $\Z$ to $\Z_p$.
\end{proof}
\begin{implementation}
\module{ModularRep.PaperProofs.ConlonBasicSet} constructs scalar
extension and the mark map on isomorphism classes of finite $A$-sets.
The scalar extension uses Mathlib's equivalences for tensor products
with finitely supported functions. The intertwining identities are
proved in this module.
The two published theorems are separate hypotheses:
\lean{PadicConlonMarkDetection} and
\lean{PublishedBurnsideMarkInjectivity}. Both range over all finite
$A$-sets. Under these hypotheses,
\lean{corollary_2_4_conlonMark} proves the implication from
isomorphic $p$-adic permutation lattices to an equivariant bijection.
The proofs of the published theorems are not part of that declaration.
\end{implementation}

If $A$ is $p$-hypoelementary, Proposition~\ref{lib:actions:basic-set}
and Corollary~\ref{lib:actions:conlon} give an equivariant bijection
between the ordinary basic set and the modular characters it indexes.
Under the same hypothesis, Proposition~\ref{lib:actions:block-restriction}
gives the corresponding bijection on an invariant block.
\begin{implementation}
The corresponding declarations are
\lean{equivariantSetEquiv_of_kZero_naturality} and
\lean{equivariantSetEquiv_restrictBlock_of_kZero_naturality}
in the modules for basic sets and block restriction.
They assume the specified integral basic set, actions, naturality,
hypoelementary group and published mark theorems.
They prove existence of a bijection without prescribing its pairs.
No unitriangularity hypothesis is used.
\end{implementation}

\section{Extending a bijection over an orbit}

Let $\pi_X:X\to I$ and
$\pi_Y:Y\to I$ be equivariant maps, and assume $I$ is transitive.

\begin{proposition}[Extension from one fibre]
\label{lib:actions:fibre-transport}
Fix $i_0\in I$ and an $A_{i_0}$-equivariant bijection
$\beta_0:X_{i_0}\to Y_{i_0}$. There is a unique $A$-equivariant bijection
$\beta:X\to Y$ that preserves the maps to $I$ and restricts to $\beta_0$ on
$X_{i_0}$.
\end{proposition}
\begin{proof}
For each $i\in I$, choose $t_i\in A$ with $t_i i_0=i$, and define
\[
 \beta(x)=t_i \beta_0(t_i^{-1}x)\qquad(x\in X_i).
\]
The expression has image in $Y_i$. Replacing $\beta_0$ by its inverse gives
an inverse map on every fibre and hence on the total sets. At $i=i_0$,
the chosen $t_{i_0}$ belongs to $A_{i_0}$, so equivariance of $\beta_0$
shows that the restriction is exactly $\beta_0$.

Two choices of $t_i$ differ by an element of $A_{i_0}$. Equivariance
of $\beta_0$ under that stabiliser makes the formulas identical. For
$a\in A$, the element $t_{ai}^{-1}at_i$ also fixes $i_0$. Applying
$A_{i_0}$-equivariance to this element proves $\beta(ax)=a\beta(x)$.
Finally, an equivariant extension of $\beta_0$ must send
$t_i(t_i^{-1}x)$ to $t_i \beta_0(t_i^{-1}x)$, which proves uniqueness.
\end{proof}
\begin{implementation}
\module{Formalisation.FibreTransport}, with \lean{Formalisation.transportedEquiv},
\lean{Formalisation.transportedEquiv_equivariant},
\lean{Formalisation.transportedEquiv_independent_of_transporter} and
\lean{Formalisation.transportedEquiv_preserves_base}.
The equivalence of the total sets uses Mathlib's \lean{Equiv.ofFiberEquiv}.
The transporter formula, its equivariance and its independence of
the choices are proved in this module.
The uniqueness assertion follows from the formula for $\beta$, as in the
proof above. It is not a separate theorem in this module.
\end{implementation}

For an arbitrary $A$-set $I$, choose one representative $i$ of each
orbit and an $A_i$-equivariant bijection $X_i\to Y_i$. Applying the
proposition on each orbit gives an $A$-equivariant bijection $X\to Y$
preserving the maps to $I$. The module
\module{Formalisation.IndexedAssembly} combines families satisfying
$\beta_{ai}(ax)=a\beta_i(x)$. The transporter formula above proves precisely
this compatibility from the bijections on the representative fibres.

The same formula preserves any specified $A$-invariant relation
$\mathcal R\subseteq X\times Y$, provided the initial bijections
satisfy it. More explicitly, assume
\begin{align*}
 (x,y)\in\mathcal R&\ \Longrightarrow\ (ax,ay)\in\mathcal R
   &&(a\in A),\\
 (x,\beta_i(x))&\in\mathcal R &&(x\in X_i)
\end{align*}
with the second condition imposed on each representative fibre.
Every $x\in X$ has the form $au$ for
some $u$ in one of those fibres, and
$(x,\beta(x))=(au,a\beta_i(u))\in\mathcal R$. Thus the local relation and
its invariance are separate hypotheses in an application.

\begin{example}[A block correspondence on an orbit]
Suppose $I$ is an orbit of blocks, $X_i$ is its Brauer character set and
$Y_i$ its set of weight classes. An equivariant bijection for one block,
under that block's stabiliser, extends to the whole orbit by
Proposition~\ref{lib:actions:fibre-transport}. It still sends each $X_i$
into $Y_i$. A modular character triple relation is preserved by the
extension provided it holds for each pair on the chosen fibre and is
invariant under the automorphisms being used.
\end{example}

\section{Involutions and cancellation}

\begin{proposition}[Classification of finite sets with an involution]
\label{lib:actions:involutions}
Let $\sigma$ and $\tau$ be involutions on finite sets $X$ and $Y$.
There is a bijection $\beta:X\to Y$ with $\beta\sigma=\tau \beta$ if and only if
$|X|=|Y|$ and $|X^\sigma|=|Y^\tau|$.
\end{proposition}
\begin{proof}
Each orbit of an involution has size one or two. If $f$ is the number
of fixed points and $t$ the number of orbits of size two, then
$|X|=f+2t$. Thus the two stated equalities imply that the numbers of
orbits of each size agree. Match the singleton orbits and the orbits of
size two. A choice of one point in each paired orbit of size two determines
the bijection on its other point by equivariance. Combining these maps gives
$\beta$. Conversely, any such $\beta$ preserves total size and carries the fixed
points bijectively to the fixed points.
\end{proof}

Equivalently, the two counts determine the cycle type of an involution
and hence its conjugacy class in the symmetric group.
\begin{implementation}
\module{Formalisation.C2Cancellation}, including
\lean{card_eq_fixed_add_twice_cycleCount} and
\lean{exists_equivariantEquiv_of_card_eq_of_fixed_card_eq} in
\lean{Formalisation.C2Cancellation}.
These results apply Mathlib's fixed point counts and classification
of permutations by cycle type,
\lean{Equiv.Perm.isConj_iff_cycleType_eq}, to involutions.
\end{implementation}

The pair $(|X|,|X^\sigma|)$ determines a finite set with involution up
to equivariant bijection. Both entries add under invariant disjoint
unions. Consequently, if $X\simeq Y$ equivariantly and invariant subsets
$X'\subseteq X$, $Y'\subseteq Y$ are equivariantly bijective, then
$X\setminus X'\simeq Y\setminus Y'$ equivariantly. The argument
constructs a bijection of the complements from their two counts.

For two parts exchanged by the involution, a different formula is useful.
Choose a bijection $f:X_0\to Y_0$, and define the map on $X_1$ by
$\tau f\sigma$. It takes $X_1$ to $Y_1$, is inverse to the analogous
map constructed from $f^{-1}$ and commutes with the involutions.
These constructions are in \module{Formalisation.BlockCancellation}
and \module{Formalisation.C2Cancellation}.

\begin{proposition}[Cancellation over a block orbit]
\label{lib:actions:block-cancellation}
Let $I$ be a finite set with an involution, and let $X\to I$ and
$Y\to I$ be equivariant maps of finite sets with involutions.
Write $I_0,\ldots,I_r$ for the orbits in $I$, and let $U_j$ and
$V_j$ be their inverse images in $X$ and $Y$. Suppose that
$X\simeq Y$ equivariantly and that $U_j\simeq V_j$ equivariantly
for $1\leq j\leq r$. Then there is an equivariant bijection
$U_0\to V_0$ preserving the map to $I_0$.
\end{proposition}
\begin{proof}
Cancel the known parts in the total size and fixed point count.
Proposition~\ref{lib:actions:involutions} gives an equivariant
bijection between the remaining parts. If $I_0$ is a singleton,
it already preserves the map to $I_0$.

Otherwise write $I_0=\{i,\sigma i\}$. The involutions exchange
the two fibres on each side, so each fibre has half the size of
its corresponding remaining part. These sizes agree. Choose a
bijection $f:X_i\to Y_i$ and define the map on the other fibre by
$\tau f\sigma$, where $\sigma$ and $\tau$ denote the source and
target involutions. The two formulas give an equivariant bijection
preserving both fibres.
\end{proof}

For an application to blocks, the maps to $I$ assign a block to each
character or weight class. The hypotheses then require equivariant
comparisons on each of the other block orbits.

\begin{implementation}
In \module{Formalisation.BlockCancellation},
\lean{cancel_equivariant_equiv_of_parts} handles the finite
disjoint union of known parts, and
\lean{cancel_exchanged_blocks_preserving} constructs the map on
the remaining pair of exchanged fibres. The resulting map need not
be the restriction of the initially given equivalence $X\simeq Y$.
\end{implementation}

\begin{example}
A set with eight elements and two fixed points has three orbits of size
two. If an invariant part consists of one such orbit, the complement has
the pair of counts $(6,2)$. Any other complement with these counts is
isomorphic as a set with involution.
\end{example}

\section{Fixed representatives in a product}

Let $D\rtimes R$ act on a set $X$. Suppose that
\[
 (D\rtimes R)_x=D_x\rtimes R_x
\]
for some $x\in X$. If $r\in R$ satisfies $rx\in Dx$, then $rx=x$.
Indeed, write $rx=dx$. The element $(d^{-1},r)$ fixes $x$, so the
stabiliser equality implies $r\in R_x$.

\begin{proposition}[A fixed tuple in a product orbit]
Let $n>0$, let $O_0,\ldots,O_{n-1}$ be $D$-orbits in $X$, and let
$T:X\to X$. Suppose a cyclic group $E=\langle\tau\rangle$ acts on
$X^n$ by
\[
 \tau(y_0,\ldots,y_{n-1})=(Ty_{n-1},Ty_0,\ldots,Ty_{n-2}).
\]
Assume that $\tau$ preserves $\prod_jO_j$. Choose $x\in O_0$ such
that
\[
 T^n x=rx\quad\text{for some }r\in R,
 \qquad (D\rtimes R)_x=D_x\rtimes R_x.
\]
Then $(x,Tx,\ldots,T^{n-1}x)$ belongs to $\prod_jO_j$ and is fixed
by $E$.
\end{proposition}
\begin{proof}
Varying the $j$th coordinate of a tuple in the product shows that
$T(O_j)\subseteq O_{j+1}$, with indices taken modulo $n$.
Consequently $T^n x\in O_0=Dx$. The preceding observation gives
$rx=x$, and hence $T^n x=x$.
Successive transport places $T^jx$ in $O_j$.
The formula defining $\tau$ now shows that it fixes the tuple,
so its cyclic group $E$ fixes the tuple as well.
\end{proof}

Apply this construction on each cycle when there are several cycles
of coordinates. The groups $D$ and $R$ may vary from cycle to cycle.
If the actions extend compatibly to $D_{\mathrm{prod}}\rtimes E$,
where $D_{\mathrm{prod}}$ is the product of the coordinate groups,
the resulting tuple $\theta$ satisfies
\[
 (D_{\mathrm{prod}}\rtimes E)_\theta
   =(D_{\mathrm{prod}})_\theta\rtimes E.
\]
Indeed, every element of $E$ fixes $\theta$, and $(d,e)$ fixes
$\theta$ exactly when $d$ does.

\begin{implementation}
\module{ModularRep.ComponentReturnFull} allows separate groups and
orbit data on every cycle.
The theorem \lean{component_return_full_of_product_orbit_stable} in
\lean{ModularRep.ManuscriptVerification.ComponentReturnFull}
derives the conditions on coordinate orbits from stability of the
product orbit and constructs the fixed tuple.
In a representation theoretic application, the product character
sets, their diagonal actions and the action after a full cycle must
satisfy the stated action formulas.
\end{implementation}

\chapter{Character counts from finite matrices}
\label{chap:library-computations}

Restriction of ordinary characters to regular elements gives a
matrix whose rank is the number of irreducible Brauer characters.
If an involution fixes the block, a second rank calculation gives
the number of Brauer characters it fixes. Both calculations follow
from linear algebra once the matrix entries and the action on
conjugacy classes are identified.

\section{The rank of ordinary restrictions}

\begin{lemma}
Let $M,N,V$ be vector spaces over a field $K$, with $N$ finite
dimensional. Let $D:M\to N$ and $E:N\to V$ be linear maps.
If $D$ is surjective and $E$ is injective, then
\[
 \dim_K\im(ED)=\dim_K N.
\]
\end{lemma}
\begin{proof}
Surjectivity gives $\im(ED)=\im E$. Injectivity identifies $\im E$
with $N$.
\end{proof}

We recall the character theoretic application over the complex
numbers, with a fixed choice of roots for Brauer characters.
For an $\ell$-block $b$, write $\Irr(b)$ and $\IBr(b)$ for its ordinary
irreducible characters and irreducible Brauer characters.

\begin{theorem}[Rank of ordinary restrictions]
Let $b$ be an $\ell$-block of a finite group $G$. If $\chi^0$ denotes
restriction to the $\ell$-regular elements, then
\begin{equation}
 \dim_{\mathbb C}\langle\chi^0:\chi\in\Irr(b)\rangle_{\mathbb C}
   =|\IBr(b)|.
 \label{lib:computations:restriction-rank}
\end{equation}
\end{theorem}
\begin{proof}
The irreducible Brauer characters are linearly independent.
The decomposition formula is
\[
 \chi^0=\sum_{\varphi\in\IBr(b)}d_{\chi\varphi}\varphi
 \qquad(\chi\in\Irr(b)).
\]
The decomposition matrix of $b$ has rank $|\IBr(b)|$.
Thus its rows span the coordinate space on $\IBr(b)$.
Apply the lemma with $D$ sending the ordinary character basis to
these rows and $E$ sending the coordinate basis to the Brauer
characters.
\end{proof}

For linear independence, the decomposition formula and full rank
of the decomposition matrix, see, for example,
\cite[Theorem~2.6 and Corollaries~2.9--2.11]{Navarro1998}.
The passage to one block follows from the block diagonal form of
the decomposition matrix
\cite[Theorem~3.3 and the following discussion]{Navarro1998}.
The related equality between $|\IBr(G)|$ and the number of regular
conjugacy classes is Corollary~2.10 in that account.

The same rank argument works over a field $K$ of characteristic
zero if the characters take values in $K$, the decomposition formula
holds there, the Brauer characters are linearly independent over
$K$, and the decomposition rows span the coordinate space over $K$.
Theorem~\ref{lib:brauer:independence} gives independence under its
stated field and root hypotheses. The decomposition and spanning
conditions are additional hypotheses.

\begin{implementation}
\module{Formalisation.RestrictionRank} proves
\lean{finrank_range_comp_eq_card} and
\lean{finrank_span_matrix_combinations_eq_card}.
The range calculation applies Mathlib's
\lean{LinearMap.range_comp_of_range_eq_top} and
\lean{LinearMap.finrank_range_of_inj}, followed by its dimension
formula for a finite coordinate space.
The second theorem assumes independent Brauer vectors and
decomposition rows spanning their coordinate space.
It proves the linear algebra implication over $K$.
The identification of these vectors and coefficients with characters
and decomposition numbers is separate.
\end{implementation}

Let $c_1,\ldots,c_t$ represent all $\ell$-regular conjugacy classes.
Evaluation at these representatives is injective on class functions
defined on the $\ell$-regular elements.
Consequently the matrix
\[
 M=(\chi(c_j))_{\chi\in\Irr(b),\,1\leq j\leq t}
\]
has row rank $|\IBr(b)|$. More generally, one may use any list
meeting every regular conjugacy class. To apply the formula to a
specified matrix, its rows must be the indicated ordinary characters
and its columns must evaluate those characters at such a list.

\section{Fixed points from symmetrised rows}

\begin{proposition}[An involution permuting a basis]
\label{lib:computations:fixed-rank}
Let $K$ have characteristic zero. Let $(v_i)_{i\in I}$ be a finite
basis of a $K$-vector space $V$, and let $T:V\to V$ be a linear map
such that $T(v_i)=v_{\sigma(i)}$ for an involution $\sigma$ of $I$.
Set
\[
 r=\dim_K V,\qquad s=\dim_K\im(1+T).
\]
Then $|I^\sigma|=2s-r$.
\end{proposition}
\begin{proof}
The permutation matrix of $T$ has diagonal entry one at a fixed
index and zero otherwise. Hence $\tr(T)=|I^\sigma|$ in $K$.
Since $T^2=1$, the map $P=(1+T)/2$ is a projection onto $\im(1+T)$.
The trace of a projection is its rank, so
\[
 s=\tr(P)=\tfrac12(r+\tr(T)).
\]
Characteristic zero gives the asserted equality of integers.
\end{proof}

The basis need not be known explicitly. Suppose $(g_a)_{a\in A}$
spans $V$ and $E:V\to K^C$ is an injective linear map.
Let $M$ be a matrix and let $\pi$ permute $C$, with
\[
 E(g_a)(c)=M_{a,c},\qquad
 E(Tg_a)(c)=M_{a,\pi(c)}.
\]
The vectors $(1+T)g_a$ span $\im(1+T)$. Their images under $E$ are
the rows of the \emph{symmetrised matrix}
\[
 M^+_{a,c}=M_{a,c}+M_{a,\pi(c)}.
\]
Thus $M$ has row rank $r$ and $M^+$ has row rank $s$.

For a block fixed by an involution in $\Aut(G)$, take the basis
to be its irreducible Brauer characters and the spanning vectors
to be its ordinary restrictions. Once $\pi$ describes the action
on the chosen regular class representatives, the two ranks give
\[
 \bigl(|\IBr(b)|,\ |\IBr(b)^\sigma|\bigr)=(r,2s-r).
\]
The action formula must use the same characters, roots and class
representatives as the evaluation formula.

\begin{implementation}
The trace and projection argument is
\lean{fixed_card_add_card_eq_two_mul_plus_rank} in
\module{ModularRep.PaperProofs.SporadicFi24P3Definition44NamedCarrierInvolutionBasisRank}.
It applies to any finite basis over a field of characteristic zero.
Its proof uses Mathlib's \lean{Matrix.trace_permutation} and
\lean{LinearMap.IsProj.trace}. The local argument identifies the
basis action with a permutation matrix and relates the projection
to the image of $1+T$.
The comparison with the symmetrised row rank is
\lean{finrank_range_id_add_eq_symmetrizedRows} in
\module{ModularRep.PaperProofs.SporadicFi24P3Definition44NamedCarrierSymmetrizedRowRankEquality}.
Its hypotheses are the spanning, injectivity and evaluation
identities above. The matrix identity in this latter theorem
allows any map on the column index set.
\end{implementation}

\section{Certifying a matrix rank}

\begin{proposition}[A minor and spanning rows]
Let $M$ be a finite matrix over a field $K$. Suppose a square minor
on $r$ selected rows and $r$ selected columns has nonzero determinant.
If every row of $M$ is a linear combination of the selected rows,
then $M$ has row rank $r$.
\end{proposition}
\begin{proof}
Restricting the selected rows to the selected columns gives the
rows of an invertible matrix, so they are linearly independent.
By hypothesis they span the row space of $M$. They therefore form
a basis of that space.
\end{proof}

\begin{implementation}
\module{ModularRep.PaperProofs.FiniteRowRankCertificate}
records the selected rows and columns, the determinant of the
minor, its nonvanishing and the coefficients expressing every row
as a combination of the selected rows.
\lean{FiniteRowRankCertificate.rows_finrank} proves the rank from
the determinant identity and the equations expressing every row
in terms of the selected rows.
It uses Mathlib's \lean{Matrix.linearIndependent_rows_of_det_ne_zero}
for the selected minor and \lean{finrank_span_eq_card} for the
dimension of the span. The local proof shows that the selected rows
span the full row space.
\end{implementation}

Together with the evaluation and action identities, these finite
equations give character counts and counts of fixed characters.
To use the rank calculation for character counts, one must identify the full
matrix with the restrictions of all the intended characters to all the relevant
regular classes. The selected minor proves that the selected rows are linearly
independent, and the given row expressions prove that they span the full row space.

\section{Finite scalar products and character selection}

Let $K$ be a field of characteristic zero, let $H$ be a finite group,
and let $(\mathcal C_j)_{j\in J}$ be distinct
conjugacy classes with representatives $u_j$, where $J$ is finite.
If a function $f:H\to K$ is constant on each $\mathcal C_j$ and vanishes
outside their union, partitioning the finite sum gives
\[
 \frac{1}{|H|}\sum_{g\in H}f(g)
   =\sum_{j\in J}\frac{|\mathcal C_j|}{|H|}f(u_j).
\]
Every coefficient $w_j=|\mathcal C_j|/|H|$ is nonzero. In a character calculation,
the required constancy and vanishing concern the product of the two
functions in the scalar product.

\begin{implementation}
\module{ManuscriptIBAW.Library.FiniteClassPairing} proves
\lean{FiniteClassPairing.sum_localize} for a finite set with an explicit
class map and distinct selected class labels. It uses constancy on
the selected fibres and vanishing elsewhere. The theorem
\lean{FiniteClassPairing.weight_ne_zero} proves nonvanishing of the
normalised fibre sizes in characteristic zero, using a representative
in each fibre. These results do not require a character theory.
\end{implementation}

Let $K$ be a field, let $J$ be finite, and let the vectors
$v_i\in K^J$, indexed by a set $I$, span $K^J$. Suppose that
$w_j\ne0$ for every $j$ and that $z_r\in K^J$ is nonzero for each
$r\in J$. Let $a_{i,r}\in K$ and $\varepsilon_i\in K$ satisfy
\[
 \sum_{j\in J}w_jv_i(j)z_r(j)=\varepsilon_i a_{i,r}
 \qquad(i\in I,\ r\in J).
\]
Then no column of $(a_{i,r})$ is zero. Indeed, if the column at $r$
were zero, the linear functional
$v\mapsto\sum_jw_jv(j)z_r(j)$ would vanish on a spanning family.
It would therefore vanish on every coordinate vector, giving
$w_jz_r(j)=0$ for all $j$. This contradicts $z_r\ne0$.

Suppose now that $K$ has characteristic zero, every $a_{i,r}$ is a
natural number viewed in $K$, and each row has sum one. Every row
then has exactly one entry equal to one and all others zero. Since
every column is nonzero, choose $s(r)\in I$ with $a_{s(r),r}=1$.
The map $s:J\to I$ is injective and satisfies
\[
 a_{s(r),j}=\begin{cases}1&j=r,\\0&j\ne r.\end{cases}
\]
Injectivity follows by comparing the unique nonzero entries in two
selected rows.

\begin{implementation}
In \module{ManuscriptIBAW.Library.FiniteCharacterSelection},
\lean{FiniteCharacterSelection.coverage_of_span} proves that every
column is nonzero. The theorem
\lean{FiniteCharacterSelection.exists_delta_injection_of_span}
constructs the injective selection assuming that the entries are
nonnegative integers and each row sums to one. These are finite linear algebra and arithmetic
results over the stated fields.
\lean{span_eq_top_of_diagonal_pairings} and
\lean{span_eq_top_of_trace_expansions} derive the spanning condition
from diagonal pairing identities with nonzero diagonal entries and from
explicit finite trace expansions, respectively.
\lean{linearIndependent_of_spanning_independent_family} proves that a finite family is linearly independent if its span contains
a linearly independent family of the same size.
In the character application, the identification of these families,
expansions and scalar products with the geometric and representation
theoretic results remains a separate set of assumptions.
\end{implementation}

The row sum condition can also be proved by counting. Let $I$ and $J$
be finite sets with $I\ne\varnothing$, and let $U$ be a set.
Let $f:J\to U$, $u\in U$, and $a:I\times J\to\mathbb N$. Suppose that the sum in column $j$ is
one when $f(j)=u$ and zero otherwise, and that $|I|=|f^{-1}(u)|$.
If every two rows are related by a permutation of $J$, then every row
has sum one. Indeed, all row sums have a common value $r$. Summing
the entries by rows and by columns gives
\[
 |I|r=|f^{-1}(u)|=|I|,
\]
and $I\ne\varnothing$ implies $r=1$.

\begin{implementation}
\lean{ManuscriptIBAW.TypeC.PrincipalGGGR.row_sum_one} proves this
finite arithmetic statement. It takes the column sums, the equality
of cardinalities and the row permutations as assumptions.
\end{implementation}

\chapter{Group actions, characters and block parameters}
\label{chap:library-current}

Equivariant maps and character identities often determine an action without
requiring an explicit choice of representatives. We study this principle for
permutation modules, block stabilisers and ordinary characters over different
fields. We also describe the quotient associated with a similitude multiplier
and the relation between block idempotents and rational series.

\section{Fixed bases and permutation modules}

Let $A$ act linearly on a module $V$ over a nonzero ring $R$. If $A$ fixes every
vector of an $R$-basis, it acts trivially on $V$: a linear map is determined
by its values on a basis. If another basis is indexed by an $A$-set $Y$
and the action permutes its basis vectors in accordance with the action on
$Y$, the linear independence of that basis implies that $A$ fixes $Y$
pointwise.

\begin{proposition}
Let $X$ and $Y$ be finite $A$-sets. Suppose that
$d\colon\Z X\longrightarrow\Z Y$ is an $A$-equivariant isomorphism of
free abelian groups. If $A$ fixes $X$ pointwise, it fixes $Y$ pointwise,
and there is an $A$-equivariant bijection $X\longrightarrow Y$.
\end{proposition}

\begin{proof}
The images under $d$ of the basis vectors indexed by $X$ form a pointwise
fixed basis of $\Z Y$. The preceding argument shows that $A$ acts trivially
on $\Z Y$, and hence on $Y$. Equality of the ranks gives $|X|=|Y|$.
Any bijection between the two finite sets is then equivariant.
\end{proof}

Let $X$ be an $A$-stable integral basic set of a block $b$, and suppose
that the restriction of the decomposition homomorphism gives an
$A$-equivariant isomorphism $\Z X\longrightarrow\Z\IBr(b)$.
If $A$ fixes every character in $X$, the proposition shows that it
fixes every irreducible Brauer character of $b$ and gives an
$A$-equivariant bijection $X\longrightarrow\IBr(b)$.

\begin{implementation}
\module{ManuscriptIBAW.Library.FixedBasis} specialises Mathlib's basis
extensionality, injectivity and equivalence of basis index sets to the
stated actions. \module{ManuscriptIBAW.Library.FixedBasicSet} applies
this argument to Mathlib's integral permutation modules. The block application is
\module{ManuscriptIBAW.TypeC.EvenBasicSet}.
\end{implementation}

\section{Equivariant bijections for involutions}

\begin{proposition}
Let $A$ be a group acting on finite sets $X$ and $Y$. Suppose that $A$ is generated by a
set $T$ and an element $\delta$, that every element of $T$ fixes both sets
pointwise, and that $\delta$ induces an involution on each set. There is
an $A$-equivariant bijection $X\longrightarrow Y$ if
\[
 |X|=|Y|\qquad\text{and}\qquad |X^{\langle\delta\rangle}|
   =|Y^{\langle\delta\rangle}|.
\]
\end{proposition}

\begin{proof}
Each orbit under $\delta$ has size one or two. The two equalities give
equal numbers of orbits of each size. Match the fixed points and the
orbits of size two to obtain a bijection commuting with $\delta$.
It commutes with $T$ because $T$ acts trivially. The elements of $A$
with which the bijection commutes form a subgroup, so the generating
hypothesis gives equivariance under $A$.
\end{proof}

The hypothesis concerns the permutations induced by $\delta$. It does
not require that $\delta$ itself have order two in $A$. This distinction
allows the proposition to be used when $\delta^2$ is an inner automorphism
and the objects under consideration are conjugacy classes.

\begin{implementation}
\module{ManuscriptIBAW.Library.EquivariantBijections} proves the subgroup
and generating set arguments and uses
\lean{Formalisation.C2Cancellation} for the two involutions. That result
specialises Mathlib's conjugacy classification of finite permutations by cycle type
and its fixed point count. The local statements express the conclusion as
an equivariant bijection for the given actions.
\end{implementation}

\section{Independent signs on a direct sum}

Let $V=\bigoplus_{c\in I}V_c$ be a finite dimensional space over a field of
characteristic different from two. For each $c$, let $\varepsilon_c$ act
as $-1$ on $V_c$ and as $1$ on the other summands. A linear map commuting
with every $\varepsilon_c$ preserves every $V_c$. Indeed, in these
coordinates, commutation forces an entry between distinct summands to
equal its negative. That entry is zero because two is invertible.

If the decomposition is orthogonal for a nondegenerate alternating form
and the map is symplectic, its restriction to each summand is symplectic.
The form equation restricts to each diagonal block. Consequently the map
belongs to the product of the symplectic groups of the summands.

Suppose that $R=\prod_c R_c$ acts independently on these summands, that
$-1\in R_c$, and that the centraliser of $R_c$ in its symplectic group
is contained in $R_c$. Every element centralising $R$ commutes with the
independent signs, hence preserves each summand. Its restriction there
centralises $R_c$ and therefore belongs to $R_c$. Thus $C_{\Sp(V)}(R)\leq R$.
The index $c$ distinguishes different summands, including copies on which
the same group acts.

\begin{implementation}
\lean{ManuscriptIBAW.TypeC.independentSignSeparation} in
\module{ManuscriptIBAW.TypeC.PrincipalProducts} proves the matrix argument
for the specified symplectic coordinates over a finite field of odd order.
Mathlib supplies matrix reindexing and the block diagonal multiplication
and transpose identities. The proof combines them with the specified
embedding and Gram matrix equations to separate the summands.
\lean{ManuscriptIBAW.TypeC.PrincipalProductData.signs} is a proved consequence
of those equations. The local centraliser assertion remains a source
assumption in the application.
\end{implementation}

\section{Stabilisers under an equivariant bijection}

Let $A=M\rtimes E$ act on sets $X$ and $Y$, and let
$\beta:X\longrightarrow Y$ be an $A$-equivariant bijection. Then
$A_x=A_{\beta(x)}$, since equivariance and injectivity give
$\beta(ax)=\beta(x)$ if and only if $ax=x$. In particular, if
\[
 m(e\beta(x))=\beta(x)
 \quad\Longleftrightarrow\quad
 m\beta(x)=\beta(x)\ \text{and}\ e\beta(x)=\beta(x),
\]
the same equivalence holds with $x$ in place of $\beta(x)$. This transfers
the factorisation of a stabiliser into its $M$ and $E$ parts.

Chapter~\ref{chap:library:assembly} constructs a global bijection from
bijections over representatives of block orbits that are equivariant under
their stabilisers.
It preserves the block assignments. For character sets of a finite group
$G$, suppose that $A$ acts through the usual automorphism action and maps
onto $\Aut(G)$. Equivariance under $A$ then gives equivariance under every
automorphism: choose a preimage and apply the same identity.

\begin{implementation}
\module{ManuscriptIBAW.TypeC.OddGlobalFactorization} constructs
\lean{globalBijection} from the block basic sets, the stated quotient
actions and the Conlon and Burnside hypotheses. In that namespace,
\lean{GlobalBijection.automorphism_equivariant} proves full automorphism
equivariance, and \lean{GlobalBijection.allBrauerFactorization} transfers
ordinary stabiliser factorisation through that same bijection. The
construction uses the basic set hypotheses first, then transfers the
factorisation to the Brauer characters.
\end{implementation}

\section{Linear characters and block stabilisers}

Under the coefficient and tensor compatibility hypotheses of
Section~\ref{lib:blocks:linear-stabilisers}, let $G$ be finite,
let $\ell$ be prime and let $N\lhd G$.
Let $b$ be an $\ell$-block of $G$, and let $C$ be a
subgroup of the group of ordinary linear characters of $G/N$ of order
prime to $\ell$. Inflate these characters to $G$ and let them act by
tensoring. Each character in $C_b$ is trivial on $Z(G)$, so it factors
through $G/NZ(G)$. In particular,
\[
 |C_b|\leq |G:NZ(G)|.
\]
The same assertion holds for the corresponding linear Brauer characters.
The proof uses the common central character of the irreducible Brauer
characters of $b$. A character stabilising $b$ is trivial on the
$\ell'$-part of $Z(G)$ by this central character, and on its $\ell$-part
because its order is prime to $\ell$.

Now let a finite cyclic group $E$ act on $C$, and suppose that $C\rtimes E$
acts on the blocks in a way whose restriction to $C$ is the tensor
action just described. For $J=(C\rtimes E)_b$, let $D$ be the kernel
of the projection $J\longrightarrow E$. Sending an element of $D$ to
its $C$ coordinate gives an injection $D\longrightarrow C_b$.
Thus $|D|\leq |G:NZ(G)|$. If this index is at most two, $D$ is a normal
$2$-subgroup of $J$ and $J/D$ is cyclic.

Let $P/D$ be the Sylow $2$-subgroup of $J/D$. Then $P\lhd J$ is a
$2$-group and $J/P$ is cyclic of odd order. For every subgroup $U\leq J$,
the subgroup $U\cap P$ is a normal $2$-subgroup and
$U/(U\cap P)$ embeds in $J/P$. Consequently $U/O_2(U)$ is cyclic of
odd order. Thus every subgroup of $J$ is $2$-hypoelementary, as required
in the application of Conlon's theorem.

\begin{implementation}
The central character argument is in
\module{ModularRep.LinearBlockStabilizer}.
\module{ManuscriptIBAW.Library.TensorStabilizer} identifies the kernel of the
projection to $E$ with its character coordinate. It combines the resulting
bound with the hypoelementary group calculation in ModularRep.
Compatibility of the block action with tensoring
simple modules remains an explicit hypothesis.
\end{implementation}

\newpage
\begin{samepage}
\section{The multiplier quotient}

\begin{proposition}
Let $\mu\colon G\longrightarrow U$ be a surjective homomorphism to an
abelian group. It induces an isomorphism
\[
 G/(\ker\mu\,Z(G))\ \cong\ U/\mu(Z(G)).
\]
If $U=\F_q^\times$ and $\mu(Z(G))=(\F_q^\times)^2$, this quotient has
order at most two, and has order two when $q$ is odd.
\end{proposition}
\end{samepage}

\begin{proof}
The composite $G\longrightarrow U\longrightarrow U/\mu(Z(G))$ is
surjective with kernel $\ker\mu\,Z(G)$, so the first assertion is the
first isomorphism theorem. In $\F_q^\times$, the kernel of the squaring
map is the group of roots of $x^2=1$. It has at most two elements and
has exactly two in odd characteristic. The index of the image of an
endomorphism of a finite group equals the order of its kernel.
\end{proof}

For the conformal symplectic group of positive rank, the multiplier is specified by
$\langle gv,gw\rangle=\mu(g)\langle v,w\rangle$. Its kernel is the
symplectic subgroup. Surjectivity and the image of the centre are the
structural hypotheses in this application. The proposition determines the quotient index
from these hypotheses.

\begin{implementation}
\module{ManuscriptIBAW.Library.MultiplierQuotient} obtains the isomorphism
by identifying the kernel of the composite and applying Mathlib's first
isomorphism theorem. The finite field calculation combines Mathlib's
index formula for an endomorphism with its bounds and exact counts for
roots of unity.
\module{ManuscriptIBAW.TypeC.ConformalMultiplier} applies them to the
matrix groups with the specified symplectic form.
\end{implementation}

\section{Ordinary characters over a common cyclotomic field}

Let $G$ be finite, let $N$ be a positive multiple of $|G|$, and set
$L=\Q(\zeta_N)$. Fix embeddings $\iota_{\C}\colon L\longrightarrow\C$
and $\iota_K\colon L\longrightarrow K$, where $K$ has characteristic
zero. For character independence and the criterion for realisability,
see, for example, Serre, Section 12.1, Propositions 32--33. The cyclotomic
realisation theorem appears there as Theorem 24 (Brauer), with its
corollary, in Section 12.3 \cite{Serre1977}.

Assume that a realisation statement supplies, for each complex
irreducible character $\chi$, characters $\chi_L\in\Irr_L(G)$ and
$\chi_K\in\Irr_K(G)$. If the chosen basic set associates a simple $K[G]$-module $V_\chi$ with
$\chi$, the character of its representation $\rho_\chi$ must equal $\chi_K$. Thus the required identities are
\[
 \iota_{\C}(\chi_L(g))=\chi(g),\qquad
 \iota_K(\chi_L(g))=\chi_K(g)=\tr_K(\rho_\chi(g))
 \qquad(g\in G).
\]
These identities relate the two realisations to the selected labels.
They can also be obtained by constructing those labels from the
realisation statement.

Suppose that a group of automorphisms permutes a chosen set of complex
characters. Injectivity of $\iota_{\C}$ transports the character
identities to $L$, and applying $\iota_K$ gives the same identities
over $K$. Equality of the traces of the corresponding irreducible
representations gives their isomorphism over $K$, and hence equality
of their classes in the exact Grothendieck group. Thus the action on
the ordinary labels used by the decomposition homomorphism follows
from the action on the complex characters and the stated realisations.

The argument uses only the two embeddings of $L$, so it requires neither
an embedding of $\C$ into $K$ nor algebraic closure of $K$.
Existence of the realisations is an assumption. Their compatibility with
automorphisms follows from these trace identities.

\begin{implementation}
The common field is Mathlib's cyclotomic field. Trace separation uses
Mathlib's character averaging formula for the dimension of the intertwiner
space and its theorem that a nonzero intertwiner between irreducibles is
bijective. ModularRep combines these results over the stated field $K$.
In \module{ManuscriptIBAW.Characters.CyclotomicOrdinaryLabels},
\lean{OrdinaryLabelRealisation.ordinary_value} is the trace identity.
For labels obtained from source realisations,
\lean{SerreRealisationSource.labelRealisation} proves it.
\end{implementation}

\section{Primitive blocks and rational series}

Let $(b_B)_{B\in\mathcal B}$ be a finite complete family of nonzero
orthogonal central idempotents of a ring. Assign to each $B$ a parameter
$t_B$ in a finite group $H^*$. For $t\in H^*$, define
\[
 e_t=\sum_{\substack{B\in\mathcal B\\t_B\sim t}} b_B,
\]
where $\sim$ denotes conjugacy in $H^*$. Orthogonality gives
\[
 b_Be_t=
 \begin{cases}
 b_B,&t_B\sim t,\\
 0,&t_B\not\sim t.
 \end{cases}
\]
Since $b_B\ne0$, we have $b_Be_t=b_B$ if and only if $t_B$ is
conjugate to $t$.

In Jordan reduction the idempotents are the blocks of the
specified group algebra. Each $t_B$ is required to be semisimple and of
order prime to $\ell$. Let $s$ index a geometric series, and write $t_s\in H^*$ for its
assigned element of the finite dual group. The geometric series idempotent at $s$ is identified
with this same sum $e_{t_s}$. Completeness of the chosen semisimple
parameters then places every block in one of the series.

The label used to test strict quasi-isolation is also identified with
the image of $t_B$ under a specified injective homomorphism from $H^*$
to the algebraic dual group. The intended map includes rational points
in the algebraic group. Compatibility of series with automorphisms is a
separate identity for the induced map of group algebras and the dual action.
These identifications specify the blocks, labels and action to which
Jordan reduction applies. Their interpretation as rational series in
algebraic groups remains an assumption.

\begin{implementation}
\module{ManuscriptIBAW.Jordan.SeriesBinding} defines
\lean{rationalSeriesIdempotent} and the identification record
\lean{SeriesBinding} using the block family in ModularRep.
Mathlib supplies the group algebra and finite sum operations. The local
calculation applies orthogonality to this specified sum, while its
identification with the geometric series remains a hypothesis.
The multiplication formula is
\lean{block_mul_rationalSeriesIdempotent}. The membership equivalence and
coverage of the blocks are \lean{SeriesBinding.block_belongs_iff}
and \lean{SeriesBinding.every_block_belongs}, all in
\lean{ManuscriptIBAW.Jordan}.
\end{implementation}


\backmatter
\begin{thebibliography}{99}
\addcontentsline{toc}{chapter}{Bibliography}
\setlength{\itemsep}{2pt}
\setlength{\parsep}{0pt}

\ManualReference{Alperin1987}
\ManualReference{Benson1998}
\ManualReference{Bouc2000}
\ManualReference{Conlon1968}
\ManualReference{Isaacs1976}
\ManualReference{Navarro1998}
\ManualReference{NavarroTiep2011}
\ManualReference{Serre1977}

\end{thebibliography}

\end{document}
